\documentclass[11pt]{article}
\usepackage[T1]{fontenc}
\usepackage[utf8]{inputenc}
\usepackage{lmodern,amsmath,amssymb,amsthm,mathtools,mathrsfs,geometry,microtype,graphicx,booktabs,longtable,array}
\usepackage[hidelinks]{hyperref}
\hypersetup{pdftitle={Natural boundaries of anisotropic Ising susceptibility via real-cycle pole estimates},pdfauthor={Dehao Lin},pdfsubject={Fixed-ratio natural boundaries and isotropic full-sum estimates},pdfkeywords={Ising model, susceptibility, natural boundary, anisotropic}}
\geometry{margin=27mm}
\setlength{\emergencystretch}{2em}
\raggedbottom
\numberwithin{equation}{section}
\newtheorem{theorem}{Theorem}[section]
\newtheorem{lemma}[theorem]{Lemma}
\newtheorem{proposition}[theorem]{Proposition}
\newtheorem{corollary}[theorem]{Corollary}
\theoremstyle{definition}
\newtheorem{definition}[theorem]{Definition}
\newtheorem{remark}[theorem]{Remark}
\newcommand{\C}{\mathbb C}
\newcommand{\R}{\mathbb R}
\newcommand{\Z}{\mathbb Z}
\newcommand{\T}{\mathbb T}
\newcommand{\dd}{\,\mathrm d}
\newcommand{\ii}{\mathrm i}
\newcommand{\eps}{\varepsilon}
\newcommand{\Pf}{\operatorname{Pf}}
\newcommand{\Res}{\operatorname{Res}}
\newcommand{\dist}{\operatorname{dist}}
\newcommand{\cS}{\mathcal S}
\newcommand{\cA}{\mathscr A}
\newcommand{\OH}{\Omega_H}
\newcommand{\OL}{\Omega_L}
\newcommand{\XH}{\mathcal X_H}
\newcommand{\XL}{\mathcal X_L}
\newcommand{\norm}[1]{\left\lVert#1\right\rVert}
\title{Natural boundaries of anisotropic Ising susceptibility\newline
via real-cycle pole estimates}
\author{Dehao Lin\\{\small School of Physics, Sun Yat-sen University}\\{\small ORCID: \href{https://orcid.org/0009-0001-4551-8490}{0009-0001-4551-8490}}\\{\small\texttt{lindh9@mail2.sysu.edu.cn}}}
\date{29 October 2026}
\begin{document}
\maketitle
\begin{abstract}
For every fixed finite positive coupling ratio, we determine natural
boundaries of the complete high-temperature and pure-phase connected
low-temperature bulk susceptibilities. The high-temperature domain is
an explicit bounded lens; the low-temperature domain is an explicit
right epigraph. An endpoint boundary criterion identifies both domains,
and one real-analytic profile difference controls all exceptional
boundary points. On regular positive high-temperature edges and both
low-temperature edge signs, a fixed Vandermonde subproduct gives one
absolute estimate for the entire differentiated particle tail.
At negative high-temperature edges, where odd particles cannot have an
exact endpoint resonance, direct pole expansion on the original real
cycle and an elementary angular concentration estimate give an exact
beta-period model. Selected Taylor orders isolate one particle index
and prove nonextension of the complete odd sum. All lattice sites,
anisotropic factors, physical field normalization and pure-phase
subtraction are retained. At equal coupling a shared constrained-coordinate
proof treats both parities and both high-temperature boundary signs.
Finite-cell contraction, a fixed collision subproduct and an exact
phase-preserving current control the complete differentiated tails.
The resulting equal-coupling obstruction transfers to the physical
temperature domains through the hyperbolic-sine covering.
\end{abstract}
\section*{Introduction}
The complex-temperature susceptibility of the square-lattice Ising model
is governed by the interaction between singularities of its individual
form-factor integrals and the convergence of their infinite sum.
Orrick, Nickel, Guttmann and Perk \cite[Section 3.2]{ONGP} analyzed the
order dependence of the isotropic singularity amplitudes, while
Boukraa, Hassani and Maillard \cite[Section 6.2]{BHM} studied the
anisotropic cosine-grid singularity equations and their accumulation
in two complex couplings. Fixing the positive physical coupling ratio
selects a one-complex-variable temperature slice and its two physical
continuations.

We prove natural boundaries for the full bulk susceptibilities on
the explicit fixed-ratio domains stated below. The proof uses endpoint
profiles to determine the domains, a retained Vandermonde subproduct
to control all higher particle orders at regular edges, and a direct
comparison of pole integrals on the real cycle at negative
high-temperature edges. Section~\ref{sec:iso-internal} proves the
equal-coupling case by a unified contour deformation and
phase-preserving current argument for both particle parities.

Section~\ref{sec:physical} fixes the physical state, field normalization
and complete lattice sum. Section~\ref{sec:geometry} determines the
spectral domains and the dense regular arcs. Sections~\ref{sec:regular}
and \ref{sec:realcycle}--\ref{sec:jets} prove the two local
noncancellation mechanisms, and Section~\ref{sec:closure} includes the
exceptional anisotropic boundary points. Section~\ref{sec:iso-internal}
then supplies the isotropic estimates and the temperature covering.

\subsection*{AI Declaration}
ChatGPT was used extensively in the mathematical research, proof development,
and preparation of this manuscript.

\tableofcontents

\section{Fixed-temperature slice and the physical observable}\label{sec:physical}

Fix $J_x>0$ and $\lambda=J_y/J_x\in(0,\infty)$. Throughout,
\begin{equation}\label{eq:variables}
 z=2\beta J_x,\qquad S=\sinh z,\quad T=\sinh(\lambda z),\quad
 A=\cosh z\cosh(\lambda z),\quad p=ST.
\end{equation}
The ratio $\lambda$ is fixed before any complex continuation; in particular,
the generally varying ratio $T/S$ is not the fixed coupling ratio.
There is a unique $z_c>0$ with $\sinh z_c\sinh(\lambda z_c)=1$.
Define strictly increasing functions on $[0,\infty)$ by
\begin{equation}\label{eq:physicalPQ}
 P_\lambda(x)=\cosh x\tanh(\lambda x),\qquad
 Q_\lambda(x)=\cosh(\lambda x)\tanh x.
\end{equation}
Their values at $z_c$ are both one. Set
\begin{align}
 Y_\lambda&=\frac{\pi}{2\max(1,\lambda)},\nonumber\\
 H_\lambda(y)&=\min\{P_\lambda^{-1}(\cos y),
                      Q_\lambda^{-1}(\cos(\lambda y))\},
                      &&|y|<Y_\lambda,\label{eq:Hdef}\\
 L_\lambda(y)&=\max\{P_\lambda^{-1}(|\cos y|),
                      Q_\lambda^{-1}(|\cos(\lambda y)|)\},
                      &&y\in\R.\label{eq:Ldef}
\end{align}
The domains of interest are specified independently of susceptibility:
\begin{equation}\label{eq:domains}
 \OH=\{x+\ii y:|y|<Y_\lambda,\ |x|<H_\lambda(y)\},\qquad
 \OL=\{x+\ii y:x>L_\lambda(y)\}.
\end{equation}
The value of $H_\lambda$ at each endpoint is its continuous limit zero.

\begin{theorem}[Full bulk susceptibility at every positive ratio]\label{thm:aniso}
Let $0<\lambda<\infty$. The high-temperature physical
bulk susceptibility germ extends holomorphically to $\OH$, and
$\partial\OH$ is its natural boundary. The low-temperature physical
pure-phase connected bulk susceptibility germ extends holomorphically to
$\OL$, and the entire graph $\partial\OL=\{L_\lambda(y)+\ii y:y\in\R\}$
is its natural boundary. Here natural boundary means that no boundary
point has a neighborhood carrying a holomorphic function agreeing with
the indicated continuation on its domain side. Both statements concern
the entire lattice and all particle indices. They apply separately to
$\mathcal X=\beta^{-1}\chi$ and to $\chi$ with the field convention below.
\end{theorem}

The proof for unequal couplings is developed in the following sections.
Section~\ref{sec:iso-internal} proves the isotropic slice directly from
the same full physical product. No limiting estimate across the two
geometries is needed. No simultaneous uniformity in $\lambda$ is asserted:
all local constants may depend on the fixed ratio and selected boundary
point. Rational and irrational ratios are included.

\subsection{Field normalization and state}
Let the magnetic field have energy units, so the finite-volume Hamiltonian
contains $-h\sum_v\sigma_v$, with $\sigma_v\in\{-1,1\}$ and magnetic
moment set to one. Direct differentiation of the finite partition
function gives
\[
 \left.\frac{\partial\langle\sigma_0\rangle}{\partial h}\right|_{h=0}
 =\beta\sum_v\bigl(\langle\sigma_0\sigma_v\rangle
                         -\langle\sigma_0\rangle\langle\sigma_v\rangle\bigr).
\]
Thus the normalized response is the covariance sum $\mathcal X$, and
$\chi=\beta\mathcal X$. A dimensionless field $\eta=\beta h$ instead
has response $\mathcal X$. A specified nonunit magnetic moment multiplies
both by its square.

On $0<z<z_c$ the Gibbs state is unique and its magnetization is zero:
\[
 \XH(z)=\sum_{(m,l)\in\Z^2}\langle\sigma_{00}\sigma_{ml}\rangle.
\]
On $z>z_c$ take the thermodynamic plus state, or equivalently the minus
state for this connected observable:
\[
 \XL(z)=\sum_{(m,l)\in\Z^2}
       \bigl(\langle\sigma_{00}\sigma_{ml}\rangle_+-M^2\bigr),
 \qquad M^2=(1-p^{-2})^{1/4}.
\]
The fourth root is positive on the displayed real interval.
For the energy-field response in this phase, the derivative is the
zero-field pure-phase linear response, equivalently the limit from
$h>0$ in the plus phase (or from $h<0$ in the minus phase).
The covariance formula is obtained with that phase fixed before taking
the zero-field limit. A symmetric
mixture has the same two-spin expectation but zero one-spin mean; using
its zero mean as the subtraction would leave the nondecaying term $M^2$
and would define a different, divergent sum. This distinction is retained
through the thermodynamic form-factor limit below.

The order facts needed here hold for arbitrary real site fields.
Indeed the conditional probability of a plus spin at $v$ is
\[
 \left(1+\exp\left\{-2\beta
       \left(h_v+\sum_{u\sim v}J_{vu}\sigma_u\right)\right\}\right)^{-1},
\]
with the fixed boundary spins included in the sum.  It increases with
every neighboring spin, boundary spin and field.  Updating two ordered
configurations at the same site with the same uniform random number
therefore preserves their order.  In a finite box, for fixed finite
parameters, a full sweep has a uniformly positive probability of
setting every spin to plus.  Coupling two copies with those same
parameters consequently proves convergence to the unique stationary
law; the Gibbs law is stationary under each conditional update.
Taking the limit in the ordered coupling proves finite-volume field
and boundary monotonicity.  Plus and minus volume limits inherit it
and bound every Gibbs state, by conditioning such a state outside each
finite box.  Finite-support field changes of the plus state are exactly
its bounded local tilts, as follows by taking the same limit in the
finite-volume tilt formula.  This proves the monotonicity statements
also at negative fields, where GHS itself is not being applied.
Differentiating finite-volume monotonicity in a single site field
also gives nonnegative two-spin covariances; their volume limits
remain nonnegative.

The covariance sum also equals the literal uniform-field derivative after
the thermodynamic limit. Here is a precise justification once its finite
absolute sum, proved below, is available. The GHS inequality for arbitrary
nonnegative couplings and site fields \cite[Theorem 1.5]{CGN} says that
truncated two-spin correlations decrease as any nonnegative site field
increases. Pinning the spins removed from a larger plus box by fields
tending to $+\infty$ consequently shows that the plus-box covariances
increase with the box. Their sums converge monotonically to $\mathcal X$.
If $m_V(h)=\langle\sigma_0\rangle_{V,+,h}$, GHS also makes $m_V$
concave for $h\ge0$. Thus
\[
 m_V(h)-m_V(0)\le h m'_V(0)
 =\beta h\sum_{v\in V}\operatorname{Cov}_{V,+,0}(\sigma_0,\sigma_v)
 \le\beta h\mathcal X.
\]
Passing to the plus-state limit gives the upper right-derivative bound
$m'(0+)\le\beta\mathcal X$. Conversely, turn on $h$ only in a finite
set $E$ in the infinite plus state. This is a bounded local change of
measure, so its derivative at zero is exactly
$\beta\sum_{v\in E}\operatorname{Cov}_{+,0}(\sigma_0,\sigma_v)$.
Ferromagnetic field monotonicity makes the uniform-field magnetization
at least this local-field magnetization, with equality at $h=0$.
Letting $E$ increase to the lattice gives the reverse derivative bound.
Hence $m'(0+)=\beta\mathcal X$. In HT the zero-field plus and minus
states coincide and have mean zero. Spin flip gives the same derivative
for the minus state from $h<0$. To check the plus state on that side as
well, let $m_E(h)$ be the finite-support tilt of the zero-field state.
Field monotonicity gives
\[
 m_-(h)\le m_+(h)\le m_E(h)\qquad(h<0),
\]
and division by the negative $h$ reverses both inequalities. For each
fixed $E$, the lower bound on the quotient tends to
$\beta\sum_{v\in E}\operatorname{Cov}_0(\sigma_0,\sigma_v)$;
the upper bound tends to $\beta\mathcal X$. Exhausting the lattice
therefore proves the bilateral HT derivative. Other thermodynamic states
lie between the extremal ones. In LT the stated right-sided plus-phase
convention applies. The upper and lower derivative bounds identify
the response after the thermodynamic limit; absolute summability
controls the exhaustion of the lattice.

\subsection{Published real-temperature form-factor interface}
Away from criticality on either positive real interval set
\[
 B=T/S>0,\qquad C=A/S,\qquad
 \cosh\gamma(q)=W(q)=C-B\cos q,\qquad \gamma(q)>0.
\]
Let $f_N(m,l)$ be
\begin{equation}\label{eq:corrFF}
 \frac{B^{N^2/2}}{N!}\int_{\T^N}
 \prod_{i=1}^N\frac{e^{-|m|\gamma(q_i)+\ii lq_i}}{\sinh\gamma(q_i)}
 \prod_{i<j}\left[
 \frac{\sin((q_i-q_j)/2)}{\sinh((\gamma(q_i)+\gamma(q_j))/2)}
 \right]^2\prod_{i=1}^N\frac{\dd q_i}{2\pi}.
\end{equation}
The exact infinite-plane correlation interface used here is
\begin{align}
 \langle\sigma_{00}\sigma_{ml}\rangle
   &=\xi_H\sum_{N\ge1,\ N\text{ odd}}f_N(m,l),
   &\xi_H&=\frac{(1-p^2)^{1/4}}{\sqrt p},\label{eq:HTcorr}\\
 \langle\sigma_{00}\sigma_{ml}\rangle_+
   &=\xi_L\left(1+\sum_{N\ge2,\ N\text{ even}}f_N(m,l)\right),
   &\xi_L&=(1-p^{-2})^{1/4}=M^2.\label{eq:LTcorr}
\end{align}
All powers are fixed by their positive real values. These are the
anisotropic vacuum-to-state spin form factors of
\cite{BL,IL}, in the bulk normalization of \cite{ONGP}. In particular,
$B^{N^2/2}$ is part of the $N$th term. Replacing it by a fixed prefactor
would change both the observable and the infinite-index problem.

The finite-cylinder formulas in \cite[(35)--(39)]{BL} concern arbitrary separations,
not just a single lattice row. Vacuum-to-state matrix elements are also
covered by \cite[Theorem 6.1]{IL} in the LT phase and by its
Remark 6.2, equations (6.8)--(6.10), with the HT vacuum change.
On a fixed real noncritical compact set, the one-body dispersion is
analytic in a common angular strip and has a positive real gap. The
finite-circumference correction functions $\eta$, $\Lambda^{-1}$ and
$\log\xi_T$ tend uniformly to zero: their sum-minus-integral expressions
are controlled by the exponentially convergent trapezoidal rule for
periodic analytic functions. They are uniformly bounded before the limit.

The particle limit is uniform as well. For real roots
$0<\rho(q)=e^{-\gamma(q)}\le r<1$, the Schur factor
$|(\rho_i-\rho_j)/(1-\rho_i\rho_j)|$ is at most a constant strictly below
one. The matched-Pfaffian argument proved in
Section~\ref{sec:regular} applies to the discrete angular grids as an
algebraic identity, with their normalized counting measures. Its
single-body finite-size corrections cost only $C^N$. Consequently the
finite-circumference coefficients, at each fixed separation, have a
common bound $\exp(-cN^2+O(N\log(N+1)))$, including their factorial
normalization. This is summable independently of circumference, so the
particle sum passes through the cylinder limit; taking each fixed-$N$
Riemann-sum limit gives \eqref{eq:corrFF}.

Every local cylinder limit is a translation-invariant, spin-flip invariant
Gibbs state for these finite-range positive couplings. Apply
\cite[Theorem 1]{Raoufi} with the vertex-transitive subgroup consisting
of translations: unequal horizontal and vertical couplings satisfy all
four stated hypotheses. Such a state is the equal mixture of the plus
and minus states. Its even spin correlations are those of the plus
state. This identifies the two-point function in \eqref{eq:LTcorr};
the connected pure-phase observable is obtained only after removing
the zero-particle term $\xi_L$.

The spontaneous magnetization can be identified from this same
interface.  The maximal plus Gibbs state is extremal: in any convex
decomposition into Gibbs states, each component is dominated by it;
equality of the average forces equality on every increasing cylinder
event, hence equality of the measures.  Its tail is consequently
trivial.  To see the needed implication directly, conditioning on a
tail event of probability strictly between zero and one gives two
Gibbs states and a nontrivial such decomposition.  The reverse
martingale theorem now gives
\[
 \left\|\mathbb E_+(\sigma_0\mid
          \sigma\text{ outside }V)-\langle\sigma_0\rangle_+\right\|_1
       \longrightarrow0.
\]
For $v\notin V$ this norm bounds
$|\operatorname{Cov}_+(\sigma_0,\sigma_v)|$, so the two-point
function tends to $\langle\sigma_0\rangle_+^2$ at large separation.
On a real noncritical interval, the positive dispersion gap and the
uniform particle majorant just proved make every positive-particle
term tend to zero as $|m|\to\infty$ with $l=0$.
Equations \eqref{eq:HTcorr}--\eqref{eq:LTcorr} therefore give
$\langle\sigma_0\rangle_+=0$ in HT and
$\langle\sigma_0\rangle_+^2=\xi_L$ in LT.  The plus mean is
nonnegative by spin flip and extremal order.  In HT the plus and
minus means coincide at every site; their monotone coupling must
then agree at every site almost surely.  Thus the states coincide,
and their extremal bounds prove uniqueness.  This also justifies the
phase and subtraction conventions used above directly from the
stated physical inputs.

\subsection{Symmetric double contours and every lattice site}
For separate contour variables $x,y$ put
\begin{equation}\label{eq:Dpoly}
 D(x,y;z)=A-\frac S2(x+x^{-1})-\frac T2(y+y^{-1}),\qquad
 \cA(v,w)=\frac{v-w}{1-vw}.
\end{equation}
With $X=\prod_i x_i$ and $Y=\prod_i y_i$, the same correlation
coefficient is
\begin{equation}\label{eq:doublecorr}
 \frac{p^{N/2}}{N!}\oint\cdots\oint
 X^{|m|-1}Y^{|l|-1}
 \prod_{i<j}\cA(x_i,x_j)\cA(y_i,y_j)
 \prod_i\frac{\dd x_i\,\dd y_i}{(2\pi\ii)^2D(x_i,y_i;z)}.
\end{equation}
The circles are positively oriented and sufficiently close to, but inside,
the unit circle. To check the normalization, fix $y=e^{\ii q}$ and
take the $x$ residue at the interior root $\rho=e^{-\gamma}$.
Since
\[
 D_x(\rho,y)=\frac{S(1-\rho^2)}{2\rho^2},\qquad
 \frac{\sqrt p}{D_x(\rho,y)}=\sqrt B\frac{\rho}{\sinh\gamma},
\]
and subtraction of two dispersion equations gives
\begin{equation}\label{eq:cancelpair}
 \cA(\rho_i,\rho_j)\cA(y_i,y_j)
 =-B\frac{(y_i-y_j)^2\rho_i\rho_j}
           {y_iy_j(1-\rho_i\rho_j)^2}
 =B\left[\frac{\sin((q_i-q_j)/2)}
               {\sinh((\gamma_i+\gamma_j)/2)}\right]^2=:P_{ij},
\end{equation}
the residue factors supply $B^{N/2}$ and the pair factors supply
$B^{N(N-1)/2}$. This recovers precisely $B^{N^2/2}$.
The remaining factors in \eqref{eq:doublecorr} give the stated Fourier
and spatial decay factors. Reflection in each coordinate is legitimate;
equality under interchange of the two separations at fixed unequal
couplings is not assumed.

The exact full-site multiplier in the double contour is
\begin{equation}\label{eq:allsites}
 \frac{(1+X)(1+Y)}{XY(1-X)(1-Y)}
 =\frac1{XY}+\frac{2(X^{-1}+Y^{-1})}{(1-X)(1-Y)}.
\end{equation}
It is the product of the two sums $\sum_{m\in\Z}X^{|m|-1}$ and
$\sum_{l\in\Z}Y^{|l|-1}$. It counts the origin once, each reflected
axis correctly, and all four open quadrants. If $C_N$ is the full-site
coefficient and $T_N$ the contour with the second numerator but without
its factor two, then $C_N=f_N(0,0)+2T_N$. The onsite sums are one in
the HT phase and $1-M^2$ after the LT subtraction. No axis or onsite
contribution is discarded.

Summing the Fourier variable and the geometric spatial series also gives
the complete angular representation
\begin{equation}\label{eq:fullangular}
 C_N=\frac1{N!}\int_{\sum q_i=0\bmod2\pi}
 G_N(q)\frac{1+Z}{1-Z}\prod_{i=1}^{N-1}\frac{\dd q_i}{2\pi},\quad
 G_N=B^{N/2}\prod_i\frac1{\sinh\gamma_i}\prod_{i<j}P_{ij},\quad
 Z=\prod_i e^{-\gamma_i}.
\end{equation}
There are exactly $N-1$ independent normalized angular measures, with
$q_N=-\sum_{i<N}q_i$ modulo $2\pi$.

\subsection{A holomorphic full sum on an independent spectral domain}
Define
\begin{equation}\label{eq:physicalspectralset}
 \cS_\lambda=\{z:\exists u,v\in[-1,1],\ A(z)=S(z)u+T(z)v\}.
\end{equation}
It is exactly the set where \eqref{eq:Dpoly} has a zero on the product
of two unit circles. Section~\ref{sec:geometry} proves that the components
of $\C\setminus\cS_\lambda$ containing $(0,z_c)$ and $(z_c,\infty)$
are exactly \eqref{eq:domains}, and that both are simply connected.

We record the scalar algebra used below.  For even $N$ the Schur
identity is
\[
 \operatorname{Pf}\left(\frac{x_i-x_j}{1-x_ix_j}\right)_{i,j=1}^N
       =\prod_{i<j}\frac{x_i-x_j}{1-x_ix_j}.
\]
For a short proof put $u_i=(1+x_i)/(1-x_i)$.  Each entry becomes
$(u_i-u_j)/(u_i+u_j)$.  Induction in $N$ compares the two rational
functions at every possible simple pole $u_i+u_j=0$: the residues
agree by the induction hypothesis, because the two cross factors
involving any third index multiply to one.  Their difference has no
poles, is homogeneous of degree zero, and is alternating, so it is
zero.  The case $N=2$ starts the induction; removable special values
follow by rational continuation.  For odd $N$, append $x_{N+1}=-1$;
all its entries are the removable constant one.  This proves the
matching identities at both parities.

Here is the needed convergence argument independently of any proposed
natural-boundary conclusion. On a compact subset of the spectral
complement the torus polynomial has a nonzero minimum modulus. Move
both radii a small common amount inside one, within a product annulus
where this gap persists. The resulting contours can be held fixed on
a parameter neighborhood. Their Schur entries are bounded. The Schur
Pfaffian identity, followed by Hadamard's determinant inequality, bounds
the product of the two Pfaffians by
$C^N N^{(N+1)/2}$, with the harmless extra odd dimension obtained by
adjoining a constant ghost entry. After the $1/N!$ normalization and
the bounded one-body and full-site factors this is summable in $N$.
It is locally normally summable together with every fixed parameter
derivative by Cauchy's inequality on a slightly larger neighborhood.

On a physical real compact set the same contours in
\eqref{eq:doublecorr} additionally give
$r_x^{N|m|}r_y^{N|l|}$ with $r_x,r_y<1$. This proves absolute summability
in the lattice sites and particles and justifies their interchange.
Thus \eqref{eq:fullangular} and the double contour define the same
physical germs, with a proven full-site sum.

Write $F_N=p^{-N/2}C_N$ for the unscaled double contour, whose integrand
uses \eqref{eq:Dpoly} and \eqref{eq:allsites}. The complete sums are
\begin{align}
 \XH&=(1-p^2)^{1/4}\sum_{N\ge1,\ N\text{ odd}}p^{(N-1)/2}F_N,
                                                        \label{eq:globalHT}\\
 \XL&=(1-p^{-2})^{1/4}\sum_{N\ge2,\ N\text{ even}}p^{N/2}F_N.
                                                        \label{eq:globalLT}
\end{align}
The powers inside both sums are integers. The identity
$A^2=(1+S^2)(1+T^2)$ implies that every $p^2=1$ point lies in
$\cS_\lambda$. Also, $p$ never vanishes in $\OL\subset\{\Re z>0\}$;
in $\OH$ its only zero is the double zero at the origin. Hence the
prefactors in these formulas have unique holomorphic branches on the
specified simply connected domains, fixed by the real physical values.
Formula \eqref{eq:globalHT} removes the apparent HT $1/\sqrt p$
singularity at zero; $\XH(0)=1$. In particular, the one-particle term is
\[
 \XH^{(1)}=\frac{(1-p^2)^{1/4}}{A-S-T},
\]
which gives $1+2\beta(J_x+J_y)+O(\beta^2)$ at the origin.

For the angular formulas the spectral root is globally specified by
\[
 \rho=W-\sqrt{W^2-1},\qquad |\rho|<1,
 \qquad W=(A-T\cos q)/S.
\]
The square root has asymptotic $W$ at infinity. For $z$ in the spectral
complement, $W$ avoids $[-1,1]$ for every real $q$, so the inside root
is unique. At $S=0$ its root equation
$S\rho^2-2(A-T\cos q)\rho+S=0$ gives its removable value zero.
A local logarithm $\gamma=-\log\rho$ is used only where nonzero and
needed. No unsupported global logarithm or change of transfer vacuum
is imposed by periodic complex coupling shifts.



\section{The two physical spectral components and their regular edges}
\label{sec:geometry}

Fix $\lambda>0$.  We continue to write
$z=x+\ii y$, $S=\sinh z$, $T=\sinh(\lambda z)$,
$A=\cosh z\cosh(\lambda z)$, and $p=ST$.
Define independently of susceptibility
\begin{equation}\label{eq:spectralset}
 \cS_\lambda=\{z:A(z)=S(z)u+T(z)v
                       \text{ for some }u,v\in[-1,1]\}.
\end{equation}
It is closed, by compactness of the square of representing pairs.
It is also exactly the zero set of the double-contour polynomial
on the product of unit circles.  We will identify its two physical
complement components directly from the four endpoint equations.

\subsection{Endpoint curves determine the physical components}

Put
\begin{align}
 F_1(x,y)&=\cosh^2(\lambda x)\cos^2y
             -\cosh^2x\sinh^2(\lambda x),\label{eq:F1}\\
 F_2(x,y)&=\cosh^2x\cos^2(\lambda y)
             -\sinh^2x\cosh^2(\lambda x),\label{eq:F2}\\
 Z_\lambda&=\{x+\ii y:F_1(x,y)F_2(x,y)=0\}.
 \label{eq:endpointlocus}
\end{align}

\begin{lemma}[Endpoint boundary criterion]\label{lem:endpoint-sandwich}
For every positive fixed ratio,
\begin{equation}\label{eq:endpoint-sandwich}
 Z_\lambda\subset\cS_\lambda,
 \qquad \partial\cS_\lambda\subset Z_\lambda.
\end{equation}
\end{lemma}
\begin{proof}
For $x>0$, neither $S$ nor $\tanh(z/2)$ vanishes or has a pole.
For $\epsilon\in\{1,-1\}$ define
\begin{equation}\label{eq:endpointroot}
 r_\epsilon(z)=e^{\epsilon\lambda z}\tanh(z/2),
 \qquad
 \frac{r_\epsilon+r_\epsilon^{-1}}2
       =\frac{A-\epsilon T}{S}.
\end{equation}
The second identity follows by expanding the two exponentials.
Moreover,
\[
 |\tanh(z/2)|^2=\frac{\cosh x-\cos y}{\cosh x+\cos y},
 \qquad
 |r_\epsilon|=1
 \ \Longleftrightarrow\ 
 \cos y=\epsilon\cosh x\tanh(\lambda x).
\]
Thus an equation $A=Su+\epsilon T$ with $u\in[-1,1]$ is
equivalent to $F_1=0$, with the appropriate sign $\epsilon$:
the roots of $r+r^{-1}=2u$ are on the unit circle, and conversely
a unit root gives a real $u\in[-1,1]$.
Axis exchange gives the corresponding equivalence for $F_2=0$.
At $x=0$, $S,T$ are imaginary and $A$ is real, so membership in
$\cS_\lambda$ is exactly $\cos y\cos(\lambda y)=0$; this is
also $F_1F_2=0$.  Reflection $z\mapsto-\overline z$ handles
$x<0$.  In particular $Z_\lambda\subset\cS_\lambda$ everywhere.

Now take $z_0\in\partial\cS_\lambda$.  If the real-linear map
$(u,v)\mapsto S(z_0)u+T(z_0)v$ has rank two and its representing
pair is in the open square, the inverse real-linear map supplies
representing pairs in that open square for all nearby $z$.
This would make $z_0$ an interior point of $\cS_\lambda$.
Hence its representing pair has an endpoint coordinate.
If the map has rank at most one, a nonempty affine fiber through
the square meets the boundary of the square: move along a nonzero
kernel direction until first reaching its boundary.  Thus there
is an endpoint representation in this case as well.  (Rank zero
does not actually occur at a candidate, since simultaneous zeros
of $S,T$ have $A=\pm1$.)  The endpoint equivalences just proved
place $z_0$ in $Z_\lambda$.
\end{proof}

Let $z_c>0$ satisfy $\sinh z_c\sinh(\lambda z_c)=1$.
The functions
\[
 P_\lambda(x)=\cosh x\tanh(\lambda x),\qquad
 Q_\lambda(x)=\cosh(\lambda x)\tanh x
\]
strictly increase from zero to infinity on $[0,\infty)$ and
both equal one at $z_c$.  In particular,
\begin{align*}
 F_1&=\cosh^2(\lambda x)
              \bigl(\cos^2y-P_\lambda(x)^2\bigr),\\
 F_2&=\cosh^2x
              \bigl(\cos^2(\lambda y)-Q_\lambda(x)^2\bigr).
\end{align*}

\begin{theorem}[Exact physical components]\label{thm:components}
The components of $\C\setminus\cS_\lambda$ containing
$(0,z_c)$ and $(z_c,\infty)$ are precisely the domains
$\OH$ and $\OL$ in \eqref{eq:domains}.  They are simply connected.
Their boundaries are the complete lens boundary and the complete
graph $x=L_\lambda(y)$, respectively, and $0\le L_\lambda\le z_c$.
\end{theorem}
\begin{proof}
In the displayed lens, both $F_1,F_2$ are positive.
In the displayed right epigraph they are both negative.
Thus both domains are disjoint from $Z_\lambda$, whereas their
entire boundaries lie in $Z_\lambda$.  At the two lens tips this
last statement follows from the first vanishing cosine and $x=0$.
The lens is star shaped about zero: its positive profile decreases
with $|y|$.  The epigraph is connected, since $L_\lambda\le z_c$.

By Lemma~\ref{lem:endpoint-sandwich}, $\partial\cS_\lambda$ does
not meet either connected domain.  Therefore the intersection of
$\cS_\lambda$ with each domain is both relatively open and relatively closed.
At zero it is absent, and it is also absent at every real $z>z_c$,
because
\[
 (A-S-T)(A+S+T)=(ST-1)^2>0
\]
there, with all factors except possibly $A-S-T$ positive.
Thus both domains lie in the spectral complement.  Their entire
boundaries lie in $\cS_\lambda$, so no path in that complement can
leave either domain.  They are the claimed components.

Star shapedness gives simple connectivity of the lens.  The
homeomorphism $(x,y)\mapsto(x-L_\lambda(y),y)$ maps the epigraph
to the right half-plane and gives its simple connectivity.
\end{proof}

\subsection{A single analytic profile difference controls every exception}

For this purpose extend $P_\lambda,Q_\lambda$ to all real $x$.
They are odd real-analytic bijections of $\R$, with everywhere
positive derivatives.  For example
\[
 P_\lambda'(x)=\sinh x\tanh(\lambda x)
                  +\lambda\cosh x\operatorname{sech}^2(\lambda x)>0.
\]
Their inverse functions are therefore real analytic on all of $\R$.
Set
\begin{equation}\label{eq:signed-profiles}
 R(y)=P_\lambda^{-1}(\cos y),\qquad
 U(y)=Q_\lambda^{-1}(\cos(\lambda y)),\qquad
 E(y)=R(y)^2-U(y)^2.
\end{equation}
In these terms $L_\lambda=\max(|R|,|U|)$, and
$H_\lambda=\min(R,U)$ on the central strip.
The intersections of the two endpoint graphs in $x\ge0$ are
exactly the zeros of $E$.

\begin{lemma}[Isolated profile intersections]\label{lem:profile-discrete}
If $\lambda\ne1$, the zeros of $E$ are locally discrete on $\R$.
\end{lemma}
\begin{proof}
Near zero the inverse-function Taylor expansions give
\[
 E(y)=z_c\left(\frac{\lambda^2}{Q_\lambda'(z_c)}
                      -\frac1{P_\lambda'(z_c)}\right)y^2+O(y^4)
     =z_c\bigl(\lambda\coth(\lambda z_c)-\coth z_c\bigr)y^2
       +O(y^4).
\]
For the second equality put $a=\sinh z_c$ and use
$\sinh(\lambda z_c)=a^{-1}$:
$P_\lambda'(z_c)=a(1+\lambda a)/\cosh z_c$ and
$Q_\lambda'(z_c)=(1+\lambda a)/(a\cosh z_c)$.
The coefficient is nonzero for $\lambda\ne1$, because $t\coth t$
strictly increases on $(0,\infty)$:
its derivative has numerator $\sinh t\cosh t-t>0$.
Thus the globally real-analytic function $E$ is not identically
zero, so its zeros are locally discrete.
\end{proof}

The analytic-profile argument applies also at simultaneous cosine zeros.
Such zeros exist exactly when $\lambda=r/s$ in reduced form
with $r,s$ odd, at $y=(2m+1)\pi s/2$.  They are precisely the
points where the LT boundary can reach $x=0$.  At each such $y_0$,
the two nonnegative endpoint profiles have leading terms
$|y-y_0|/\lambda$ and $\lambda|y-y_0|$; the signed profiles in
\eqref{eq:signed-profiles} remain analytic there.

We next check that every remaining physical boundary point has all
the regularity properties used by the particle estimates.
Write $D=\Im(\overline S T)$.
Suppose $x>0$, $F_1=0$ and $F_2\ne0$.  By the endpoint criterion,
\begin{equation}\label{eq:endpointrepresentation}
 A=Su+\epsilon T,\qquad
 u\in[-1,1],\qquad \epsilon=\operatorname{sign}(\cos y).
\end{equation}
The following identity is decisive:
\begin{equation}\label{eq:endpoint-defect}
 D^2(1-u^2)=\tanh^2(\lambda x)F_2(x,y)^2.
\end{equation}
Here is a short derivation including the endpoint cases.
Put $K=\cosh z\,T-\epsilon\cosh(\lambda z)$.
Direct expansion gives $K^2=S^2(u^2-1)$, so
\[
 \bigl[\Re(\overline K T)\bigr]^2=D^2(1-u^2).
\]
Writing $c=\cosh x$, $h=\sinh x$,
$C=\cosh(\lambda x)$, $H=\sinh(\lambda x)$,
the relation $F_1=0$ is $\cos y=\epsilon cH/C$, and hence
\begin{align*}
 \Re(\overline K T)
 &=c\cos y\,|T|^2-\epsilon HC\\
 &=\epsilon\frac HC\left(c^2(C^2-\cos^2(\lambda y))-C^2\right)
 =-\epsilon\frac HC F_2.
\end{align*}
This proves \eqref{eq:endpoint-defect}.  Its right side is positive,
so $D\ne0$ and $u\in(-1,1)$.

The corresponding edge coordinate is also automatically nonramified.
Indeed
\[
 \frac{r_\epsilon'}{r_\epsilon}=\epsilon\lambda+\frac1S,
 \qquad
 \Re\frac{r_\epsilon'}{r_\epsilon}
     =\epsilon\left(\lambda+
        \frac{\sinh x\sinh(\lambda x)\cosh(\lambda x)}{\cosh x}\right)
     \ne0
\]
on this endpoint curve.  Thus $r_\epsilon'\ne0$.
Since $u=(r_\epsilon+r_\epsilon^{-1})/2\in(-1,1)$,
also $r_\epsilon\ne\pm1$, and
\[
 g_\epsilon(z)=\frac{A-\epsilon T}{S}
 \quad\text{has}\quad g_\epsilon'(z)\ne0.
\]
Axis exchange gives the same statements when $F_2=0$, $F_1\ne0$;
reflection gives them for $x<0$.

\begin{proposition}[Dense regular physical boundary arcs]
\label{prop:regular_dense}
For every $\lambda\ne1$, the regular edge points used in the
particle arguments are dense in both physical boundaries.
Their excluded set is locally discrete.  At each such regular
point the two spectral directions are real-linearly independent,
exactly one representing cosine is $\pm1$, the other is strictly
inside $(-1,1)$, and its holomorphic edge coordinate is nonramified.
\end{proposition}
\begin{proof}
Away from the isolated profile intersections and the two imaginary
lens tips, the active boundary graph has $x\ne0$ and exactly one
$F_i$ equal to zero.  The preceding calculation gives all the stated
properties.  The active profile is analytic there; the boundary
has no isolated point because it is a continuous graph or a lens
curve.  Removing a locally discrete subset therefore leaves a dense
set.  Finally, at a rank-two endpoint representation, the Cramer
coordinates vary continuously and the strict interior coordinate
stays interior.  More explicitly, on an $F_1$ chart the real
representing coordinate $v(z)$ satisfies
$v(z)-\epsilon=|S(z)|^2\Im g_\epsilon(z)/D(z)$.
Thus the equation that this coordinate equals its endpoint is the
nondegenerate real analytic arc just described;
the spectral set occupies one closed side and the complement
occupies the other.  This verifies the local domain-side condition
used below.
\end{proof}

On the right side of the HT lens, both cosines are positive, so
the endpoint sign is positive.  Reflection $z\mapsto-\overline z$
sends $(u,v)$ to $(-u,-v)$ and gives negative endpoints on its left
side.  On the LT boundary the sign is $\operatorname{sign}(\cos y)$
for an $F_1$ edge, and $\operatorname{sign}(\cos(\lambda y))$ for
an $F_2$ edge.  Thus all the signs treated below occur on the
specified physical components.

Each nonramified edge coordinate maps a sufficiently small regular
arc onto a real interval in $(-1,1)$.  Values
$\cos(2\pi k/q)$ with $q$ an odd prime and
$1\le k\le(q-1)/2$ are dense in that interval, since the mesh
tends to zero along the primes.  At a positive endpoint the first
odd resonance is $q$ and the first even resonance is $2q$.
At a negative endpoint exact all-endpoint momentum conservation
requires an even order, whose first resonance is again $2q$.
There is no exact odd endpoint collision there; that is the
separate issue resolved by the negative-HT boundary-jet argument.

\subsection{Canonical roots, prefactors, and the isotropic slice}

\begin{lemma}[Canonical spectral root]\label{lem:canonicalroot}
For $z\notin\cS_\lambda$ and real $q$, the equation
\[
 S\rho^2-2(A-T\cos q)\rho+S=0
\]
has a unique root with $|\rho|<1$, interpreted as $\rho=0$ if
$S=0$.  The root is holomorphic in $z$, locally jointly
holomorphic with the angle, and has a uniform modulus gap below
one on each compact parameter set and the real angle circle.
On both positive noncritical real intervals it is the physical
positive decaying root.
\end{lemma}
\begin{proof}
For $S\ne0$, put $W=(A-T\cos q)/S$.  The spectral-set definition
gives $W\notin[-1,1]$.  The reciprocal roots of
$\rho+\rho^{-1}=2W$ cannot be on the unit circle, so exactly one
is inside.  It is simple, and the implicit-function theorem makes
it locally holomorphic.  Uniqueness glues these local roots.
At $S=0$, the coefficient $A-T\cos q$ is nonzero, so the same
theorem extends the root through zero.  Compactness gives the
uniform gap and a thin complex angular neighborhood.  Finally
$A-S-T>0$ on the positive noncritical real line by the identity
used above, so $W(q)>1$ and the interior root is positive.
\end{proof}

Only local logarithms $\gamma=-\log\rho$ are used where the root
is nonzero.  No global logarithm through a root zero is required.
The identity
\[
 A^2-(\epsilon_1S+\epsilon_2T)^2=(p-\epsilon_1\epsilon_2)^2
\]
shows $\{p^2=1\}\subset\cS_\lambda$.
The only zero of either $S$ or $T$ in the HT lens is the common
zero at $z=0$; neither vanishes in the LT epigraph, where $x>0$.
Thus the physical phase prefactors have the holomorphic branches
asserted in Section~\ref{sec:physical}; the integer-power HT
formula handles the origin.

At $\lambda=1$, write $s=\sinh z$.  For $s\ne0$,
\[
 A/S=s+s^{-1}=u+v\in[-2,2]
 \quad\Longleftrightarrow\quad |s|=1,
\]
so
\begin{equation}\label{eq:isotropicgeometry}
 \cS_1=\{z:|\sinh z|=1\},\qquad
 H_1(y)=\operatorname{arsinh}(\cos y),\quad |y|<\pi/2,
 \qquad L_1(y)=\operatorname{arsinh}(|\cos y|).
\end{equation}
The precise biholomorphism on the HT lens, universal covering on the
LT epigraph, and boundary ramification transfer are proved with the
internal isotropic analytic argument in Section~\ref{sec:iso-internal}.

\section{Regular spectral edges: whole amplitudes and the entire tail}
\label{sec:regular}

At a regular boundary point exactly one of the two real spectral
coordinates is $+1$ or $-1$, the other lies strictly between them, and
the corresponding edge map is a nonramified holomorphic local
coordinate. Interchange the transfer axes when necessary. First take
the endpoint $q_0=0$ and write
\begin{equation}\label{eq:positiveedge}
 W(q;u)=b+u+B(u)(1-\cos q),\qquad
 b=\cos\phi,\quad 0<\phi<\pi,\quad \Im B(0)>0.
\end{equation}
Here $u=W(0)-b$ vanishes at the selected point. The opposite sign of
$\Im B$ is obtained by complex conjugation. On the edge, $W(q)$ is
in the upper half-plane except at $q=0$. Its canonical root
$\rho(q)=e^{-\gamma(q)}$ lies in the lower half-disk, with
$\rho(0)=e^{-\ii\phi}$. Compactness gives
\begin{equation}\label{eq:halfdisc}
 \Im\rho(q)\le-\delta<0,\qquad |\rho(q)|\le1
 \quad(q\in\R).
\end{equation}
There is no real-angle point with $W=\pm1$. The root, a local
$\gamma$, and all one-body and canceled-pair denominators are therefore
holomorphic on a common fixed complex parameter and angular
neighborhood. On the physical-domain approach $u=\ii\eps$, $\eps>0$,
\begin{equation}\label{eq:attenuation}
 \Re\gamma(q;\ii\eps)\ge c\bigl(\eps+\dist(q,0)^2\bigr).
\end{equation}
Indeed $\Im W=\eps+\Im B(\ii\eps)(1-\cos q)$ is bounded below by
that expression, while
$\Im W=\sinh(\Re\gamma)\sin(\Im\gamma)$ on the compact chosen branch.
The same local side belongs to the physical component bordering the
regular edge, by the geometry of Section~\ref{sec:geometry}.

\subsection{A strict Schur bound, with the anisotropy retained}
The elementary identity
\begin{equation}\label{eq:schurdifference}
 |1-vw|^2-|v-w|^2
 =(1-|v|^2)(1-|w|^2)+4\Im v\Im w
\end{equation}
and \eqref{eq:halfdisc} give a constant $\rho_0<1$ such that
$|\cA(\rho_i,\rho_j)|\le\rho_0$ for every pair of real angles.
By continuity a slightly larger constant still below one works on a
fixed complex parameter neighborhood. This assertion continues to hold
there even where a separate inequality $|\rho_i|\le1$ no longer does.
Every complete canceled pair $P_{ij}$ is analytic and bounded on the
same neighborhood.

It is not true that $|P_{ij}|\le1$ there. For example, on
\eqref{eq:positiveedge},
\[
 P(q,-q)=\frac{B\sin^2q}
               {[b+B(1-\cos q)]^2-1}.
\]
When $\Im B$ is small and the real part of the bracket is close to one,
this can be arbitrarily large. The proof below retains the strict bound
only for the uncanceled $\rho$-Schur factors.

\begin{lemma}[Matched-Pfaffian numerator estimate]\label{lem:matching}
For every fixed integer $r\ge0$, on a common parameter disk,
\[
 |\partial_u^rG_N(q;u)|
 \le C_r\exp\{-\kappa N^2+C N\log(N+1)\}
\]
uniformly in the real angles and all positive $N$. The constants may
depend on the edge but not on $N$.
\end{lemma}
\begin{proof}
For even $N$, expand only the Schur Pfaffian in the $y_i=e^{\ii q_i}$
variables. Each of its $(N-1)!!$ matchings contains $N/2$ edges. Combine
each matched $y$ factor with its corresponding $\rho$ factor using
\eqref{eq:cancelpair}; bound that complete pair by a constant $C$.
The other $N(N-2)/2$ factors are bounded by $\rho_0$. Thus
\[
 \left|\prod_{i<j}P_{ij}\right|
 \le (N-1)!!\,C^{N/2}\rho_0^{N(N-2)/2}.
\]
For odd $N$, adjoin a ghost variable $-1$ to the $y$ Pfaffian.
Its Schur entry with every real variable is the removable constant one.
There are $N!!$ matchings, with $(N-1)/2$ real matched edges and one
ghost edge, leaving $(N-1)^2/2$ uncanceled $\rho$ factors. This gives
\[
 \left|\prod_{i<j}P_{ij}\right|
 \le N!!\,C^{(N-1)/2}\rho_0^{(N-1)^2/2}.
\]
The identities initially hold away from reciprocal $y$ pairs and extend
by the analytic canceled formula. No division by a vanishing $\rho$
factor occurs. One-body factors, including $B^{N/2}$, cost $C^N$.
These are Gaussian bounds in $N$ up to $\exp(O(N\log(N+1)))$.
Cauchy's inequality on the common disk proves the numerator-jet
statement. It is applied only to $G_N$, not across a pole of
$(1+Z)/(1-Z)$.
\end{proof}

\subsection{All indices in the differentiated tail}
Put $Q_0(q)=\sum_i\dist(q_i,0)^2$. Equation
\eqref{eq:attenuation} implies
\[
 |1-Z|\ge1-|Z|\ge c\min\{1,N\eps+Q_0(q)\}.
\]
All fixed parameter derivatives of $\gamma$ are bounded. Differentiating
its exponential and the rational factor gives, for every fixed $j$ and
$0\le r\le j$,
\begin{equation}\label{eq:globaljet}
 \left|\partial_u^r\frac{1+Z}{1-Z}\right|
 \le C_jN^j\bigl(1+Q_0(q)^{-j-1}\bigr).
\end{equation}
We retain a fixed number of numerator collision factors before using
Lemma~\ref{lem:matching}. This gives one integrable majorant for every
higher particle index.

Fix an integer $m\ge1$, and write
\[
 \mathcal V_m(q)=\prod_{1\le i<\ell\le m}
                   \sin^2((q_i-q_\ell)/2).
\]
For $N\ge m$, the quotient $G_N/\mathcal V_m$ has an analytic
extension through the zeros of $\mathcal V_m$. Indeed, each complete
pair internal to the first $m$ labels is replaced by
$B/\sinh^2((\gamma_i+\gamma_\ell)/2)$. These pairs and the first
$m$ one-body factors are bounded by a constant depending only on $m$.
The $m(N-m)$ cross pairs cost $C^{m(N-m)}$, and the complementary
numerator on the last $N-m$ labels satisfies
Lemma~\ref{lem:matching}; the empty numerator is interpreted as one.
Thus, on a common complex parameter disk,
\[
 \left|G_N/\mathcal V_m\right|
 \le C_m\exp\{-\kappa(N-m)^2
          +C(N-m)\log(N-m+1)+Cm(N-m)\}.
\]
Cauchy's estimate on a fixed smaller disk, applied to this analytic
quotient, proves for every fixed $r\ge0$ that
\begin{equation}\label{eq:fixed-subvandermonde}
 |\partial_u^rG_N(q;u)|
 \le C_{m,r}\mathcal V_m(q)
       \exp\{-\kappa N^2+C_mN\log(N+1)\}.
\end{equation}
All constants are independent of $N$ and of the real angles.

Let $n$ be the first resonant index of the selected parity, and set
\[
 k=\frac{n^2-1}{2}\quad(n\text{ odd}),\qquad
 k=\frac{n^2}{2}-1\quad(n\text{ even}),\qquad m=n+1.
\]
Every larger same-parity $N$ satisfies $N\ge n+2=m+1$, so the first
$m$ angles are independent integration variables. Choose their
representatives in $[-\pi,\pi]$ and put $R^2=\sum_{i=1}^m q_i^2$.
Then $Q_0\ge R^2$ and
$\mathcal V_m\le C_mR^{m(m-1)}$. The fixed $m$-dimensional
majorant
\[
 \mathcal V_m(q)\bigl(1+R^{-2j-2}\bigr)
\]
is integrable for every $0\le j\le k$: near zero its radial power is
$R^{m^2-2j-3}\dd R$, and $m^2>2j+2$. Every unused normalized
angular measure has mass one. Applying the product rule to
\eqref{eq:fullangular}, and using \eqref{eq:globaljet} and
\eqref{eq:fixed-subvandermonde}, therefore gives
\begin{equation}\label{eq:hightail}
 |\partial_u^jC_N(\ii\eps)|
 \le\frac{C_{n,j}N^j}{N!}
      \exp\{-\kappa N^2+C_nN\log(N+1)\},
 \qquad N>n,\quad N\equiv n\pmod2,
 \quad0\le j\le k.
\end{equation}
The bound is summable in $N$, uniformly as $\eps\downarrow0$. Hence
\begin{equation}\label{eq:alltail}
 \sum_{\substack{N>n\\N\text{ of the selected parity}}}
 |\partial_u^jC_N(\ii\eps)|=O_j(1),\qquad0\le j\le k.
\end{equation}
The same summable majorant controls every higher particle order
of the selected parity and every derivative through order $k$.

\subsection{The first amplitude of the complete form factor}
Choose $\phi=2\pi a/p_0$ with $p_0$ an odd prime and
$1\le a\le(p_0-1)/2$. Then the first odd resonance is $n=p_0$ and
the first even resonance is $n=2p_0$. For all lower indices of the
selected parity, the only possible unit-modulus product configuration
has $Z=e^{-\ii N\phi}\ne1$, so those entire form factors are
holomorphic on a neighborhood of the selected point.

The Gaussian Vandermonde period in the actual independent-angle measure
is
\begin{equation}\label{eq:mehta}
 G_n^{\rm Gauss}:=
 \int_{q_n=-\sum_{i<n}q_i}e^{-Q(q)}\Delta(q)^2
                         \dd q_1\cdots\dd q_{n-1}
 =\frac{n!\,\pi^{(n-1)/2}P_n}{\sqrt n},\qquad
 P_n=\prod_{r=1}^{n-1}\frac{r!}{2^r}.
\end{equation}
Here is an elementary evaluation of the Gaussian period. The monic
polynomials
\[
 H_r(x)=(-1)^r2^{-r}e^{x^2}\frac{\dd^r}{\dd x^r}e^{-x^2}
\]
are orthogonal for $e^{-x^2}\dd x$: integrating by parts $r$ times
against a polynomial of degree less than $r$ gives zero. The same
calculation against $H_r$ gives
$\int H_r(x)^2e^{-x^2}\dd x=2^{-r}r!\sqrt\pi$.
Since $\bigl(\det[H_{r-1}(x_i)]_{r,i=1}^n\bigr)^2=\Delta(x)^2$, expansion of
the two determinants and orthogonality give the ordinary
$n$-variable integral $n!\pi^{n/2}P_n$.
Put $x_i=q_i+m$ with $\sum q_i=0$. The coordinate Jacobian is $n$,
while integration in $m$ gives $\sqrt{\pi/n}$. Dividing by their
product proves \eqref{eq:mehta}. In particular, replacing the actual independent
coordinate measure by hyperplane surface measure would introduce an
incorrect factor $\sqrt n$.

Write $s_\phi=\sin\phi$ and $\alpha=(n^2-1)/2$. The singular part
of the complete first form factor, on its unique collision chart, has
leading expression
\begin{equation}\label{eq:leadingintegral}
 \frac{2\ii^{1-n^2}B(0)^{n^2/2}}
 {n!(2\pi)^{n-1}2^{n(n-1)}s_\phi^{n^2-1}}
 \int\frac{\Delta(q)^2}{nu+(B(0)/2)Q(q)}
                                   \dd q_1\cdots\dd q_{n-1}.
\end{equation}
This follows directly from
\[
 \gamma(q;u)=\ii\phi-\frac{\ii}{s_\phi}
                  \left(u+\frac{B(0)}2q^2\right)
                  +O(u^2+u q^2+q^4)
\]
and the sine divided differences in the numerator. The collision
chart is unique, and its complement in the fixed-$n$ torus integral
is analytic.

Differentiate \eqref{eq:leadingintegral} to the specified order and
scale $q=\sqrt\eps\,v$ on $u=\ii\eps$. The coefficient of $Q$ in
$\gamma$ has positive real part, so the Gaussian/Laplace period is
well-defined with its prescribed branch. The radial tails after the
top derivative have radial integrands $O(R^{-3})$ for odd $n$ and
$O(R^{-2})$ for even $n$; their integrated tails are respectively
$O(R^{-2})$ and $O(R^{-1})$. This gives dominated convergence.
The error estimate can also be made directly before scaling.
After factoring out $\Delta(q)^2$, the analytic numerator differs
from its center value by $O(|u|+Q)$, and the pole denominator differs
from $nu+B(0)Q/2$ by $O((|u|+Q)^2)$. On $u=\ii\eps$ both
denominators have modulus at least $c(\eps+Q)$ in a sufficiently
small fixed chart. Differentiating the exact difference of the two
quotients $k$ times therefore bounds its absolute value by
$C\Delta(q)^2(\eps+Q)^{-k}$. Its radial integral is
$O(1+|\log\eps|)$ for odd $n$ and $O(1)$ for even $n$.
Multiplication by $\eps$ or $\sqrt\eps$, respectively, sends this
error to zero. The off-chart terms are holomorphic on a fixed
neighborhood. This proves the differentiated asymptotic without
an unbounded formal expansion.
Using \eqref{eq:mehta} and the Laplace integral yields
\begin{align}
 \lim_{u\to0,\ u=\ii\eps}u\,\partial_u^kC_n(u)
 &=\frac{2^{1-k}\sqrt{B(0)}P_n n^{k-3/2}}
        {\pi^{(n-1)/2}\sin^{2k}\phi}\ne0,
 &k&=\frac{n^2-1}{2},\quad n\text{ odd},\label{eq:oddfirst}\\
 \lim_{u\to0,\ u=\ii\eps}u^{1/2}\partial_u^kC_n(u)
 &=-\ii\frac{2^{1-\alpha}\sqrt{B(0)}P_n n^{k-1}}
        {\pi^{(n-2)/2}\sin^{n^2-1}\phi}\ne0,
 &k&=\frac{n^2}{2}-1,\quad n\text{ even}.\label{eq:evenfirst}
\end{align}
Here $u^{1/2}=e^{\ii\pi/4}\sqrt\eps$, and $\sqrt{B(0)}$ is fixed
by the beta-period continuation from
$\Re(-\ii B(0)/(2s_\phi))>0$. For even $n$ this sign is determined
by the period; it is not an arbitrary extra square-root branch in the
original integer power of $B$.

For the odd physical HT coefficients the square-root convention is
particularly explicit. On $\OH$, both $\cos y$ and $\cos(\lambda y)$
are positive, and $\sin y$ and $\sin(\lambda y)$ have the same sign.
Consequently
\[
 \Re(T\overline S)=\sinh x\sinh(\lambda x)\cos y\cos(\lambda y)
 +\cosh x\cosh(\lambda x)\sin y\sin(\lambda y)>0
\]
away from the origin, where $B$ has its positive removable value
$\lambda$. Thus $\Re B>0$ throughout the HT lens; its physical square
root is the principal branch. The same is true after axis exchange,
since $\Re(1/B)>0$. This fixes the odd amplitudes on both transfer-axis
charts without an additional sign assumption.

The lower derivatives of the even first term are bounded. For the odd
first term, the $(k-1)$st derivative can be logarithmic, and all lower
derivatives are still $o(u^{-1})$. Together with
\eqref{eq:alltail} this proves that the singularity belongs to the
complete sum, not merely to one term. Multiplication by the analytic
physical prefactor and by $\beta=z/(2J_x)$ preserves it: both are
nonzero at a regular selected point, and the product-rule contributions
from lower derivatives are smaller. A nonramified local change from
$u$ to $z$ also preserves the nonzero top coefficient.

\subsection{All low-temperature signs; the high-temperature distinction}
At a negative endpoint put $q=\pi+q'$. Then
$W=b+u+(-B)(1-\cos q')$. For even $N$, the constraint
$\sum q_i=0\pmod{2\pi}$ remains exactly the same constraint for
$q'_i$, and $B^{N^2/2}=(-B)^{N^2/2}$ because $N^2/2$ is even.
Thus the whole even argument, including its tail and amplitude, applies
to both endpoint signs. The other pair of edges is handled by exchanging
the axes in the symmetric double contour. Rational phase values of the
form above are dense on every regular arc. Section~\ref{sec:geometry}
then supplies the density needed on all of $\partial\OL$.

For odd $N$, the negative endpoint has a different constraint:
$\sum_i(q_i-\pi)$ is an odd multiple of $\pi$. In particular, with
representatives $\delta_i=q_i-\pi\in[-\pi,\pi]$,
\begin{equation}\label{eq:momentummismatch}
 \sum_i\delta_i^2\ge\frac{\pi^2}{N}.
\end{equation}
There is no exact odd all-endpoint collision there. The same Gaussian
numerator bound and \eqref{eq:momentummismatch} imply that the sum of
every fixed parameter derivative is bounded as the negative edge is
approached. This proves existence of all boundary jets, but gives no
positive radius for their Taylor series. The following two sections
resolve this separate infinite-index issue by proving zero radius at
a dense rational-phase family.



\section{Beta periods on the original real cycle at a negative HT edge}
\label{sec:realcycle}

At a negative endpoint the odd momentum constraint displaces the closest
collision by $\pi/n$.  We keep the actual real integration cycle and
compare its exact pole with its quadratic Taylor polynomial.  The
Vandermonde angular measure supplies the needed control of the cubic and
quartic errors.  The angular estimate reuses the constrained Gaussian period,
ordinary gamma moments, and the arithmetic--geometric mean inequality.

\subsection{The exact pole and its direct Taylor estimate}

Fix a regular negative endpoint and exchange the transfer axes if needed.
Use the nonramified coordinate
\begin{equation}\label{eq:rc-coordinate}
 t=\frac{C+B-b}{B},\qquad C=A/S,\qquad
 W(\pi+\delta;t)=b+B(t)(t-1+\cos\delta),\qquad b=\cos\phi.
\end{equation}
After conjugation if necessary, assume
\begin{equation}\label{eq:rc-phase}
 0<\phi<\pi,\qquad \Im B(0)<0,\qquad
 \phi=2\pi r/p,\qquad (r,p)=1,\quad p\ge3\text{ odd}.
\end{equation}
The function $B$ is holomorphic and nonzero on a fixed disk.  Set
$s=\sin\phi$ and choose the holomorphic branch
\[
 g(t,\delta)=\operatorname{arcosh}W(\pi+\delta;t),\qquad
 g(0,0)=\ii\phi.
\]
It is even in $\delta$, and $g_t$ is bounded away from zero on a
sufficiently small fixed bidisk.  Constants below may depend on these
fixed local data, but not on the particle number or Taylor order.

Let $n$ tend to infinity through odd multiples of $p$, and put
\begin{gather}
 d_n=\frac\pi n,\qquad \tau_n=1-\cos d_n,\qquad
 B_n=B(\tau_n),\qquad \alpha_n=\frac{n^2-1}{2},\label{eq:rc-centres}\\
 h_n=\cos d_n-\frac{B_n b}{s^2}\sin^2d_n=1+O(n^{-2}),\qquad
 a_n=\Re h_n>0.\label{eq:rc-h}
\end{gather}
In the first winding write $\delta_i=d_n+v_i$, with
\[
 V_n=\{v\in\C^n:\sum_i v_i=0\},\qquad Q(v)=\sum_i v_i^2.
\]
The real integration measure is $\dd v_1\cdots\dd v_{n-1}$ with
$v_n=-\sum_{i<n}v_i$.  When $v$ is real, write $Q=Q(v)$ and
$\Delta(v)=\prod_{i<k}(v_i-v_k)$.  The inner region
$\sum_i\delta_i^2<\sigma$ is exactly the real ball
\begin{equation}\label{eq:rc-actual-ball}
 Q<R_n^2,\qquad R_n^2=\sigma-\frac{\pi^2}{n}.
\end{equation}
For fixed small $\sigma>0$ and large $n$, these squared radii lie
between $\sigma/2$ and $\sigma$.

The relevant pole is the unique solution of
\begin{equation}\label{eq:rc-pole-equation}
 \sum_i\bigl(g(T_n(v),d_n+v_i)-\ii\phi\bigr)=0.
\end{equation}
It is defined and holomorphic on a common small ball in $V_n$.
Indeed, after the common constant and linear terms in $v_i$ are
subtracted, the remaining sum is bounded by $C\sum_i|v_i|^2$;
the parameter derivative of the left side has modulus at least $cn$.
Rouch\'e's theorem on a circle of radius
$C\sum_i|v_i|^2/n$ about $\tau_n$, followed by uniqueness and the
implicit-function theorem, gives
\begin{equation}\label{eq:rc-pole-coarse}
 T_n(0)=\tau_n,\qquad
 |T_n(v)-\tau_n|\le\frac Cn\sum_i|v_i|^2.
\end{equation}
All scalar analytic bounds come from the original fixed bidisk.

Direct differentiation at the center gives
\[
 g(\tau_n,d_n)=\ii\phi,\qquad
 g_t(\tau_n,d_n)=-\frac{\ii B_n}{s},\qquad
 \frac12g_{\delta\delta}(\tau_n,d_n)=\frac{\ii B_nh_n}{2s}.
\]
Evenness in $\delta$ gives $g_{\delta\delta\delta}(t,d_n)=O(d_n)$.
Writing $u=T_n-\tau_n$, Taylor expansion of
\eqref{eq:rc-pole-equation}, with the linear terms canceled by
$\sum v_i=0$, therefore gives on the real ball
\[
 0=ng_t(\tau_n,d_n)u+\tfrac12g_{\delta\delta}(\tau_n,d_n)Q
 +O\left(n|u|^2+|u|Q+d_n\sum_i|v_i|^3+\sum_i v_i^4\right).
\]
Using \eqref{eq:rc-pole-coarse}, we obtain the following estimate:
\begin{equation}\label{eq:rc-pole-error}
 \begin{split}
 T_n(v)&=T_{0,n}(v)+\mathcal R_n(v),\qquad
 T_{0,n}(v)=\tau_n+\frac{h_nQ}{2n},\\
 |\mathcal R_n(v)|&\le\frac Cn\left(
       \frac1n\sum_i|v_i|^3+\sum_i v_i^4+\frac{Q^2}{n}\right).
 \end{split}
\end{equation}
In particular, since $\sum|v_i|^3\le Q^{3/2}$ and
$\sum v_i^4\le Q^2$, a sufficiently small fixed $\sigma$ gives
\begin{equation}\label{eq:rc-small-error}
 |\mathcal R_n(v)|\le\tfrac14D_n(v),\qquad
 D_n(v):=\tau_n+\frac{a_nQ}{2n}\le |T_{0,n}(v)|
 \quad(Q<\sigma).
\end{equation}

\subsection{The complete residue density}

Include every normalization in the numerator
\begin{equation}\label{eq:rc-numerator}
 F_n(t,v)=\frac{B(t)^{n^2/2}}{n!(2\pi)^{n-1}}
 \prod_i\frac1{\sinh g(t,d_n+v_i)}
 \prod_{i<k}\left[
 \frac{\sin((v_i-v_k)/2)}
 {\sinh((g(t,d_n+v_i)+g(t,d_n+v_k))/2)}\right]^2.
\end{equation}
The square root of $B$ is its physical one.  The coefficient of the
exact parameter pole is
\begin{equation}\label{eq:rc-residue}
 I_{n,j}=\int_{Q<R_n^2}\frac{\rho_n(v)}{T_n(v)^{j+1}}\dd v,
 \qquad
 \rho_n(v)=-\frac{2F_n(T_n(v),v)}
 {\sum_i g_t(T_n(v),d_n+v_i)}.
\end{equation}
The minus sign comes from expansion of $(t-T_n)^{-1}$ at $t=0$.
The ratio $\rho_n/\Delta^2$ is holomorphic and nonzero near the
center, where its value is
\begin{equation}\label{eq:rc-center-density}
 \rho_{n,0}=-\frac{2B_n^{\alpha_n-1/2}}
 {n\,n!(2\pi)^{n-1}2^{n(n-1)}s^{2\alpha_n}}.
\end{equation}
Indeed, the numerator divided by $\Delta^2$ at the center is
$-\ii B_n^{n^2/2}/[n!(2\pi)^{n-1}2^{n(n-1)}s^{n^2}]$,
and the phase derivative is $-\ii nB_n/s$.

\begin{lemma}[Direct full-density estimate]\label{lem:rc-density}
On the real ball $Q<\sigma$,
\begin{equation}\label{eq:rc-density-log}
 \rho_n(v)=\rho_{n,0}\Delta(v)^2e^{\ell_n(v)},\qquad
 |\ell_n(v)|\le C(\sqrt{nQ}+nQ).
\end{equation}
The logarithm is holomorphic through all coincidence divisors.
\end{lemma}
\begin{proof}
Divide every analytic factor by its nonzero center value before
taking logarithms.  The factor $B(T_n)^{n^2/2}$ costs
$Cn^2|T_n-\tau_n|\le CnQ$.  Also
\[
 \sum_i|g(T_n,d_n+v_i)-\ii\phi|
 \le C\left(n|T_n-\tau_n|+d_n\sum_i|v_i|+Q\right).
\]
Each phase occurs in $n-1$ pair denominators, so their total
logarithmic cost is at most $C(nQ+\sqrt{nQ})$.
The sine divided differences cost
$C\sum_{i<k}(v_i-v_k)^2=CnQ$, and the one-body denominators have
the smaller single-index bound.  Finally, cancellation of linear
terms in $\sum_i g_t(T_n,d_n+v_i)$ gives its ratio to
$ng_t(\tau_n,d_n)$ as $1+O(Q/n)$.  These estimates prove the claim
for the complete residue density.
\end{proof}

\subsection{An angular estimate from the Gaussian period}

Separate radius and angle in the positive real measure
$\Delta(v)^2\dd v$ on the zero-sum plane.  Write
$v=\sqrt Q\,\theta$, where $\sum\theta_i=0$ and
$\sum\theta_i^2=1$, and let $\nu_n$ be the normalized angular
measure proportional to $\Delta(\theta)^2$.

\begin{lemma}[Uniform angular exponential moment]\label{lem:rc-angular}
For $Z_n(\theta)=n\max_i\theta_i^2$, there are absolute
$c_0,C_0>0$ such that
\begin{equation}\label{eq:rc-angular-moment}
 \sup_{n\ge3}\mathbb E_{\nu_n}
       \bigl[(1+Z_n)e^{c_0Z_n}\bigr]\le C_0.
\end{equation}
One may take $c_0=1/32$.
\end{lemma}
\begin{proof}
Use the Gaussian probability measure
\[
 \frac1{\mathcal G_n}e^{-Q(v)}\Delta(v)^2\dd v,
 \qquad
 \mathcal G_n=\frac{n!\pi^{(n-1)/2}P_n}{\sqrt n},\qquad
 P_n=\prod_{r=1}^{n-1}\frac{r!}{2^r},
\]
whose normalization is \eqref{eq:mehta}.  Homogeneity shows that
$Q$ and $\theta$ are independent, $\theta$ has law $\nu_n$, and
$Q$ has the gamma density
$q^{\alpha_n-1}e^{-q}/\Gamma(\alpha_n)$, where
$\alpha_n=(n^2-1)/2$.

We first bound the one-coordinate marginal density $f_n(x)$.
Fix $v_1=x$, put $m=n-1$ and $c=n/(n-1)$, and write
\[
 v_{i+1}=u_i-x/m\quad(1\le i\le m),\qquad \sum_i u_i=0.
\]
In the actual independent coordinates
$(x,u_1,\ldots,u_{m-1})$, the Jacobian is one.  Moreover
\[
 Q(v)=cx^2+Q(u),\qquad
 \Delta(v)^2=\Delta(u)^2\prod_{i=1}^m(cx-u_i)^2.
\]
Consequently, with $\mathbb E_m$ denoting the analogous constrained
Gaussian measure in $m$ variables,
\begin{equation}\label{eq:rc-gaussian-marginal}
 f_n(x)=\frac{\mathcal G_{n-1}}{\mathcal G_n}e^{-cx^2}
       \mathbb E_m\prod_{i=1}^m(cx-u_i)^2,
 \qquad
 \frac{\mathcal G_{n-1}}{\mathcal G_n}
 =\frac{2^{n-1}}{\sqrt\pi\,n!}\sqrt{\frac n{n-1}}.
\end{equation}
The arithmetic--geometric mean inequality and $\sum_i u_i=0$
give
\[
 \prod_{i=1}^m(cx-u_i)^2
 \le\left(c^2x^2+\frac{Q(u)}m\right)^m.
\]
Here $Q(u)$ is gamma with shape $\alpha_m=(m^2-1)/2$.
For $0\le k\le m$ its moments satisfy
\[
 \mathbb E_m\left(\frac{Q(u)}m\right)^k
 =\frac{\alpha_m(\alpha_m+1)\cdots(\alpha_m+k-1)}{m^k}
 \le n^k,
\]
because $(\alpha_m+m-1)/m\le n$.  Since $c^2\le9/4<3$,
the binomial expansion therefore yields
\begin{equation}\label{eq:rc-marginal-product-bound}
 \mathbb E_m\prod_{i=1}^m(cx-u_i)^2
 \le (3x^2+n)^{n-1}.
\end{equation}
Use $n!\ge(n/e)^n$, $c\ge1$, and
$\sqrt{n/(n-1)}/(2\sqrt\pi)<1$.  With $t=x^2/n$,
\eqref{eq:rc-gaussian-marginal}--\eqref{eq:rc-marginal-product-bound}
give
\[
 f_n(x)\le\frac1n(2e)^n(1+3t)^{n-1}e^{-x^2}.
\]
For $t\ge16$,
$\log(2e)+\log(1+3t)\le t/2$: it holds at $16$, and the
derivative of the right side minus the left side is positive
thereafter.  Thus
\begin{equation}\label{eq:rc-coordinate-tail-density}
 f_n(x)\le\frac1n e^{-x^2/2}\qquad(x^2\ge16n).
\end{equation}
The Gaussian measure is invariant under coordinate permutations;
the independent-coordinate measure is a constant multiple of
Euclidean measure on the zero-sum plane.  Hence a union bound and
\eqref{eq:rc-coordinate-tail-density} imply, for $A^2\ge16n$,
\begin{equation}\label{eq:rc-coordinate-union-tail}
 \Pr\left(\max_i v_i^2\ge A^2\right)
 \le2\int_A^\infty e^{-x^2/2}\dd x
 \le\frac2A e^{-A^2/2}.
\end{equation}

To pass to the angular variable, the inverse gamma moment gives
\[
 \Pr\left(Q<\frac{n^2}8\right)
 \le\frac{n^2}8\mathbb E\frac1Q
 =\frac{n^2}{4(n^2-3)}\le\frac38\qquad(n\ge3).
\]
Independence of radius and angle, and
$\max_i v_i^2=QZ_n/n$, consequently give
\[
 \frac58\Pr(Z_n\ge z)
 \le\Pr\left(\max_i v_i^2\ge\frac{nz}8\right).
\]
For $z\ge128$, apply \eqref{eq:rc-coordinate-union-tail} with
$A^2=nz/8\ge16n$ and $A\ge1$ to conclude
\begin{equation}\label{eq:rc-angular-tail}
 \Pr(Z_n\ge z)\le\frac{16}5e^{-nz/16},\qquad z\ge128.
\end{equation}
Integrating this tail proves \eqref{eq:rc-angular-moment}, for
example with $c_0=1/32$, uniformly in $n\ge3$.
\end{proof}

The precise consequences we need are
\begin{equation}\label{eq:rc-angular-power-bounds}
 \sum_i v_i^4\le\frac{Q^2Z_n}{n},\qquad
 \sum_i|v_i|^3
 \le\sqrt{Q\sum_i v_i^4}
 \le\frac{Q^{3/2}\sqrt{Z_n}}{\sqrt n}.
\end{equation}
These inequalities are pointwise; the exponential moment handles
all angular directions, including those where $Z_n$ is large.

\subsection{The exact complex quadratic period}

Write
\[
 P_n=\prod_{r=1}^{n-1}\frac{r!}{2^r},\qquad
 \mathcal G_n=\int_{V_n\cap\R^n}e^{-Q}\Delta(v)^2\dd v
 =\frac{n!\pi^{(n-1)/2}P_n}{\sqrt n}.
\]
This is the constrained Gaussian period \eqref{eq:mehta}, in the
actual independent-coordinate measure.  Homogeneity gives radial
measure $\mathcal G_nQ^{\alpha_n-1}\dd Q/\Gamma(\alpha_n)$.
For $j+1>\alpha_n$ and $\Re h>0$, therefore,
\begin{equation}\label{eq:rc-complex-beta}
 J_h:=\int_{V_n\cap\R^n}
   \frac{\Delta(v)^2\dd v}{(\tau_n+hQ/(2n))^{j+1}}
 =\mathcal G_n(2n)^{\alpha_n}h^{-\alpha_n}
   \tau_n^{\alpha_n-j-1}
   \frac{\Gamma(j+1-\alpha_n)}{\Gamma(j+1)}.
\end{equation}
One proof uses the scalar Laplace formula for the negative power,
followed by the radial gamma integral; absolute convergence follows
from $\Re h>0$.  All integrations remain over the original real
variables.  The gamma ratio in \eqref{eq:rc-complex-beta} is exact.

Use $a_n=\Re h_n$ for the positive majorant.  The corresponding
period satisfies
\begin{equation}\label{eq:rc-phase-cost}
 \frac{J_{a_n}}{|J_{h_n}|}
 =\left(\frac{|h_n|}{a_n}\right)^{\alpha_n}
 =\exp(O(n^{-2})).
\end{equation}
Indeed, $a_n=1+O(n^{-2})$ and $\Im h_n=O(n^{-2})$, so the
logarithm in parentheses is $O(n^{-4})$. This comparison bounds
the complex quadratic period by its positive majorant on the same
real integration cycle.

\subsection{Uniform comparison and the two first windings}

We need uniformity only over
\begin{equation}\label{eq:rc-needed-window}
 \log j<n<\frac{\sqrt j}{(\log j)^{1/4}}.
\end{equation}
In this window $n\to\infty$ and $\varepsilon_n:=n^2/j\to0$
uniformly.  Let $I_{n,j}$ be the exact first-winding pole coefficient
in \eqref{eq:rc-residue}.

\begin{proposition}[Uniform estimate on the original real cycle]
\label{prop:negative-real-cycle-coefficients}
Uniformly for the odd multiples of $p$ in
\eqref{eq:rc-needed-window},
\begin{equation}\label{eq:rc-one-winding-model}
 I_{n,j}=\rho_{n,0}J_{h_n}
           \left(1+O\left(\frac n{\sqrt j}\right)\right).
\end{equation}
The sum of the two first-winding pole coefficients is consequently
\begin{equation}\label{eq:rc-two-winding-model}
 M(n,j)\left(1+O\left(\frac n{\sqrt j}\right)\right),
\end{equation}
where the complete exact model is
\begin{equation}\label{eq:rc-model}
 \boxed{\displaystyle
 M(n,j)=-\frac{2^{2-\alpha_n}B_n^{\alpha_n-1/2}P_n
                    n^{\alpha_n-3/2}}
 {\pi^{(n-1)/2}s^{2\alpha_n}}
 h_n^{-\alpha_n}\tau_n^{\alpha_n-j-1}
 \frac{\Gamma(j+1-\alpha_n)}{\Gamma(j+1)}.}
\end{equation}
\end{proposition}
\begin{proof}
Normalize the positive measure associated with $J_{a_n}$ and set
\[
 x=\frac{a_nQ}{2n\tau_n},\qquad \varepsilon=\frac{n^2}{j}.
\]
Radius and angle are independent.  The angular law is $\nu_n$,
and the radial density is proportional to
$x^{\alpha_n-1}(1+x)^{-j-1}$.  Since $2n^2\tau_n$ and $a_n$
are bounded above and below, $Q\asymp x/n$ uniformly.  Equations
\eqref{eq:rc-pole-error} and \eqref{eq:rc-angular-power-bounds}
give
\begin{equation}\label{eq:rc-weighted-pole-error}
 (j+1)\frac{|\mathcal R_n|}{D_n}
 \le C\frac j{n^2}
   \frac{x^{3/2}\sqrt{Z_n}+x^2(Z_n+1)}{1+x}.
\end{equation}
Moreover \eqref{eq:rc-small-error} implies
\begin{equation}\label{eq:rc-integrand-error}
 \left|\frac{e^{\ell_n}}{T_n^{j+1}}
                 -\frac1{T_{0,n}^{j+1}}\right|
 \le D_n^{-j-1}
 \left[\exp\left\{|\ell_n|+
               \frac{2(j+1)|\mathcal R_n|}{D_n}\right\}-1\right].
\end{equation}
This follows by the convergent logarithm of
$1+\mathcal R_n/T_{0,n}$, whose modulus is at most
$2|\mathcal R_n|/D_n$.

Fix $K=32$.  On $x\le K\varepsilon$,
\eqref{eq:rc-density-log} and \eqref{eq:rc-weighted-pole-error}
bound the exponent in \eqref{eq:rc-integrand-error} by
\[
 C\{\sqrt\varepsilon(1+\sqrt{Z_n})+
                       \varepsilon(Z_n+1)\}
 \le C'\sqrt\varepsilon(1+Z_n).
\]
For sufficiently large $j$, Lemma~\ref{lem:rc-angular} and the
mean-value theorem for the exponential show that its angular
expectation, after subtracting one, is $O(\sqrt\varepsilon)$.
This estimates the whole angular measure in the central radial
region.

On the actual ball, $x\le Ln$ for a fixed $L$.  The deterministic
bound $Z_n\le n$ in \eqref{eq:rc-weighted-pole-error} gives
\[
 \frac{2(j+1)|\mathcal R_n|}{D_n}
 \le Cj\frac{x^{3/2}/n^{3/2}+x^2/n}{1+x}.
\]
Divided by $j\log(1+x)$, this tends to zero uniformly for
$0<x\le Ln$: splitting at $1$ and $\sqrt n$ bounds the ratio
by $O(n^{-1})$, $O(n^{-1/2})$, and $O(1/\log n)$,
respectively.  Also $|\ell_n|\le C(\sqrt x+x)$.
On $K\varepsilon\le x\le Ln$,
\[
 \frac{\sqrt x+x}{j\log(1+x)}=o(1)
\]
uniformly.  For $x\le1$, use
$j\sqrt x\ge\sqrt K\,n\sqrt j$; for $x\ge1$, use
$x\le Ln$ and $n/j\to0$.  Hence for all sufficiently large
$j$ in the stated window,
\begin{equation}\label{eq:rc-tail-absorption}
 |\ell_n|+\frac{2(j+1)|\mathcal R_n|}{D_n}
 \le\frac j4\log(1+x)
 \quad(K\varepsilon\le x\le Ln).
\end{equation}

Write
$Z_\beta(\alpha,J)=\int_0^\infty
x^{\alpha-1}(1+x)^{-J-1}\dd x$.
The normalized contribution of this remaining radial region is
therefore bounded by a constant times
\begin{equation}\label{eq:rc-radial-tail}
 (1+K\varepsilon)^{-j/4}
 \frac{Z_\beta(\alpha_n,j/2)}{Z_\beta(\alpha_n,j)}
 \le (1+K\varepsilon)^{-j/4}4^{\alpha_n}
 \le e^{-cn^2}.
\end{equation}
Indeed, $\alpha_n\le j/4$ eventually, and its integrality gives
the exact ratio
\[
 \frac{Z_\beta(\alpha_n,j/2)}{Z_\beta(\alpha_n,j)}
 =\prod_{r=0}^{\alpha_n-1}\frac{j-r}{j/2-r}\le4^{\alpha_n}.
\]
Use $\log(1+K\varepsilon)\ge K\varepsilon/2$ to obtain the
last inequality.  The omitted part of the unperturbed period
outside the actual ball has the same bound: its positive majorant
has no perturbation factor, and the ball endpoint exceeds
$K\varepsilon$ for large $j$.

Combining these estimates and \eqref{eq:rc-phase-cost} gives
\[
 \frac{|I_{n,j}-\rho_{n,0}J_{h_n}|}
 { |\rho_{n,0}J_{h_n}|}
 \le C\left(\frac n{\sqrt j}+e^{-cn^2}\right).
\]
Because $n>\log j$, the exponential term is absorbed into
$n/\sqrt j$, uniformly.  This proves
\eqref{eq:rc-one-winding-model}.

Reflection $\delta\mapsto-\delta$ preserves the full integrand
and independent-angle measure, and exchanges the two first-winding
balls.  They therefore contribute equally.  Multiplication of
\eqref{eq:rc-center-density} by the period
\eqref{eq:rc-complex-beta} and this factor two gives
\eqref{eq:rc-model} exactly.
\end{proof}

The factors $B_n^{\alpha_n-1/2}$ retain their quadratic dependence
of phase on the particle number.  The proposition does not assert
positivity across indices.  The subsequent coefficient comparison
will prove noncancellation for the complete odd sum.  Its low and
high particle ranges use separate absolute estimates, so no relative
asymptotic outside \eqref{eq:rc-needed-window} is needed.

\section{Negative high-temperature edges: boundary jets and the complete odd sum}
\label{sec:jets}

The absence of an odd-particle endpoint resonance at a negative spectral
edge does not imply continuation of the infinite odd sum. We prove its
nonanalyticity directly from its boundary Taylor coefficients,
defined by inward limits on the physical side.

\subsection{Setting and the existence of the boundary jets}

Fix a regular negative endpoint. After complex conjugation if necessary,
the local analytic data can be written
\begin{equation}\label{eq:jets_setting}
 W(t,\pi+\delta)=b+B(t)\bigl[t-(1-\cos\delta)\bigr],\qquad
 b=\cos\phi,\quad \phi=\frac{2\pi r}{p}\in(0,\pi),
\end{equation}
where $p\ge3$ is odd, $(r,p)=1$, $B$ is holomorphic and nonzero,
$\Im B(0)<0$, and $t$ is a nonramified local temperature coordinate.
The physical inward direction is $t=-\eps$, $\eps>0$. The branch
$\cosh\gamma(q;t)=W(t,q)$ has
$\gamma(\pi;0)=\ii\phi$. All assertions below concern this branch.
At a regular negative edge, the earlier root construction gives a common
branchfree parameter neighborhood for the entire real angular circle.

For clarity, write the complete coefficient as
\begin{equation}\label{eq:jets_full_coefficient}
 C_n(t)=\int_{\mathcal T_n}H_n(t,q)
        \coth\!\left(\frac{\Gamma_n(t,q)}2\right)\dd q,
 \qquad \Gamma_n=\sum_{i=1}^n\gamma(q_i;t),
\end{equation}
where $\mathcal T_n$ is the zero-total-momentum torus,
$q_n=-\sum_{i<n}q_i$ modulo $2\pi$, the measure $\dd q$ is the
ordinary measure in the $n-1$ independent angles, and
\begin{equation}\label{eq:jets_numerator}
 H_n(t,q)=\frac{B(t)^{n^2/2}}{n!(2\pi)^{n-1}}
 \prod_{i=1}^n\frac1{\sinh\gamma(q_i;t)}
 \prod_{i<k}\left[
  \frac{\sin((q_i-q_k)/2)}
       {\sinh((\gamma(q_i;t)+\gamma(q_k;t))/2)}
 \right]^2.
\end{equation}
The branch of $B^{1/2}$ is the one in the physical representation. The
normalization in \eqref{eq:jets_numerator} is retained throughout this
section. The physical susceptibility differs from the complete odd sum
by the analytic prefactor already identified in the physical
representation; that prefactor is nonzero at a regular endpoint.

Set
\[
 S(q)=\sum_{i=1}^n\dist(q_i,\pi)^2.
\]
The spectral geometry and the matched Schur--Pfaffian estimate imply,
with constants independent of $n$,
\begin{align}
 \Re\gamma(q;0)&\ge c_0\dist(q,\pi)^2,
 &|\gamma(q;0)-\ii\phi|&\le C\dist(q,\pi)^2,\label{eq:jets_attenuation0}\\
 \sup_{|t|<r_0,\ q\in\mathcal T_n}|H_n(t,q)|
 &\le E_n,
 &E_n&=\exp\{-\kappa n^2+C n\log(n+1)+C\}.
 \label{eq:jets_gaussian_num}
\end{align}
Shrinking $r_0$ is allowed. In particular, the strict Schur contraction
used in \eqref{eq:jets_gaussian_num} persists on this fixed complex
parameter disk; it does not require $|e^{-\gamma}|\le1$ at every
complex parameter in the disk.

For odd $n$, choose the shortest lifts $q_i=\pi+\delta_i$ modulo
$2\pi$. Momentum conservation gives
$\sum_i\delta_i=w\pi$ with $w$ odd. Therefore
\begin{equation}\label{eq:jets_momentum_minimum}
 S(q)\ge\frac{\pi^2}{n}.
\end{equation}
This is the source of an $n^{-2}$ coefficientwise analytic disk, rather
than a common disk for the entire particle sum.

\begin{lemma}[Coefficientwise disks and boundary derivatives]
\label{lem:jets_boundary_derivatives}
There is $c>0$ such that each odd coefficient is holomorphic on
$|t|<c/n^2$, agrees there with \eqref{eq:jets_full_coefficient}, and
obeys
\begin{equation}\label{eq:jets_small_disk}
 \sup_{|t|<c/n^2}|C_n(t)|\le CnE_n,
 \qquad |[t^j]C_n(t)|\le (Cn^2)^j CnE_n.
\end{equation}
The complete odd sum on $t<0$ has derivatives of every fixed order at
$t=0$ from that side, and its normalized boundary jets are
\begin{equation}\label{eq:jets_boundary_coefficients}
 a_j=\lim_{\eps\downarrow0}\frac1{j!}
       \partial_t^j\sum_{n\ \mathrm{odd}}C_n(-\eps)
     =\sum_{n\ \mathrm{odd}}[t^j]C_n(t).
\end{equation}
The sum on the right converges absolutely for each fixed $j$.
\end{lemma}

\begin{proof}
The parameter derivative of each one-particle phase is bounded on a
fixed disk. Thus, using \eqref{eq:jets_attenuation0} and
\eqref{eq:jets_momentum_minimum}, a sufficiently small perturbation
$|t|<c/n^2$ preserves $\Re\Gamma_n\ge c'/n$. The modulus of the
hyperbolic cotangent is then at most $Cn$. This proves the first bound;
Cauchy's estimate proves the second.

For the inward real approach, \eqref{eq:jets_setting} gives
\[
 \Re\gamma(q;-\eps)\ge
 c\bigl(\eps+\dist(q,\pi)^2\bigr).
\]
Indeed $\Im W=(-\Im B(-\eps))(\eps+1-\cos\delta)$ is positive,
and writing $\gamma=v+\ii\psi$ expresses it as
$\sinh v\sin\psi$, with $\psi$ in a compact subinterval of
$(0,\pi)$. A disk centered at $-\eps$ of radius
$c(\eps+n^{-2})$ consequently preserves the same denominator gap,
up to a fixed factor. Cauchy's estimate there bounds the $j$th
derivative by $j!(Cn^2)^j CnE_n$, uniformly for small $\eps$.
This is summable for every fixed $j$. Differentiation, the boundary
limit, and the particle sum therefore commute by dominated
convergence. This argument proves existence of all boundary jets; it
does not assert a common holomorphic disk for their sum.
\end{proof}

\subsection{Removing all analytic remainders}

Choose $\sigma>0$ sufficiently small, and then a fixed
$c_{\rm out}>0$ such that
$C_\Gamma c_{\rm out}<c_0\sigma/2$ and
$C_\Gamma(\sigma+c_{\rm out})<\pi/p$.
Here $C_\Gamma$ bounds the phase increments below. On
$|t|<c_{\rm out}/n$,
\begin{equation}\label{eq:jets_c_over_n_estimate}
 \Re\Gamma_n\ge c_0S-C_\Gamma c_{\rm out},\qquad
 |\Gamma_n-n\ii\phi|\le C_\Gamma(S+c_{\rm out}).
\end{equation}
If $S\ge\sigma$, the first estimate supplies a fixed positive real
part by this choice of $c_{\rm out}$. If $S<\sigma$ and $p\nmid n$, the
second estimate supplies a fixed phase gap: the distance from
$n\phi$ to $2\pi\Z$ is at least $2\pi/p$. Thus the whole
coefficient for $p\nmid n$, and the part with $S\ge\sigma$ for any
$n$, are holomorphic on this $c_{\rm out}/n$ disk with norm at most $CE_n$.
The angular cutoff is fixed and independent of $t$.

For $p\mid n$, put $D_n=\Gamma_n-n\ii\phi$. On the inner region
$S<\sigma$, $|D_n|<\pi$ on the same disk if the constants are
sufficiently small. The identity
\begin{equation}\label{eq:jets_coth_subtraction}
 \coth(D_n/2)=\frac2{D_n}+R(D_n),
 \qquad R \text{ holomorphic on } |D_n|<2\pi,
\end{equation}
places the $R$ part in the same analytic remainder class.

We also subtract the exact simple parameter pole, not only its
quadratic approximation. In the inner region use the unique small
lifts $q_i=\pi+\delta_i$ and put $v_i=1-\cos\delta_i$. Let
\[
 d_n(t,q)=\frac1nD_n(t,q),\qquad
 a_0=\frac{B(0)}{\sinh(\ii\phi)}\ne0.
\]
On a sufficiently small fixed parameter disk,
\[
 \partial_t d_n(t,q)=a_0+O(|t|+\sigma),
\]
uniformly in $n$ and in $S<\sigma$. Reducing the disk and $\sigma$
makes this derivative lie within $|a_0|/4$ of $a_0$. The quantitative
inverse-function estimate gives a unique zero
$t=T_n(q)=O(S/n)$ and the factorization
\begin{equation}\label{eq:jets_exact_pole_factor}
 D_n(t,q)=n(t-T_n(q))A_n(t,q),\qquad
 A_n(t,q)=\int_0^1\partial_t d_n
       (T_n(q)+u(t-T_n(q)),q)\dd u.
\end{equation}
The factor $A_n$ is analytic and uniformly bounded away from zero on
a smaller fixed disk containing every such $T_n(q)$. The proof can
also be obtained directly by Rouch\'e's theorem on a circle of radius
$CS/n$: $|d_n(0,q)|\le CS/n$, while the derivative differs uniformly
from the nonzero constant $a_0$.

Set $F_n=H_n/A_n$. The exact identity
\begin{equation}\label{eq:jets_divided_difference}
 \frac{2H_n(t,q)}{D_n(t,q)}=
 \frac{2F_n(T_n(q),q)}{n(t-T_n(q))}
 +\frac2n\frac{F_n(t,q)-F_n(T_n(q),q)}{t-T_n(q)}
\end{equation}
has a holomorphic second term, including at $t=T_n(q)$. A
fixed-disk Cauchy bound on $\partial_tF_n$ bounds this divided
difference by $CE_n$. In particular no lower bound on the distance
between the pole and a chosen integration contour is being assumed.

\begin{lemma}[Exact reduction to the inner pole integrals]
\label{lem:jets_remainder}
For every $j\ge0$ the boundary coefficient in
\eqref{eq:jets_boundary_coefficients} is the sum of the exact inner
pole coefficients for odd $n$ divisible by $p$, together with a
remainder $r_j$ satisfying
\begin{equation}\label{eq:jets_remainder_bound}
 |r_j|\le \exp\left\{\frac12j\log(j+1)+Cj+C\right\}.
\end{equation}
For a resonant $n$, the inner pole coefficient is exactly
\begin{equation}\label{eq:jets_exact_residue_coefficient}
 a_{j,n}^{\mathrm{pole}}=
 -\int_{S<\sigma}
  \frac{2H_n(T_n(q),q)}{\partial_t\Gamma_n(T_n(q),q)}
  T_n(q)^{-j-1}\dd q.
\end{equation}
\end{lemma}

\begin{proof}
All removed terms are bounded on a disk $c_{\rm rem}/n$, where
$c_{\rm rem}>0$ is a fixed sufficiently small fraction of $c_{\rm out}$.
The exact pole subtraction is constructed on its fixed disk before
restriction; the pole need not lie in the exterior disk. The preceding
decompositions are identities on the $c/n^2$ disk where the original
integral is holomorphic. Taking its Taylor
coefficients gives \eqref{eq:jets_exact_residue_coefficient}; the
minus sign is the coefficient of $(t-T)^{-1}$ at $t=0$. Every
removed term is analytic on $|t|<c_{\rm rem}/n$ with norm bounded by $CE_n$.
Their total coefficient bound is
\[
 \sum_{n\ge1}(Cn)^j E_n.
\]
Absorb $Cn\log(n+1)$ into $(\kappa/2)n^2+C'$, and optimize
$j\log n-(\kappa/2)n^2$. This gives
\eqref{eq:jets_remainder_bound}, including summability of the tail.
\end{proof}

\subsection{Uniform bounds for the winding sectors}

The exact pole is quantitatively close to the average
$\bar v=n^{-1}\sum_i v_i$. At $t=\bar v$, the linear terms in the
analytic expansion of
$\operatorname{arcosh}(b+B(\bar v)(\bar v-v_i))$ sum to zero.
The same uniform inverse-function estimate as above consequently gives
\begin{equation}\label{eq:jets_pole_lower}
 T_n(q)=\bar v+O\bigl(\operatorname{Var}(v)\bigr),\qquad
 |T_n(q)|\ge (1-C\sigma)\frac{S(q)}{2n}.
\end{equation}
Here $\operatorname{Var}(v)=n^{-1}\sum_i(v_i-\bar v)^2\le
C\sigma\bar v$, and
$\bar v\ge (1-C\sigma)S/(2n)$. The second estimate concerns the
modulus of a generally complex pole; it does not assert that the
pole is real.

Write $s=\sin\phi>0$, $\alpha_n=(n^2-1)/2$, and
\begin{equation}\label{eq:jets_H0}
 H_{0,n}=\frac{|B(0)|^{n^2/2}}
 {n!(2\pi)^{n-1}2^{n(n-1)}s^{n^2}}.
\end{equation}
The complete numerator, evaluated at the pole, satisfies
\begin{equation}\label{eq:jets_numerator_vandermonde}
 |H_n(T_n(q),q)|\le H_{0,n}\Delta(\delta)^2 e^{CnS},
 \qquad \Delta(\delta)=\prod_{i<k}(\delta_i-\delta_k).
\end{equation}
To check the constant and the dependence on $n$, divide every sine
numerator by $(\delta_i-\delta_k)/2$ and every hyperbolic
denominator by $\sinh(\ii\phi)$. Analytic logarithms of the
resulting near-one ratios are bounded by
$C(\delta_i-\delta_k)^2$ and
$C(|T_n|+\delta_i^2+\delta_k^2)$, respectively. Their sums are
$O(nS+n^2|T_n|)=O(nS)$. The one-body factors cost $O(S)$.
Finally the logarithm of
$B(T_n)^{n^2/2}/B(0)^{n^2/2}$ costs $O(n^2|T_n|)$.
This proves \eqref{eq:jets_numerator_vandermonde} with the exact
scale \eqref{eq:jets_H0}, rather than an unspecified constant raised
to the power $n^2$.

In winding sector $w$, where $w$ is odd, put
\[
 \delta_i=\frac{w\pi}{n}+\eta_i,\quad
 \sum_i\eta_i=0,\quad Q=\sum_i\eta_i^2,\quad
 m_{n,w}=\frac{\pi^2w^2}{n}.
\]
Then $S=Q+m_{n,w}$ and $\Delta(\delta)=\Delta(\eta)$.
There are at most $C\sqrt n$ sectors with nonempty $S<\sigma$.
The constrained Gaussian integral, with the same independent-angle
measure as above, is
\begin{equation}\label{eq:jets_Mehta}
 I_n=\int_{\R^{n-1}}e^{-Q}\Delta(\eta)^2\dd\eta
 =\frac{\pi^{(n-1)/2}2^{-n(n-1)/2}}{\sqrt n}
   \prod_{r=1}^n r!.
\end{equation}
For example, separating the center of mass in the ordinary Gaussian
Vandermonde integral gives a Jacobian $n$ and a center Gaussian
$\sqrt{\pi/n}$, whose product explains the constrained
$1/\sqrt n$ factor. Homogeneity therefore gives, whenever
$j+1>\alpha_n$,
\begin{equation}\label{eq:jets_beta_integral}
 \int_{\R^{n-1}}\frac{\Delta(\eta)^2}{(Q+m)^{j+1}}\dd\eta
 =I_n m^{\alpha_n-j-1}
       \frac{\Gamma(j+1-\alpha_n)}{\Gamma(j+1)}.
\end{equation}

\begin{lemma}[Higher windings and large particle numbers]
\label{lem:jets_windings_large_n}
For $j+1>\alpha_n$, the exact pole coefficient in winding $w$
satisfies
\begin{align}
 |a_{j,n,w}^{\mathrm{pole}}|
 &\le \frac{C H_{0,n}}n e^{Cn\sigma}
 \left(\frac{2n}{1-C\sigma}\right)^{j+1}
 I_n m_{n,w}^{\alpha_n-j-1}
 \frac{\Gamma(j+1-\alpha_n)}{\Gamma(j+1)}.
 \label{eq:jets_winding_bound}
\end{align}
For every resonant odd $n$ and every $j\ge0$, without a restriction
involving $\alpha_n$,
\begin{equation}\label{eq:jets_large_n_bound}
 \log\sum_w|a_{j,n,w}^{\mathrm{pole}}|
 \le 2j\log n-\frac12n^2\log n+Cn^2+Cj+C.
\end{equation}
Empty sectors are understood to contribute zero.
\end{lemma}

\begin{proof}
The residue denominator has modulus at least $cn$. Combine this
with \eqref{eq:jets_pole_lower} and
\eqref{eq:jets_numerator_vandermonde}, replace $S$ in its exponential
by $\sigma$, and extend the positive radial majorant to the whole
real zero-sum plane. Equation \eqref{eq:jets_beta_integral} proves
\eqref{eq:jets_winding_bound}. Its exact dependence on $w$ relative
to the same expression with $w=1$ is
$|w|^{-2(j+1-\alpha_n)}$.

For the second estimate, use the simpler bound $|T_n|\ge c/n^2$
and the exact ball mass
\[
 \int_{Q<\sigma}\Delta(\eta)^2\dd\eta
 =I_n\frac{\sigma^{\alpha_n}}{\Gamma(\alpha_n+1)}.
\]
The logarithm of this ball mass is
$-\tfrac12n^2\log n+O(n^2)$; this follows by applying Stirling's
inequalities to \eqref{eq:jets_Mehta} and to
$\Gamma(\alpha_n+1)$. Multiplication by $H_{0,n}$ costs $O(n^2)$.
The number of windings and the remaining polynomial factors are
absorbed into the displayed error. This proves
\eqref{eq:jets_large_n_bound} for arbitrarily large $n$.
\end{proof}

\subsection{The exact first-winding model}

Let
\begin{equation}\label{eq:jets_model_parameters}
 \tau_n=1-\cos(\pi/n),\quad B_n=B(\tau_n),\quad
 h_n=\cos(\pi/n)-\frac{B_n b}{s^2}\sin^2(\pi/n),\quad
 P_n=\prod_{r=1}^{n-1}r!\,2^{-r}.
\end{equation}
In particular $h_n=1+O(n^{-2})$. Proposition~\ref{prop:negative-real-cycle-coefficients} in
Section~\ref{sec:realcycle} gives, uniformly for odd multiples of $p$
in the window
\[
 \log j<n<\frac{\sqrt j}{(\log j)^{1/4}},
\]
\begin{equation}\label{eq:jets_uniform_model}
 a_{j,n,1}^{\mathrm{pole}}+a_{j,n,-1}^{\mathrm{pole}}
 =M(n,j)(1+e_{n,j}),
\end{equation}
where
\begin{equation}\label{eq:jets_model}
 M(n,j)=-\frac{2^{2-\alpha_n}B_n^{\alpha_n-1/2}P_n
                     n^{\alpha_n-3/2}}
                   {\pi^{(n-1)/2}s^{2\alpha_n}}
 h_n^{-\alpha_n}\tau_n^{\alpha_n-j-1}
 \frac{\Gamma(j+1-\alpha_n)}{\Gamma(j+1)},
\end{equation}
and
\begin{equation}\label{eq:jets_model_errors}
 |e_{n,j}|\le C\frac n{\sqrt j}.
\end{equation}
This estimate is for the actual inner ball of squared radius
$\sigma-\pi^2/n$ on the original real angular cycle. It introduces
no additional boundary remainder.

The constants in \eqref{eq:jets_model} follow from the central
numerator, pole derivative and constrained beta period.
At the collision $\gamma_i=\ii\phi$, the numerator divided by
$\Delta^2$ is
\[
 -\frac{\ii B_n^{n^2/2}}
       {n!(2\pi)^{n-1}2^{n(n-1)}s^{n^2}},
 \qquad \Gamma_t=-\frac{\ii nB_n}s.
\]
Multiplication by $-2/\Gamma_t$, followed by the exact complex quadratic beta integral, gives half of \eqref{eq:jets_model}.
Reflection $\delta\mapsto-\delta$ preserves the full integrand and
exchanges the first two winding sectors, giving the displayed
factor of two. The factor $B_n^{\alpha_n-1/2}$ retains its quadratic
particle-dependent phase. No positivity across particle indices is
being assumed.

\subsection{A discrete concavity estimate with controlled errors}

Write $L(n,j)=\log|M(n,j)|$ and let $d=2p$ be the step between
successive allowed odd multiples of $p$.  We need only coarse
asymptotics for the position of the largest term.  In particular no
constant term in the asymptotic expansion of that position is needed.
The exact model gives
\begin{equation}\label{eq:jets_L_expansion}
 L(n,j)=2j\log n-\frac{n^2}{2}\log j-j\log(\pi^2/2)
       +O\left(n^2+\frac{n^4}{j}+\frac{j}{n^2}\right),
\end{equation}
when $n\to\infty$ and $n^2/j\to0$.  The finite differences needed
below are obtained from the exact factors, not by differentiating this
error term.

\begin{lemma}[Exact-factor differences]\label{lem:jets_discrete_concavity}
For integer $j>\alpha_{n+d}$, uniformly as $n\to\infty$ and
$n^2/j\to0$,
\begin{align}
 L(n+d,j)-2L(n,j)+L(n-d,j)
 &=-d^2\left(\frac{2j}{n^2}+\log j\right)
       +O_d\left(1+\frac{n^2}{j}+\frac{j}{n^4}\right),
       \label{eq:jets_second_difference}\\
 \frac{L(n+d,j)-L(n-d,j)}{2d}
 &=\frac{2j}{n}-n\log j
       +O_d\left(n+\frac{n^3}{j}+\frac{j}{n^3}\right).
       \label{eq:jets_crossing_expansion}
\end{align}
In particular the second difference is at most $-cd^2\log j$
throughout
\begin{equation}\label{eq:jets_middle_window}
 \log j<n<\frac{\sqrt j}{(\log j)^{1/4}},
\end{equation}
whenever both neighboring indices are in this interval, for all
sufficiently large $j$.
\end{lemma}

\begin{proof}
Put $A=\pi^2/2$ and
$e_n=\log(\tau_n n^2/A)$.  The elementary expansions of $\tau_n$,
$B(\tau_n)$ and $h_n$ are convergent in $n^{-2}$ near zero.
Consequently the exact logarithm of \eqref{eq:jets_model} can be
organized as
\begin{equation}\label{eq:jets_exact_log_decomposition}
 \begin{split}
 L(n,j)={}&\log P_n+(2j-n^2/2+1)\log n
       -\alpha_n\log j-j\log A\\
       &+R_n-je_n+E_j(\alpha_n),\\
 E_j(a)={}&\sum_{r=0}^{a-1}-\log(1-r/j).
 \end{split}
\end{equation}
Here $e_n=O(n^{-2})$ has such a convergent expansion, while
$R_n$ is a quadratic polynomial in $n$ plus a convergent series
$O(n^{-2})$.  Its coefficients contain the fixed quantities
$B(0)$ and $s$, but their values are not needed.  This description
follows by grouping all powers of $n$ and $\tau_n=A n^{-2}e^{e_n}$;
the remaining factors $B(\tau_n)$ and $h_n$ are nonzero near
$n^{-1}=0$.  Thus $R_n$ has centered first difference $O_d(n)$
and second difference $O_d(1)$, while the corresponding differences
of $je_n$ are $O_d(j/n^3)$ and $O_d(j/n^4)$.

The factorial products have the exact quotients
\[
 \log P_{n+d}-\log P_{n-d}
   =\sum_{r=n-d}^{n+d-1}\bigl(\log(r!)-r\log2\bigr),
\]
\[
 \log P_{n+d}-2\log P_n+\log P_{n-d}
   =\sum_{k=0}^{d-1}\sum_{\ell=1}^{d}
        \log(n-d+k+\ell)-d^2\log2.
\]
Stirling's inequality applied to the finitely many factorials in the
first expression gives $2dn\log n+O_d(n)$; the second expression
is $d^2\log n+O_d(1)$.  Taking differences of the explicit elementary
function $(2j-n^2/2+1)\log n$ cancels these $\log n$ terms and leaves
respectively $2j/n+O_d(n+j/n^3)$ after division by $2d$, and
$-2d^2j/n^2+O_d(1+j/n^4)$.

It remains to control $E_j$.  On $0\le r\le\alpha_{n+d}$,
\[
 0\le-\log(1-r/j)\le Cn^2/j,\qquad
 \frac{\dd}{\dd r}[-\log(1-r/j)]\le C/j.
\]
The difference $E_j(\alpha_{n+d})-E_j(\alpha_{n-d})$ has
$2dn$ summands, so it is $O_d(n^3/j)$.  For the second difference,
the two increments of $\alpha_n$ have lengths
$dn+d^2/2$ and $dn-d^2/2$.  The $d^2$ excess summands cost
$O_d(n^2/j)$.  Pair each remaining upper summand with its lower
counterpart: their arguments differ by $O_d(n)$, so the derivative
bound gives another total $O_d(n^2/j)$.  Finally the differences of
$-\alpha_n\log j$ are exactly $-n\log j$ after division by $2d$,
and $-d^2\log j$.  This proves both stated formulas.

For completeness, ordinary Stirling bounds give
$\log P_n=\tfrac12n^2\log n+O(n^2)$, and
$E_j(\alpha_n)=O(n^4/j)$.  Substitution into
\eqref{eq:jets_exact_log_decomposition} proves
\eqref{eq:jets_L_expansion}.  In \eqref{eq:jets_middle_window},
$n^2/j=o(1)$ uniformly, and $j/n^4$ is smaller than $j/n^2$ by
at least the factor $(\log j)^{-2}$.  The negative second-difference
bound therefore holds throughout the stated interval.
\end{proof}

\subsection{Selecting an integer order with one dominating index}

Fix a large allowed particle index $n$, and for integer
$j>\alpha_{n+d}$ put
$F_n(j)=L(n+d,j)-L(n-d,j)$.  The gamma factors in the exact model
satisfy a one-step recurrence, giving the exact identity
\begin{equation}\label{eq:jets_crossing_derivative}
 F_n(j+1)-F_n(j)
 =\log\frac{\tau_{n-d}}{\tau_{n+d}}
   +\log\frac{j+1-\alpha_{n+d}}{j+1-\alpha_{n-d}}.
\end{equation}
Uniformly for $\tfrac12n^2\log n\le j\le\tfrac32n^2\log n$,
this increment is $4d/n+O_d(1/(n\log n))$, hence positive and
$O_d(1/n)$.

Equation~\eqref{eq:jets_crossing_expansion} proves an actual sign
change on this integer interval.  Indeed, at
$j=c n^2\log n+O(1)$ with $c=1/2$ or $3/2$, its leading expression is
\[
 \frac{2j}{n}-n\log j
   =2(c-1)n\log n-n\log\log n-n\log c+O(1/n),
\]
whereas the error in \eqref{eq:jets_crossing_expansion} is $O_d(n)$.
Thus the two endpoints have opposite signs for large $n$.
Choose $j_n$ as the first integer in this interval with $F_n(j_n)\ge0$.
The recurrence then gives $0\le F_n(j_n)=O_d(1/n)$.  Substituting
this in \eqref{eq:jets_crossing_expansion}, and using
$\log j_n=2\log n+O(\log\log n)$, yields
\begin{equation}\label{eq:jets_selected_j}
 j_n\sim n^2\log n,\qquad F_n(j_n)=O_d(1/n).
\end{equation}
The selected index $n$ and its neighbors are strictly inside
\eqref{eq:jets_middle_window} for these $j_n$.

Let $\ell_k=L(n+kd,j_n)$ for all allowed indices in that window.
Lemma~\ref{lem:jets_discrete_concavity} says that successive differences
of $\ell_k$ decrease by at least $cd^2\log j_n$.
Since $\ell_1-\ell_{-1}=F_n(j_n)=O_d(1/n)$, the two adjacent gaps
satisfy $\ell_0-\ell_{\pm1}\ge c'd^2\log j_n$ for large $n$.
Summing the decreasing successive differences in both directions gives
\[
 \ell_k-\ell_0\le-c''d^2 k^2\log j_n+O_d(|k|/n),\qquad k\ne0.
\]
Therefore
\begin{equation}\label{eq:jets_single_index}
 \sum_{\substack{m\ne n,\ p\mid m,\ m\ \mathrm{odd}\\
                  m\text{ in }\eqref{eq:jets_middle_window}}}
 |M(m,j_n)|
 \le |M(n,j_n)|\,O(j_n^{-c'''d^2}).
\end{equation}
No common phase of these terms is assumed.  A uniform $o(1)$ relative
error in the local model remains sufficient, because it multiplies
the already summable moduli on the left.

\subsection{Completing the all-index comparison}

For $j\to\infty$, split the odd particle indices at
\[
 n_{\mathrm{lo}}=\log j,\qquad
 n_{\mathrm{hi}}=\frac{\sqrt j}{(\log j)^{1/4}}.
\]
The low range, including every nonasymptotic index, is bounded by
Lemma~\ref{lem:jets_boundary_derivatives}:
\begin{equation}\label{eq:jets_low_range}
 \sum_{n\le n_{\mathrm{lo}}}|[t^j]C_n|
 \le\exp\{2j\log\log j+O(j)\}.
\end{equation}
All $c/n$ analytic remainders have the smaller scale
$\exp\{\tfrac12j\log j+O(j)\}$ by
\eqref{eq:jets_remainder_bound}.

For $n\ge n_{\mathrm{hi}}$, use the ball-mass bound
\eqref{eq:jets_large_n_bound}. In the variable $x=n^2$ its
controlling exponent is
\[
 j\log x-\tfrac14x\log x+Cx+Cj,
\]
which is decreasing for every $x\ge j/\sqrt{\log j}$ when $j$
is sufficiently large. At the threshold it equals at most
\begin{equation}\label{eq:jets_high_range}
 j\log j-\tfrac12j\log\log j-\tfrac14j\sqrt{\log j}
       +o(j\sqrt{\log j}).
\end{equation}
Its subsequent decrease makes the full infinite tail summable,
with the same leading bound. This estimate handles all windings in
the high range and requires no beta integral outside its
convergence range.

In the intermediate range, the error in
\eqref{eq:jets_uniform_model} is uniformly
$O((\log j)^{-1/4})=o(1)$. It remains to compare the higher odd
windings to that model. Let $M_0(n,j)$ be the modulus of
\eqref{eq:jets_model} with $B_n$ replaced by $B(0)$, $h_n$ by
$1$, and $\tau_n$ by $\pi^2/(2n^2)$. The exact formulas give
\begin{equation}\label{eq:jets_baseline_ratio}
 \left|\log\frac{M_0(n,j)}{|M(n,j)|}\right|
 \le C(j/n^2+1)=o(j)
\end{equation}
uniformly in this entire range. Equation
\eqref{eq:jets_winding_bound} is bounded by $CM_0(n,j)$ times
\[
 e^{Cn\sigma}(1-C\sigma)^{-j-1}
 |w|^{-2(j+1-\alpha_n)}.
\]
Here $n=o(j)$ and $\alpha_n/j=o(1)$ uniformly. Choose the fixed
$\sigma$ so small that $-\log(1-C\sigma)<\tfrac12\log3$.
Summing over all odd $|w|\ge3$ then gives
\begin{equation}\label{eq:jets_higher_windings_negligible}
 \sum_{|w|\ge3}|a_{j,n,w}^{\mathrm{pole}}|
 \le |M(n,j)|e^{-c j}
\end{equation}
for some $c>0$ and all sufficiently large $j$, uniformly over the
intermediate range. Thus no unexamined winding or particle range
remains between the preceding estimates.

At the selected sequence \eqref{eq:jets_selected_j}, the coarse
expansion \eqref{eq:jets_L_expansion} suffices for the exact leading
growth constant.  Indeed $j_n\sim n^2\log n$ implies
\[
 2\log n=\log j_n-\log\log j_n+\log2+o(1),\qquad
 \frac{n^2\log j_n}{2j_n}\longrightarrow1.
\]
Every error in \eqref{eq:jets_L_expansion} is $o(j_n)$ at this scale.
Consequently
\begin{align}
 \log|M(n,j_n)|
 &=j_n\log j_n-j_n\log\log j_n
   +j_n\bigl(\log(4/\pi^2)-1\bigr)+o(j_n),
 \label{eq:jets_main_growth}\\
 |M(n,j_n)|^{1/j_n}
 &\sim\frac{4j_n}{e\pi^2\log j_n}.
 \notag
\end{align}
The exact gamma ratio was retained in the finite-difference and
relative local estimates.  Here its logarithmic correction to
$-\alpha_n\log j_n$ is $O(n^4/j_n)=o(j_n)$, as included in
\eqref{eq:jets_L_expansion}.


The bounds \eqref{eq:jets_low_range} and
\eqref{eq:jets_high_range}, and all analytic remainders, are
$o(|M(n,j_n)|)$ by \eqref{eq:jets_main_growth}. The intermediate
higher windings have the same property by
\eqref{eq:jets_higher_windings_negligible} and
\eqref{eq:jets_single_index}. Finally the uniform relative error
in \eqref{eq:jets_uniform_model}, summed over that range, is at
most $o(1)\sum_m|M(m,j_n)|=o(|M(n,j_n)|)$. Therefore the
complete boundary coefficient satisfies
\begin{equation}\label{eq:jets_complete_noncancellation}
 a_{j_n}=M(n,j_n)(1+o(1)).
\end{equation}
This proves noncancellation in the infinite odd-particle sum even
though the model phases vary with $n$.

\begin{theorem}[A regular rational negative edge is singular for the full odd sum]
\label{thm:jets_negative_edge}
Under the physical-branch and nonramified-coordinate hypotheses
of \eqref{eq:jets_setting}, the complete odd sum has all inward
boundary derivatives at $t=0$, but the Taylor series of those
boundary derivatives has radius of convergence zero. More
precisely, along the integer sequence \eqref{eq:jets_selected_j},
\[
 |a_{j_n}|^{1/j_n}\sim\frac{4j_n}{e\pi^2\log j_n}\longrightarrow\infty.
\]
Consequently the complete high-temperature physical germ has no
holomorphic continuation across this regular negative edge
point. Multiplication by its nonzero analytic physical
normalization preserves this conclusion.
\end{theorem}

\begin{proof}
The jets exist and equal the absolutely convergent coefficient
sums by Lemma~\ref{lem:jets_boundary_derivatives}. Equations
\eqref{eq:jets_main_growth} and
\eqref{eq:jets_complete_noncancellation} give their stated
growth and zero Taylor radius. If there were a holomorphic
continuation across the point, its Taylor coefficients would be
these same boundary jets and would satisfy a positive-radius
Cauchy bound, a contradiction. The physical normalization can
be divided out in a neighborhood because it is analytic and
nonzero there. The use of the nonramified coordinate $t$ makes
this a statement about the fixed-ratio one-variable physical
germ itself.
\end{proof}




\section{Closure, exceptional points, and coefficient interfaces}\label{sec:closure}

\subsection{Completing the unequal-coupling theorem}
We now trace the two physical germs to their full boundaries. The
representation in Section~\ref{sec:physical} gives holomorphic functions
on the explicit, independently determined domains $\OH$ and $\OL$.
Their equality with the two physical susceptibilities holds on the
appropriate positive real intervals. The identity theorem fixes the
continued physical germs on the entire corresponding connected domains.

For each fixed $\lambda\ne1$, the profile-intersection set and the
two lens tips in Section~\ref{sec:geometry} form a locally discrete
exceptional set on the physical boundaries. Every other boundary
point has a real analytic edge arc with exactly one endpoint cosine,
a strictly interior second cosine, a nonzero real-linear determinant,
and a nonramified holomorphic edge coordinate. Restricting that
coordinate to the arc gives an open real interval. Cosines of primitive
odd-prime rational angles are dense in every such interval.

On every regular LT arc, Section~\ref{sec:regular} proves a genuine
singularity of the entire even-particle sum at each point of this dense
family, including both endpoint signs and both transfer-axis choices.
The first singular derivative has a nonzero coefficient and every
higher even index is included in the absolute bounded tail. On every
regular positive HT arc, the odd counterpart proves the analogous
whole-sum singularity. On every regular negative HT arc,
Sections~\ref{sec:realcycle}--\ref{sec:jets} show that all fixed boundary
jets exist, but a sequence of their Taylor coefficients satisfies
\begin{equation}\label{eq:boundarygrowthclosure}
 |a_{j_n}|^{1/j_n}\sim\frac{4j_n}{e\pi^2\log j_n}\longrightarrow\infty.
\end{equation}
These coefficients are those of the complete odd sum. The Gaussian
bound permits taking the boundary limit through the sum for each fixed
derivative. If a holomorphic extension existed, its derivatives would
have these limits and its Taylor series would have positive radius,
contradicting \eqref{eq:boundarygrowthclosure}. Division by the local
analytic nonzero physical prefactor proves the same assertion for
$\XH$, and multiplication by the nonzero $\beta$ gives $\chi_H$.

Finally suppose a holomorphic extension existed through any boundary
point, including an exceptional point. A small disk of extension contains
regular arcs of the same physical boundary, and hence contains a nearby
point of the appropriate dense true-singularity family. Restricting the
extension to that point contradicts its already proved nonremovability.
The complete graph $\partial\OL$ and the complete closed curve
$\partial\OH$ therefore are natural boundaries. This proves
the unequal-coupling part of Theorem~\ref{thm:aniso}.

For unequal couplings $\cS_\lambda$ has interior. The natural-boundary
conclusion concerns the explicitly determined boundaries of its two
physical complementary components. Their nonextension follows from
the complete-sum estimates and the density argument above.

\subsection{Every fixed ratio and exceptional contact}
For $\lambda\ne1$, the signed-profile difference in
\eqref{eq:signed-profiles} is a nonzero real-analytic function on
the entire real axis. Its isolated zeros include all contacts of the
two endpoint graphs, independently of whether the ratio is rational
or irrational. Away from those contacts and the two lens tips, the
endpoint-defect identity \eqref{eq:endpoint-defect} proves strictness
of the interior cosine and independence of the two spectral directions;
the explicit endpoint-root derivative proves nonramification.
The closure argument therefore includes every finite exceptional point.
The odd-over-odd rational contacts at the imaginary axis are covered by
the same analytic-profile argument.

The negative-HT proof uses $B(t)$ only as a fixed holomorphic nonzero
function near the selected regular point. Its exact centers
$1-\cos(\pi/n)$, the real-cycle beta estimate, and single-index
comparison require no arithmetic control of the temperature coordinate
or of $\lambda$. All constants are chosen after fixing $\lambda$ and
the regular boundary point. The case $\lambda=1$ is handled next by
the internal proof in Section~\ref{sec:iso-internal}; no limiting
assertion in $\lambda$ is needed.

\subsection{Exact coefficient interface on the isotropic slice}
Set $\lambda=1$ and $s=\sinh z$. Then $S=T=s$, $p=s^2$, $B=1$, and
\[
 D(x,y;z)=s\left[s+s^{-1}
                  -\tfrac12(x+x^{-1})-\tfrac12(y+y^{-1})\right].
\]
On the positive LT interval $\sqrt p=s$, so every factor
$\sqrt p/D$ in \eqref{eq:doublecorr} is exactly the normalized
isotropic factor.
The even coefficients and the full-site identity agree term by term.
In particular the function to which the isotropic result applies is
\begin{equation}\label{eq:iso-LT-interface}
 \mathcal X_L(s)=M^2\sum_{N\ge2,\ N\ {\rm even}}C_N(s)
 =1-M^2+2M^2\sum_{N\ge2,\ N\ {\rm even}}T_N(s),
 \qquad M^2=(1-s^{-4})^{1/4}.
\end{equation}
The onsite term $1-M^2$ follows from the complete connected correlation
sum; no lattice site is omitted.

In HT, \eqref{eq:fullangular} has $W=s+s^{-1}-\cos q$ and
$B=1$. In the high-temperature interval $0<s<1$, the exterior
parameter is $t=1/s>1$ and the complete angular coefficient is
$C_N(s)$ as defined by \eqref{eq:fullangular}. Consequently
\begin{equation}\label{eq:iso-HT-interface}
 \mathcal X_H(s)=\frac{(1-s^4)^{1/4}}s
       \sum_{N\ge1,\ N\ {\rm odd}}C_N(s).
\end{equation}
The roots, independent-angle measures, all particle parities and physical
prefactors therefore agree on the respective real intervals. The physical
state and energy-field identifications are supplied by
Section~\ref{sec:physical}, including at $B=1$.

The coefficient identities above identify the full physical observables
defined in Section~\ref{sec:physical}. Section~\ref{sec:iso-internal}
proves their equal-coupling boundary obstruction by a balanced
deformation, regular-coordinate current and extreme-coordinate coarea
for both parities. Finite-cell contraction and retained collision
factors control the complete differentiated tail.

\section{The isotropic slice from its full physical product}
\label{sec:iso-internal}

Here $\lambda=1$, $s=\sinh z$, and the dispersion parameter will be
written $\mathsf S$ to distinguish it from the notation of
Section~\ref{sec:physical}.  The full product formula already proved
there specializes to $B=1$ and $C=s+s^{-1}$.  We give the remaining
analytic argument directly for that product.

For $\mathsf S\in\C\setminus[-2,2]$ and $y=e^{\ii\theta}$, let
\[
 W=\mathsf S-\cos\theta,\qquad
 \rho+\rho^{-1}=2W,\quad |\rho|<1,\qquad
 \sinh\gamma=(\rho^{-1}-\rho)/2.
\]
The root is unique and holomorphic: a root on the unit circle would
force $W\in[-1,1]$, and the two roots are reciprocal.  Write
\[
 P_{ij}=-\frac{(y_i-y_j)^2\rho_i\rho_j}
                  {y_i y_j(1-\rho_i\rho_j)^2},\qquad
 Z=\prod_i\rho_i,
 \qquad a_+=0,\quad a_-=\pi.
\]
For either sign $\sigma\in\{1,-1\}$ define the exact auxiliary coefficient
\begin{equation}\label{eq:iso-product}
 \mathcal F_N^\sigma(\mathsf S)
 =\frac1{N!}\int_{\mathbb H_N^\sigma}
       \frac{1+\sigma Z}{1-\sigma Z}
       \prod_i\frac1{\sinh\gamma_i}\prod_{i<j}P_{ij}\,d\mu_N,
 \qquad
 \mathbb H_N^\sigma=\left\{\sum_i\theta_i=a_\sigma
                                    \pmod{2\pi}\right\},
\end{equation}
where
\[
 d\mu_N=\prod_{i<N}\frac{d\theta_i}{2\pi},\qquad
 \theta_N=a_\sigma-\sum_{i<N}\theta_i\pmod{2\pi}.
\]
Thus this is Euclidean hypersurface measure divided by
$\sqrt N(2\pi)^{N-1}$.  For $N=1$ the constrained integral is evaluation
at its single angle.  All pair factors are rational squares, so no
choice of half-angle or logarithm changes the coefficient.  Put
\begin{equation}\label{eq:iso-exterior-coefficients}
 C_N^\sigma(t)=\mathcal F_N^\sigma(t+t^{-1}),\qquad |t|>1.
\end{equation}
The compact normal-convergence bounds for the full product apply to
each fixed differentiated series here.  To see directly that the
auxiliary sign causes no change, on a compact dispersion neighborhood
all roots satisfy $|\rho_i|\le q<1$ and have imaginary parts of one
sign.  The latter fact follows from
$\Im W_i=\Im\rho_i(1-|\rho_i|^{-2})/2$.
The identity
\[
 |1-vw|^2-|v-w|^2
   =(1-|v|^2)(1-|w|^2)+4\Im v\Im w
\]
therefore bounds every root Schur factor
$|(\rho_i-\rho_j)/(1-\rho_i\rho_j)|$ strictly below one, uniformly.
Expand the angular Schur product by the matching identity proved
in Section~\ref{sec:physical} and cancel each matched angular pole
with its corresponding root factor.  The resulting complete matched
pairs are bounded; all but $O(N)$ root factors retain the strict
contraction.  The number of matchings costs
$\exp(O(N\log(N+1)))$, while the unmatched root factors supply
$\exp(-cN^2)$.  This gives a summable coefficient bound on the
normalized constrained torus for either sign.  All one-body factors
and $1-\sigma Z$ have fixed gaps there.  Cauchy's inequality on a
slightly larger compact neighborhood gives every fixed derivative.

\subsection{Exact physical and sign identities}

The physical real-temperature identities, with their analytic
continuations fixed by those real values, are exactly
\begin{align}
 \mathcal X_L(s)
   &=(1-s^{-4})^{1/4}
       \sum_{\substack{N\ge2\\N\ {\rm even}}}C_N^+(s),
       &&|s|>1,\label{eq:iso-physical-L}\\
 \mathcal X_H(1/t)
   &=Q_H(t)\sum_{\substack{N\ge1\\N\ {\rm odd}}}C_N^+(t),
       &Q_H(t)&=t(1-t^{-4})^{1/4},\quad |t|>1.
       \label{eq:iso-physical-H}
\end{align}
The fourth roots tend to one at infinity.  These identities follow by
putting $p=s^2$ and $B=1$ in the main physical formulas.  In the LT
identity the zero-particle term has already been subtracted in the plus
state.  The lattice origin is included: its contributions are
$1-(1-s^{-4})^{1/4}$ in LT and $1$ in HT.  In particular there is no
additional constant to append to either full coefficient sum.

Angle reflection preserves $\mathbb H_N^\sigma$, its Haar measure and
each complete factor in \eqref{eq:iso-product}.  Conjugation also
commutes with the coefficients.  Replacing $\mathsf S,y_i,\rho_i$ by
$-\mathsf S,-y_i,-\rho_i$ gives the precise negation identity
\begin{equation}\label{eq:iso-negation}
 \mathcal F_N^\sigma(-\mathsf S)
   =(-1)^N\mathcal F_N^{(-1)^N\sigma}(\mathsf S).
\end{equation}
Every pair is unchanged, every one-body factor changes sign, and
$\prod_i y_i$ and $Z$ both acquire $(-1)^N$.  This proves
\eqref{eq:iso-negation} with its constraint and product-kernel signs.
In particular the even physical coefficients are even in $t$, whereas
the odd coefficients satisfy
\begin{equation}\label{eq:iso-odd-negation}
 C_N^+(-t)=-C_N^-(t)\qquad(N\ {\rm odd}).
\end{equation}
No evenness assertion about the HT observable is made.
The two one-particle formulas, useful below, are
\begin{equation}\label{eq:iso-one-particle}
 \mathcal F_1^+(\mathsf S)=\frac1{\mathsf S-2},\qquad
 \mathcal F_1^-(\mathsf S)=\frac1{\mathsf S+2}.
\end{equation}
For example the second formula follows by taking $\theta=\pi$ and
using $(1-\rho)/(1+\rho)\sinh(\gamma)^{-1}
=2\rho/(1+\rho)^2$.

\subsection{Prime arithmetic and both HT geometries}

\begin{lemma}[The first order and the branch-mean gap]
\label{lem:iso-prime-family}
Let $p\ge5$ be prime, $1\le a\le(p-1)/2$, and
$\alpha=2\pi a/p$, $\zeta=e^{\ii\alpha}$,
$\mathsf S_*=2\cos\alpha$.
The first odd positive integer $n$ admitting
\begin{equation}\label{eq:iso-Nickel}
 \mathsf S_* = \cos(2\pi l/n)+\cos(2\pi m/n)
\end{equation}
is $n=p$.  At $p$ its unordered cosine pair is uniquely
$\{\cos\alpha,\cos\alpha\}$.
If $0<\alpha<\pi/2$, the corresponding first even order is $2p$,
and the same uniqueness holds there.
For $b=\arccos(|\mathsf S_*|-1)$, $e^{-\ii b}$ is algebraic and
is not a root of unity.  The selected points are dense in both open
upper quadrants for odd orders, and in the first quadrant for even
orders.
\end{lemma}
\begin{proof}
Write $d=(p-1)/2$, $\zeta_p=e^{2\pi\ii/p}$ and
$u_j=\zeta_p^j+\zeta_p^{-j}$ for $1\le j\le d$.
These $d$ numbers are linearly independent over $\mathbb Q$.
Indeed a relation gives a polynomial of degree at most $p-1$, with
zero constant term, vanishing at $\zeta_p$; its divisibility by
$1+X+\cdots+X^{p-1}$ forces every coefficient to vanish.
Each $u_a$ is irrational, since otherwise $\zeta_p^a$ would have
degree at most two over $\mathbb Q$.

If $p\nmid n$, the two cyclotomic fields
$\mathbb Q(\zeta_p)$ and $\mathbb Q(\zeta_n)$ have intersection
$\mathbb Q$: their compositum has degree
$\varphi(pn)=\varphi(p)\varphi(n)$.
Thus \eqref{eq:iso-Nickel} is impossible at every odd $n<p$ and
every even $n<2p$.  The equal pair realizes orders $p$ and $2p$.

Put $\delta=\operatorname{sgn}(u_a)$ and $\mathsf S_+=\delta u_a$.
The conjugates of $u_a$ are all the $u_j$, including a positive and
a negative number, so $\mathsf S_+$ has a negative conjugate.
The number $\xi=e^{-\ii b}$ satisfies
$\xi+\xi^{-1}=2\mathsf S_+-2$ and is algebraic.
If it were a root of unity, every conjugate of the right side would
lie in $[-2,2]$.  A negative conjugate of $\mathsf S_+$ makes it
strictly less than $-2$, a contradiction.  Embeddings extend to a
field containing $\xi$, so this conjugate argument is legitimate.

At odd order $p$, the cosine $-1$ is absent.  A cosine $1$ has an
impossible partner when $\mathsf S_*<0$, and when $\mathsf S_*>0$
it would make $\xi$ a root of unity.  All remaining twice-cosines
are $u_j$; independence applied to $2u_a=u_l+u_m$ gives uniqueness.
At even order $2p$ and $\mathsf S_*>0$, a cosine $-1$ is impossible
and $1$ is excluded as before.  The other twice-cosines are
$\pm u_j$, and the same independence again forces both to equal
$u_a$.  Finally choose arbitrarily large primes and a nearest
integer to $p\alpha_0/(2\pi)$ for any prescribed angle in the
stated open interval.  This proves density.
\end{proof}

All estimates below are performed at a positive dispersion point.
For LT take $0<\alpha<\pi/2$, $\sigma=+$ and $n_*=2p$.
For an HT prime point in the first quadrant take
$\sigma=+$ and $n_*=p$.  At an HT prime point in the second quadrant,
put
\begin{equation}\label{eq:iso-negative-geometry}
 \widehat\alpha=\pi-\alpha,\qquad
 \widehat\zeta=-\overline\zeta=e^{\ii\widehat\alpha},\qquad
 \widehat{\mathsf S}_*=-\mathsf S_*>0,
 \qquad \sigma=-,\quad n_*=p.
\end{equation}
Then $p\widehat\alpha$ is an odd multiple of $\pi$.
The signed equality condition has both angles in
$(2\pi\Z+\pi)/n$.  Adding $\pi$ to both angles, for odd $n$,
converts it to the ordinary condition at the negative dispersion
point.  Lemma~\ref{lem:iso-prime-family} therefore proves absence of
lower signed orders and uniqueness at $p$ as well.  The two equality
centers in positive geometry are $\theta_i=\pm\widehat\alpha$.
The exact derivative transfer back to the physical coefficient is
\begin{equation}\label{eq:iso-negative-jets}
 (C_N^+)^{(j)}((1+\eps)\zeta)
  =(-1)^{j+1}
      \overline{(C_N^-)^{(j)}((1+\eps)\widehat\zeta)}.
\end{equation}
It follows from \eqref{eq:iso-odd-negation} and conjugation; no
identification of different physical phases is used.

In the positive geometry, write the angle as $\alpha_+\in(0,\pi/2)$
and keep its corresponding sign $\sigma$.
For every finite set of particle orders the constrained equality means
\[
 \theta_0=(2\pi l+a_\sigma)/n
\]
stay a positive distance from the lower branch $-b$.
Otherwise $e^{-\ii nb}=\sigma$, contrary to
Lemma~\ref{lem:iso-prime-family}.  This is a finite-order assertion;
no uniform arithmetic gap for unbounded $n$ is asserted or used.
The radial path and derivative budget are
\begin{equation}\label{eq:iso-budget}
 t_\eps=(1+\eps)e^{\ii\alpha_+},\qquad
 \mathsf S_\eps=t_\eps+t_\eps^{-1}
 =2\cos\alpha_++2\ii\sin\alpha_+\eps+O(\eps^2),
 \qquad
 k=\left\lfloor\frac{n_*^2-1}{2}\right\rfloor.
\end{equation}
In particular $\Re(\mathsf S_\eps-\mathsf S_*)=O(\eps^2)$.

\subsection{One analytic tail estimate for both isotropic parities}
\label{sec:iso-shared-tail}

We use the signed coefficients $\mathcal F_N^\sigma(\mathsf S)$,
$\sigma\in\{1,-1\}$, on
$\mathbb H_N^\sigma=\{\sum\theta_i=a_\sigma\pmod {2\pi}\}$,
where $a_+=0$ and $a_-=\pi$.  Thus, with the independent-coordinate
Haar measure $\prod_{i<N}\dd\theta_i/(2\pi)$,
\begin{equation}\label{eq:iso-tail-coefficient}
 \mathcal F_N^\sigma(\mathsf S)=\frac1{N!}
 \int_{\mathbb H_N^\sigma}
 \frac{1+\sigma Z}{1-\sigma Z}
 \prod_i\frac1{\sinh\gamma_i}\prod_{i<l}P_{il}\,
 \prod_{i<N}\frac{\dd\theta_i}{2\pi},\qquad
 P_{il}=\frac{\sin^2((\theta_i-\theta_l)/2)}
                {\sinh^2((\gamma_i+\gamma_l)/2)}.
\end{equation}
Here $\cosh\gamma_i=\mathsf S-\cos\theta_i$, $z_i=e^{-\gamma_i}$
is the interior root, and $Z=\prod_i z_i$.
Fix $\mathsf S_*=2\cos\alpha\in(0,2)$, and approach it with
\begin{equation}\label{eq:iso-parameter-approach}
 \Im\mathsf S_\eps\asymp\eps>0,
 \qquad \Re\mathsf S_\eps-\mathsf S_*=O(\eps^2).
\end{equation}
All constants below may depend on this point and on fixed derivative
orders, but never on $N$ or $\eps$.  The radial exterior approach in
the physical coordinate satisfies \eqref{eq:iso-parameter-approach}.

\begin{theorem}[Complete isotropic higher tail]\label{thm:iso-tail}
Let $n_*\ge2$, and put $k=(n_*^2-1)/2$ when $n_*$ is odd and
$k=n_*^2/2-1$ when $n_*$ is even.  For either sign $\sigma$,
\begin{equation}\label{eq:iso-complete-tail}
 \sum_{\substack{N\ge n_*+2\\N\equiv n_*\ (2)}}
  |\partial_{\mathsf S}^j\mathcal F_N^\sigma(\mathsf S_\eps)|
 \le C_j(1+|\log\eps|),\qquad 0\le j\le k.
\end{equation}
The estimate includes every higher particle order.  Under any fixed
holomorphic change of the exterior parameter, it also holds for all
parameter derivatives through order $k$.
\end{theorem}

\subsubsection{A deformation of the complete constrained torus}

Set $b=\arccos(\mathsf S_*-1)$; then $\alpha<b<2\alpha$.
Choose a smooth periodic $p$, supported in $(0,\pi)$, with
$0\le p\le1$ and $p=1$ on $[\delta,\pi-\delta]$; put
$m(\theta)=p(-\theta)$.  The endpoint width $\delta$ is small,
so both $\pm b$ lie in the corresponding plateaus.  With a small
fixed $\tau>0$, define
\begin{equation}\label{eq:iso-balanced-deformation}
 P=\sum_i p(\theta_i),\quad Q=\sum_i m(\theta_i),\quad
 h_i=\frac\tau N(Pm(\theta_i)-Qp(\theta_i)),\quad
 y_i(\xi)=e^{i\theta_i+\xi h_i},\quad 0\le\xi\le1.
\end{equation}
The sum of the $h_i$ is zero, so the constrained product of the $y_i$
is unchanged.  For $y=e^{h+i\theta}$,
\[
 \Im\left(\mathsf S-\frac{y+y^{-1}}2\right)
       =\Im\mathsf S-\sinh h\sin\theta.
\]
The motion is inward on the upper semicircle and outward on the lower
one.  Its imaginary part therefore remains positive throughout the
homotopy.  All roots stay on the interior sheet, all one-body and pair
denominators stay nonzero, and $|Z|<1$.  The complete holomorphic
integrand can consequently be transported over this closed constrained
torus.  The weights below are introduced after the homotopy.

Write $w_i=\theta_i-ih_i$ at its endpoint, and $M=I-iDh$.
Because $\boldsymbol1^TM=\boldsymbol1^T$, the full determinant of $M$
equals its tangent determinant on the constrained hyperplane.  Moreover
\[
 Dh=\frac\tau N\operatorname{diag}(Pm'_i-Qp'_i)
       +\frac\tau N(m\otimes p'-p\otimes m').
\]
Each summand has operator norm $O(\tau)$, independently of $N$.
Thus $M^{-1}$ is uniformly bounded, $|\det M|\le C^N$, and its
fixed angular jets cost at most $C_j^NN^{C_j}$, by the determinant
and inverse-matrix differentiation rules.

Let $\eta=0$ on $(-\infty,1/4]$ and $\eta=1$ on $[1/2,\infty)$,
with $0\le\eta\le1$.  On the final contour use
\[
 \omega_F=\eta(P)\eta(Q),\quad
 \omega_L=\tfrac12(1-\eta(P))(1+\eta(Q)),\quad
 \omega_U=\tfrac12(1-\eta(Q))(1+\eta(P)).
\]
They sum to one.  Reflection of all angles exchanges $P,Q$, sends
$w_i$ to $-w_i$, and preserves every complete pair and one-body factor.
It also preserves either signed constraint.  Consequently
\begin{equation}\label{eq:iso-exact-split}
 \mathcal F_N^\sigma=F_N+2L_N.
\end{equation}
The support of $L_N$ has $P\le1/2$ and contains only the lower
semicircle and short upper endpoint slivers.  Lift the latter near
$\pi$ across $-\pi$; the rational pair factors and the constraint are
periodic.  Its angles then lie in $(-\pi-\delta,\delta)$.

\subsubsection{The mixed branch populations}

On the support of $F_N$, $P,Q\ge1/4$.  Since the bump supports stay
away from $0,\pi$, the imaginary-part identity gives
\begin{equation}\label{eq:iso-mixed-gap}
 -\log|Z|\ge c\left(N\eps+\frac{\tau PQ}{N}\right)\ge c/N.
\end{equation}
Every true lower branch has normal displacement $\tau P/N$, and every
true upper branch has displacement $\tau Q/N$.  The square-root form
of $\gamma$ at $W=1$ then gives a branch distance at least $c/N$ in
the $W$ plane and $|\partial_{\mathsf S}\log z_i|\le C\sqrt N$.
These statements persist on a common disk
$|\mathsf S-\mathsf S_\eps|\le c/N^3$: the variation of $\log Z$
is at most $CN^{3/2}c/N^3$, smaller than half the gap in
\eqref{eq:iso-mixed-gap}.  Constants can be decreased once to cover
all $N\ge2$.

We estimate the pair product by expanding the angular Schur factors into matchings.
Put $A(u,v)=(u-v)/(1-uv)$.  Subtracting the two dispersion equations
gives the exact identity $P_{il}=A(z_i,z_l)A(y_i,y_l)$.
The scalar Schur--Pfaffian identity expands the product of the
$A(y_i,y_l)$ as a sum over matchings when $N$ is even.  When $N$ is
odd, append the constant ghost entry obtained by setting its variable
to $-1$; its entries are the removable constant one.  In each term
combine a matched angular factor with its root factor to obtain the
complete $P_{il}$.  Every remaining root factor has modulus at most
one on the actual contour, since
\begin{equation}\label{eq:iso-halfdisk-identity}
 |1-uv|^2-|u-v|^2
   =(1-|u|^2)(1-|v|^2)+4\Im u\Im v\ge0
\end{equation}
in the closed lower half disk.  This is the same scalar matching
identity as in Lemma~\ref{lem:matching}, with no strict contraction
required here.

Let $B_0$ count coordinates in fixed neighborhoods of the two true
branches.  Each contributes one to $P+Q$, so $B_0\le P+Q$.
Same-branch complete pairs are bounded: the even inverse in the
branch coordinate cancels $\gamma_i+\gamma_l$, leaving
$\gamma_i-\gamma_l$.  All pairs away from opposite branches are also
bounded by compactness.  Opposite-branch coordinates lie in a common
sector of width less than $\pi$, and hence
\[
 |1-z_i z_l|\ge c\sqrt{(P+Q)/N},\qquad
 |P_{il}|\le CN/(P+Q)\le CN/B_0.
\]
At most $B_0/2$ matched pairs are exceptional.  Since
$B_0\log(N/B_0)\le N/e$, their total extra cost is $C^N$.
There are $(N-1)!!$ matchings in the even case and $N!!$ in the odd
case.  Thus both parities satisfy
\begin{equation}\label{eq:iso-mixed-pairs}
 \left|\prod_{i<l}P_{il}\right|\le C^N N^{(N+1)/2}.
\end{equation}

This bound holds on the protected complex disk as well.  Indeed
\[
 \partial_{\mathsf S}A(u,v)
 =\frac{u_{\mathsf S}(1-v^2)+v_{\mathsf S}(u^2-1)}{(1-uv)^2}.
\]
Near the branches, $|1-uv|\asymp|\gamma_i|+|\gamma_l|$,
$|u_{\mathsf S}|\le C/|\gamma_i|$, and
$|\gamma_i|,|\gamma_l|\ge cN^{-1/2}$.  The displayed derivative
is $O(N)$; outside these neighborhoods it is no larger.
Each unmatched root entry therefore grows by at most $O(N^{-2})$
on the $c/N^3$ disk.  Their $O(N^2)$-fold product costs only a
constant.  The canceled matched pairs retain their previous bounds
there by the same sector estimates.

The one-body factors have the independent majorant
\[
 f(\theta)=C\left(1+
       \operatorname{dist}(\theta,\{b,-b\})^{-1/2}\right).
\]
Near a branch $\theta_b\in\{b,-b\}$, use the imaginary part of
$W-1$ when $|\theta-\theta_b|\le C(\eps+|h|)$, and its linear
real part outside that region.  The small
complex disk is absorbed by the branch displacement.  In addition,
the uniform branch distance $c/N$ gives
$|\sinh\gamma_i|^{-1}\le C\sqrt N$ for every coordinate throughout
the mixed support and its protected disk; outside the true branch
neighborhoods the factor is bounded by a constant.  Eliminate the
fixed $N$th angle with the constraint, bound its one-body factor by
$C\sqrt N$, and use the independent $L^1$ majorants for the others.
This gives, for every $N\ge2$,
\[
 \int_{\mathbb H_N^\sigma}\omega_F
       \prod_i|\sinh\gamma_i|^{-1}\,d\mu_N
       \le C\sqrt N\,\|f\|_1^{N-1}\le C^N\sqrt N .
\]
The estimate remains valid when the harmless cutoff is bounded by one
and the remaining independent integrations are enlarged to full circles.
The product kernel costs
$CN$ by \eqref{eq:iso-mixed-gap}, and the angular Jacobian costs
$C^N$.  Cauchy's inequality on the protected disk and $N!\ge(N/e)^N$
give
\begin{equation}\label{eq:iso-mixed-jets}
 |\partial_{\mathsf S}^j F_N(\mathsf S_\eps)|
       \le C_j^N N^{-N/2+C_j},\qquad N\ge2.
\end{equation}
For each fixed $N=2,3$, the displaced branches and product gap instead
give a bounded analytic integrand on the compact mixed support.
Every fixed derivative of these pieces is bounded.

\subsubsection{The regular inverse and a cell Gaussian bound}

Write $z=e^{-i\varphi}$, $\gamma=i\varphi$, and set $v=\varphi+w$.

\begin{lemma}[One-sided inverse]\label{lem:iso-inverse}
There is a fixed complex neighborhood of the limiting one-sided
curve on which $w=U(\mathsf S,v)$ is analytic with bounded fixed
joint jets.  On a fixed convex neighborhood $\mathcal D_R$ of the
right $v$-interval, including a short negative branch segment,
$\Re U_{vv}\ge c>0$.  All the actual one-sided images, including
the endpoint slivers, lie in these neighborhoods when the fixed
geometric choices are small.
\end{lemma}
\begin{proof}
On $[-b,0]$, $g(\theta)=\arccos(\mathsf S_*-\cos\theta)$ has
$g'=-\sin\theta/\sin g\ge0$, and $g'(0)=0$.
Thus $v=g(\theta)+\theta$ is strictly increasing, also on a short
upper sliver.  At the branch,
$w=-\arccos(\mathsf S-\cos\varphi)$ is even analytic in
$\varphi$, and $\dd(w+\varphi)/\dd\varphi=1$ at zero.
The explicit right inverse is
\begin{equation}\label{eq:iso-inverse-explicit}
 U(\mathsf S,v)=\frac v2-
       \arccos\frac{\mathsf S}{2\cos(v/2)}.
\end{equation}
Its limiting real interval is strictly inside $(-2\alpha,2\alpha)$,
because $b<2\alpha$.  With $x=v/2$ and
$D=4\cos^2x-\mathsf S_*^2>0$, differentiation gives
\[
 U_{vv}(\mathsf S_*,v)
 =\frac{\mathsf S_*}{4\cos^2x\sqrt D}
      +\frac{\mathsf S_*\sin^2x}{D^{3/2}}>0.
\]
Compactness preserves a positive real part on a thin convex rectangle.
On the left arc $\varphi=-i\gamma$, $\gamma>0$; the real part of $v$
is $\theta$, and $1+\varphi_w$ cannot vanish because $\varphi_w$
is purely imaginary.  The left and right embedded arcs meet in the
regular branch chart.  Their local analytic inverses consequently
patch on a sufficiently small neighborhood of the compact arc.
Short endpoint extensions preserve injectivity.  Square-root
continuity near the branch, and ordinary continuity elsewhere, then
put every sufficiently small deformed image inside the chosen domains.
\end{proof}

Put $\phi(\mathsf S,v)=v-U(\mathsf S,v)$.  With the square-root
branch fixed by $\sqrt{4\cos^2(v/2)-\mathsf S^2}=\sin\phi-\sin U$,
implicit differentiation gives
\begin{equation}\label{eq:iso-regular-factors}
 U_v=\frac{\sin\phi}{\sqrt{4\cos^2(v/2)-\mathsf S^2}},\quad
 U_{\mathsf S}=\frac1{\sqrt{4\cos^2(v/2)-\mathsf S^2}},\quad
 a(\mathsf S,v):=\frac{U_v}{\sinh(i\phi)}
     =\frac{-i}{\sqrt{4\cos^2(v/2)-\mathsf S^2}}.
\end{equation}
The denominator is separated on the fixed inverse neighborhoods.
The complete pair is analytic there, including at the double branch:
the dispersion identity also writes it locally as
\[
 P(v_i,v_l)=-\frac{\sin^2((\phi_i-\phi_l)/2)}
                         {\sin^2((U_i+U_l)/2)},
\]
whose denominator is nonzero near $U_i=U_l=-b$.
Away from that double branch the original pair denominator is
separated.  The diagonal zero therefore gives the global regular
factorization
\begin{equation}\label{eq:iso-pair-zero}
 P(v_i,v_l)=(v_i-v_l)^2 A(\mathsf S,v_i,v_l),
\end{equation}
where $A$ and its fixed joint jets are bounded on a smaller
neighborhood of the one-sided arc squared.

\begin{lemma}[Cell contraction and retained collision factors]
\label{lem:iso-cell-gaussian}
On a fixed common complex neighborhood of the one-sided arc, for
every fixed $\ell$ the quotient
\[
 \frac{\prod_{i<l\le N}P_{il}}
      {\prod_{1\le i<l\le\ell}(v_i-v_l)^2}
\]
is analytic, and each of its fixed joint jets is bounded by
$C_{r,\ell}\exp(-\kappa N^2+C_{r,\ell}N)$ for $N\ge\ell$.
The same bound holds after multiplying by $\prod_i a(\mathsf S,v_i)$.
It applies on the enlarged right rectangle with arbitrary actual
occupancy fractions in $[0,1]$.
\end{lemma}
\begin{proof}
On the limiting arc $-\pi\le\theta\le0$, both $y=e^{i\theta}$ and
$z$ are in the closed lower half disk: $z$ is on its unit semicircle
for $-b\le\theta\le0$ and on $(0,1]$ on the remaining segment.
Equations \eqref{eq:iso-halfdisk-identity} and $P=A(z_i,z_l)A(y_i,y_l)$
give $|P|\le1$.  Apparent endpoint quotients are filled using the
analytic complete pair.  Since $P(v,v)=0$, finitely many cells,
say $K$, cover the compact regular arc so that same-cell pairs have
$|P|\le\eta<1$.  Choose a small common complex neighborhood in
$(\mathsf S,v)$: same-cell pairs then have modulus at most
$\eta'<1$, and all pairs have modulus at most $M=1+o(1)$.
For populations $N_1,\ldots,N_K$,
\[
 \sum_a\binom{N_a}{2}\ge\tfrac12(N^2/K-N),\qquad
 \prod_{i<l}|P_{il}|
 \le M^{\binom N2}(\eta'/M)^{(N^2/K-N)/2}.
\]
The neighborhood can be fixed small enough that the coefficient of
$N^2$ in its logarithm is negative.  This proves
$\exp(-\kappa N^2+CN)$, including on slightly larger fixed polydisks.

Factor the internal $\ell$ pairs by \eqref{eq:iso-pair-zero}.
They cost a fixed constant, their cross pairs with the remaining
coordinates cost $C^{\ell(N-\ell)}$, and the remaining complete
pair product retains its Gaussian bound.  Cauchy's inequality on the
fixed polydisks gives all asserted joint jets.  The bounded one-body
product adds only $C^N$.

For completeness, this common neighborhood contains the domains used
after absolute integration enlarges the right support.  First fix the
cells and inverse-domain margins on the undeformed compact arc.  Choose
the endpoint slivers and the negative branch segment inside those
margins.  In the full right angular rectangle, every actual
$\rho=P/N$ and $\upsilon=Q/N$ is in $[0,1]$, and
$|h_i|\le\tau$.  Away from the branch the change in $v_i$ is
$O(\eps+\tau)$; near it the uniform square-root bound makes the
change $O(\sqrt{\eps+\tau})$.  Choosing $\tau$ and then $\eps_0$
small puts all such images in the same fixed neighborhood.  No
one-sided occupancy restriction is used for this extension.
\end{proof}

\subsubsection{The exact current and its joint jets}

For $\chi_{\mathsf S}(\theta)=(v_i)$ define the scalar
\[
 J=\det D_\theta\chi_{\mathsf S}
    =\det M\prod_i U_{v,i}^{-1}.
\]
The physical density equals $Jf$, where
\begin{equation}\label{eq:iso-regular-density}
 f=\frac1{N!(2\pi)^{N-1}}
       \frac{1+\sigma Z}{1-\sigma Z}
       \prod_i a(\mathsf S,v_i)\prod_{i<l}P_{il}.
\end{equation}
This is an equality of scalar coefficients on the original constrained
angles: $J\prod a=(\det M)\prod(\sinh\gamma_i)^{-1}$.
No intrinsic complex measure on a new constrained $v$ hyperplane is
asserted.

For a regular-coordinate field $B$ put
\begin{equation}\label{eq:iso-angular-field}
 V=-M^{-1}(U_{\mathsf S}+U_vB).
\end{equation}
If $\sum_iU_{v,i}B_i=-\sum_iU_{\mathsf S,i}$, then $\sum_iV_i=0$.
If also $\sum_iB_i=0$, the operator
$\partial_{\mathsf S}-V\cdot\nabla_\theta$ freezes $\sum_i\varphi_i$.
Indeed $\sum\varphi_i=\sum v_i-\sum w_i$, and $\sum w_i=\sum\theta_i$.
The exact scalar current identity is
\begin{equation}\label{eq:iso-current}
 \partial_{\mathsf S}(Jb)-\operatorname{div}_\theta(VJb)
  =J\left(E_{\mathsf S}b+\sum_i B_iE_i b
                           +\Bigl(\sum_iE_iB_i\Bigr)b\right).
\end{equation}
Here $b$ may already be a parameter-dependent transported amplitude.
The commuting regular frame is
\[
 E_i=\sum_l(M^{-1})_{li}U_{v,i}\partial_{\theta_l},\qquad
 E_{\mathsf S}=\partial_{\mathsf S}
                  +\sum_l(M^{-1}U_{\mathsf S})_l\partial_{\theta_l}.
\]
It satisfies $E_i v_l=\delta_{il}$ and $E_{\mathsf S}v_l=0$.
Every commutator is a linear combination of the angular derivatives,
annihilates all $v_l$, and hence vanishes.  Thus this frame applies
also to smooth real cutoffs without analytically continuing them.
To verify \eqref{eq:iso-current}, differentiating the coordinate
Jacobian and using Jacobi's determinant formula gives
$J\operatorname{div}_vB=\partial_{\mathsf S}J-\operatorname{div}_\theta(JV)$;
the product rule gives the displayed identity.  Tangency makes its
ambient divergence exactly the constrained one, since
\[
 \sum_{i<N}(\partial_i-\partial_N)V_i=\sum_i\partial_iV_i.
\]
Integration by parts is applied to the real and imaginary parts of
these complex currents.

Choose one smooth right cutoff $r$, equal to one for
$\theta+b\ge-\delta_B/2$ and zero for $\theta+b\le-\delta_B$,
with its periodic return in the excluded upper plateau.  Decompose
$1=\prod_i r_i+(1-\prod_i r_i)$ on the one-sided term and telescope
the second part.  Each of its at most $N$ terms has a named compact
left anchor $a$, also on all cutoff jets.  It has
$|z_a|\le q_0<1$ and $|U_{v,a}|\ge c$.  Use
\begin{equation}\label{eq:iso-anchor-field}
 B_a=-\frac{\sum_iU_{\mathsf S,i}}{U_{v,a}},\qquad B_i=0\ (i\ne a).
\end{equation}
This gives tangency; phase freezing is unnecessary because
$|1-\sigma Z|\ge1-q_0$.  All fixed regular-frame jets of $B,V$ and
the cutoffs have polynomial $N$ bounds.  The cell Gaussian, the
bounded regular one-body factors, and \eqref{eq:iso-current} then
bound every fixed derivative of this whole piece by
\begin{equation}\label{eq:iso-anchor-bound}
 \frac{C_jN^{C_j}}{N!}\exp(-\kappa N^2+C_jN).
\end{equation}
In fact, eliminate the compact anchor angle with the constraint.
Its Jacobian factor is bounded, and every remaining factor has the
independent integrable majorant
$C(1+\operatorname{dist}(\theta,-b)^{-1/2})$.
Their integrals cost only $C^N$.  The product kernel and its
transported derivatives have polynomial $N$ bounds on this anchored
support.  This argument applies also at each fixed $N=2,3$.

We now work on the full right part, with all angles in
$I=[-b-\delta_B,\delta]$.  Set $u_i=\theta_i+b$ and, on a diameter
shell $r_0\le d:=\max\theta_i-\min\theta_i\le2r_0$, put
\begin{equation}\label{eq:iso-scales}
 \mu=N^{-1}\sum_i u_i,\qquad
 T=|\mu|+r_0+\eps,\qquad \Delta=r_0/\sqrt T.
\end{equation}
These are pointwise envelopes; no derivative of a mean, diameter,
or scale is ever taken.

The branch normal form, for a common outward displacement $h\ge0$,
is
\begin{equation}\label{eq:iso-branch-normal}
 \begin{gathered}
 1-W=\sin b\,u-i\beta+R,\qquad
 c(\eps+h)\le\beta\le C(\eps+h),\\
 R\in\mathbb R,\qquad |R|\le C(u^2+\eps^2+h^2),\qquad
 \varphi=\sqrt{2(1-W)}(1+O(1-W)).
 \end{gathered}
\end{equation}
Indeed $W=\mathsf S-\cosh(h+i(-b+u))$ and
$\beta=\Im\mathsf S+\sinh h\sin(b-u)$; its real part has the
displayed Taylor expansion.  The physical square root lies in a
small enlargement of the fourth quadrant.  In particular
\[
 \begin{gathered}
 |1+\varphi_w(w_i)|\le C(|u_i|+\eps)^{-1/2},\quad
 -\Im\varphi_i\ge c\sqrt{|u_i|+\eps}\quad(u_i\le0),\\
 \int_{\mu-2r_0}^{\mu+2r_0}(|u|+\eps)^{-1/2}\dd u
       \le Cr_0/\sqrt T.
 \end{gathered}
\]
The interval bound follows by comparability if $|\mu|\ge4r_0$;
otherwise enlarge to $[-6r_0,6r_0]$ and rationalize
$\sqrt{6r_0+\eps}-\sqrt\eps$.
Joining two angles along their one-variable curve with the actual
$\rho=P/N,\upsilon=Q/N$ held fixed gives
\begin{equation}\label{eq:iso-coordinate-scales}
 |v_i-v_l|\le C\Delta,\qquad |U_{v,i}|\le C\sqrt T.
\end{equation}
The joining derivative has the additional bounded factor
$1-i\tau(\rho m'-\upsilon p')$.  Near the lower branch its outward
shift is the common $\tau\rho$, so the preceding interval bound
applies; elsewhere it is regular.  For the second bound, if $P=0$
then every actual $h_i=0$, and the branch form and
\eqref{eq:iso-regular-factors} apply directly.  If $P>0$, some
actual angle is positive, so $\max u_i\ge b$ and $T\ge b/2$;
boundedness of $U_v$ proves the same result.  These arguments allow
all actual occupancy fractions in $[0,1]$.

Use the filled analytic divided differences
\[
 \mathcal U_{il}=\int_0^1U_v(v_l+t(v_i-v_l))\,dt,\qquad
 \mathcal E_{il}=\int_0^1U_{vv}(v_l+t(v_i-v_l))\,dt.
\]
The convex inverse domain gives $\Re\mathcal E_{il}\ge c$.
Define the analytic field
\begin{equation}\label{eq:iso-covariance-field}
 \begin{gathered}
 D_U=\sum_{i<l}(U_i-U_l)^2,\qquad
 H_{il}=\frac{(v_i-v_l)\mathcal U_{il}^2}{\mathcal E_{il}},
 \qquad G_i=\sum_lH_{il},\\
 A=-\sum_iU_{\mathsf S,i},\qquad B_i=AG_i/D_U.
 \end{gathered}
\end{equation}
There is no pair pole: $H_{il}$ is analytic also at $v_i=v_l$.
Away from equality it equals $(U_i-U_l)^2/(U_{v,i}-U_{v,l})$.
Antisymmetry and pairing give
\[
 \sum_iG_i=0,\qquad \sum_iU_{v,i}G_i=D_U,\qquad
 \sum_iB_i=0,\qquad \sum_iU_{v,i}B_i=-\sum_iU_{\mathsf S,i}.
\]
Thus the field has both exact transport identities.

\begin{lemma}[Protected analytic variance and joint jets]
\label{lem:iso-covariance-jets}
On each actual right shell,
\begin{equation}\label{eq:iso-covariance-lower}
 \Re D_U\ge c\sum_{i<l}(\theta_i-\theta_l)^2
                   \ge cT\Delta^2>0.
\end{equation}
For every fixed total derivative order $s$,
\begin{equation}\label{eq:iso-field-jets}
 |\mathcal D^s B_i|+|\mathcal D^s V_i|
       \le C_sN^{C_s}\Delta^{-1-s}.
\end{equation}
The field is regular at every partial diagonal.
\end{lemma}
\begin{proof}
On the actual contour $U_i=w_i=\theta_i-ih_i$, and the real
balanced shifts satisfy
$|h_i-h_l|\le C\tau|\theta_i-\theta_l|$ at fixed actual occupancies.
Hence
$\Re(U_i-U_l)^2=(\theta_i-\theta_l)^2-(h_i-h_l)^2
 \ge(1-C^2\tau^2)(\theta_i-\theta_l)^2$.
Small fixed $\tau$ proves the denominator bound.

Freeze the actual values of $T,\Delta$ and consider the polydisk
\[
 |\delta\mathsf S|,\ |\delta v_i|\le c_0\Delta/N^2.
\]
Choose $c_0$ small so it lies in the fixed inverse neighborhoods.
Throughout this disk, $|v_i-v_l|\le C\Delta$,
$|U_v|\le C\sqrt T$, and
$|U_{\mathsf S}(v_i)-U_{\mathsf S}(v_l)|\le C\Delta$.
Indeed the last bound is an analytic divided difference; the middle
one follows from its actual endpoint bound and bounded joint jets,
since $\Delta\le\sqrt T$.
Along the straight parameter and coordinate segments, the common
parameter change therefore moves each $U_i-U_l$ by $O(\Delta
|\delta\mathsf S|)$, and the two coordinate changes move it by
$O(\sqrt T(|\delta v_i|+|\delta v_l|))$.
Their total change is at most $Cc_0r_0/N^2$.
Since each original difference is $O(r_0)$, summing the $N(N-1)/2$
squared differences gives
\[
 |D_U^{\rm new}-D_U|\le Cc_0r_0^2\le\tfrac12\Re D_U.
\]
Thus $|D_U^{\rm new}|\ge cr_0^2$ on this whole polydisk.
Also $|\mathcal U_{il}|\le C\sqrt T$ and
$|\mathcal E_{il}|\ge c$ there, so
$|H_{il}|\le Cr_0\sqrt T=CT\Delta$, $|A|\le CN$, and
$|B_i|\le CN^2/\Delta$.
Cauchy's inequality now gives all fixed joint jets at once:
\begin{equation}\label{eq:iso-covariance-jet-table}
 |\partial_{\mathsf S}^a\partial_v^\beta B_i|
 \le C_sN^{2+2s}\Delta^{-1-s},\qquad s=a+|\beta|.
\end{equation}
No scale or polydisk radius has been differentiated.

The commuting-frame coefficients
$M^{-1}\operatorname{diag}U_v$ and $M^{-1}U_{\mathsf S}$ have
polynomially bounded fixed joint jets: the real angular jets of
$M$ and analytic jets of $U$ are bounded, and differentiating
$M^{-1}$ proves the assertion inductively without an inverse $U_v$.
This also bounds all fixed smooth cutoff jets.
Formula \eqref{eq:iso-angular-field}, the product rule, and
\eqref{eq:iso-covariance-jet-table} give the same field bound for $V$.
The positive actual denominator and filled analytic $H_{il}$ show
that no partial diagonal is singular.
\end{proof}


Fix now $\ell=n_*+1$ and $D_\ell=\ell(\ell-1)$.
Factor the first $\ell$ collision variables using
Lemma~\ref{lem:iso-cell-gaussian}.  For the initial smooth weighted
regular amplitude $b_0$ after removing the product kernel,
\[
 |\mathcal D^s b_0|
 \le\frac{C_sN^{C_s}}{N!}e^{-\kappa N^2+C_sN}
                     \Delta^{D_\ell-s},\qquad s\le2k.
\]
The same statement includes parameter derivatives, which preserve the
explicit collision polynomial.  Iterate the exact current with
\begin{equation}\label{eq:iso-transport-recurrence}
 b_{j+1}=\mathcal T b_j,\qquad
 \mathcal T b=E_{\mathsf S}b+\sum_iB_iE_i b
                              +\Bigl(\sum_iE_iB_i\Bigr)b.
\end{equation}
This recurrence differentiates every parameter-dependent descendant
coefficient and cutoff.  Because the field freezes $\sum\varphi_i$,
the product kernel is unchanged.  Induction using
\eqref{eq:iso-field-jets} gives
\begin{equation}\label{eq:iso-transported-amplitude}
 |\mathcal D^s b_j|
 \le\frac{C_jN^{C_j}}{N!}e^{-\kappa N^2+C_jN}
              \Delta^{D_\ell-2j-s},\qquad s\le2(k-j).
\end{equation}
Indeed $E_{\mathsf S}$ loses at most one scale power, while
$B_iE_i$ and $(E_iB_i)$ each lose at most two.  Leibniz' rule gives
the same count for joint jets.  Only polynomially many index choices
occur at each fixed order.

The boundary flux vanishes.  At fixed $N,\eps>0$ puncture
the finitely many full-equality points by angular balls of radius
$r$.  All geometric Jacobians and the product kernel are then smooth;
$V=O(r^{-1})$, and the retained collision degree is $D_\ell$.
The flux after $j$ transports is bounded by a constant times
$r^{D_\ell+N-3-2j}$, which tends to zero for every $N\ge n_*+2$
and $j\le k$.  Partial diagonals require no puncture by
\eqref{eq:iso-covariance-lower}.  The original smooth weights vanish
at the outer chart boundaries; no diameter cutoff has been
differentiated.  Thus the $j$th derivative of the complete right term
is the integral of its transported density, bounded by
\begin{equation}\label{eq:iso-transported-density}
 \frac{C_jN^{C_j}}{N!}e^{-\kappa N^2+C_jN}
 \Delta^{D_\ell-2j}
 \frac{\prod_i|1+\varphi_w(w_i)|}{|1-\sigma Z|}.
\end{equation}

\subsubsection{Coarea with the actual occupancy dependence}

We prove the integral estimate needed for
\eqref{eq:iso-transported-density}; in particular the occupancies
will not be treated as constants when differentiating the phase.

First, for $G=x-ia$, $a\ge a_0>0$, and fixed
$|\vartheta|<\pi/2$, put $\nu=\Re(e^{-i\vartheta}G)$ and
$P_\vartheta=2\pi\cos\vartheta$.  If $a_0\le1$, then
\begin{equation}\label{eq:iso-oblique-kernel}
 \frac1{|1-e^{-iG}|}
 \le\frac{C_\vartheta}
 {a_0+\operatorname{dist}(\nu,P_\vartheta\mathbb Z)}.
\end{equation}
For $a\le1$, this follows from
$|1-e^{-a-ix}|\ge c(a+\operatorname{dist}(x,2\pi\mathbb Z))$
and
$|\nu-2\pi n\cos\vartheta|\le
 \cos\vartheta|x-2\pi n|+|\sin\vartheta|a$.
For $a\ge1$, the left denominator is separated and the proposed
right denominator is bounded above.  Integrating the majorant over
an interval of length $L$ costs at most
$C_\vartheta(1+L)(1+|\log a_0|)$.  For $\sigma=-1$ replace $G$
by $G-\pi$.  On the actual contours $-\log|Z|\ge cN\eps$:
write $\varphi_i=x_i-ia_i$, and use
$\Im W_i=\sin x_i\sinh a_i\le Ca_i$ on the bounded phase range.
Apply the lemma with $a_0=\min(1,cN\eps)$.  Its logarithmic cost is
at most $C(1+L)(1+|\log\eps|)$.

\begin{lemma}[Complete right-shell coarea]\label{lem:iso-coarea}
On each fixed-mean lift of the full right rectangle
$I^N\cap\mathbb H_N^\sigma$, with the actual balanced occupancies,
the integral over any diameter shell in \eqref{eq:iso-scales} obeys
\begin{equation}\label{eq:iso-shell-coarea}
 \int_{\substack{\sum u_i=N\mu\\r_0\le d\le2r_0}}
  \frac{\prod_i|1+\varphi_w(w_i)|}{|1-\sigma Z|}
       \prod_{i<N}\dd\theta_i
 \le C^NN^C\Delta^{N-3}(1+|\log\eps|).
\end{equation}
There are at most $CN$ constrained means.  The constants are uniform
in each of them; their subsequent sum costs only this polynomial
factor.
\end{lemma}
\begin{proof}
Write $m_0=\min u_i$, $M_0=\max u_i$, and first suppose
$0\le m_0<M_0$.  With other angles fixed, the constraint fixes
$m_0+M_0$.  Put $G=\sum_i\varphi_i$ and $\vartheta=3\pi/8$.
We claim
\begin{equation}\label{eq:iso-extreme-ratio}
 \Re(e^{-i\vartheta}\dd G/\dd m_0)>0,\qquad
 \frac{|(1+\varphi_w(w_{m_0}))(1+\varphi_w(w_{M_0}))|}
      {\Re(e^{-i\vartheta}\dd G/\dd m_0)}
 \le C\frac{\sqrt{M_0+\eps}}{M_0-m_0}.
\end{equation}

For its proof only, freeze the actual values $\rho=P/N,\upsilon=Q/N$
and set
$f(u)=-b+u-i\tau\{\rho m(-b+u)-\upsilon p(-b+u)\}$,
$E=\eps+\tau\rho$.  For $u\ge0$, the branch normal form and
compactness on the rest of $I$ give
\[
 |\varphi_w(f(u))|\le C(u+E)^{-1/2},\qquad
 |\varphi_{ww}(f(u))|\asymp(u+E)^{-3/2}.
\]
The inverse supplies
$-\varphi_{ww}=U_{vv}/U_v^3$.  For $u\ge0$ the argument of $U_v$
is in a small enlargement of $[-\pi/4,0]$, while $U_{vv}$ has small
argument.  Hence $-\varphi_{ww}$ has uniformly positive
$3\pi/8$ projection.  If $a=f'=1+O(\tau)$ and
$B(u)=\varphi_w(f(u))a(u)$, then
$-B'=-\varphi_{ww}(f)a^2-\varphi_w(f)a'$ retains that projection
for small $\tau$.  Consequently
\begin{equation}\label{eq:iso-fixed-occupancy-curvature}
 \begin{gathered}
 \Re\{e^{-i\vartheta}[B(m_0)-B(M_0)]\}\ge cL,
 \\
 L=\int_{m_0}^{M_0}(u+E)^{-3/2}\dd u
 =\frac{2(M_0-m_0)}{\sqrt{m_0+E}\sqrt{M_0+E}
                         (\sqrt{m_0+E}+\sqrt{M_0+E})}.
 \end{gathered}
\end{equation}

Now restore all occupancy derivatives.  Put
$B_m^{\rm occ}=\sum_lm(\theta_l)\varphi_w(w_l)$ and
$B_p^{\rm occ}=\sum_lp(\theta_l)\varphi_w(w_l)$.
Direct differentiation of \eqref{eq:iso-balanced-deformation} gives
the exact identity
\begin{equation}\label{eq:iso-occupancy-derivative}
 \partial_{\theta_i}G
 =\varphi_w(w_i)[1-i\tau(\rho m'_i-\upsilon p'_i)]
       -\frac{i\tau p'_i}{N}B_m^{\rm occ}
       +\frac{i\tau m'_i}{N}B_p^{\rm occ}.
\end{equation}
Every actual $u_l\ge m_0$, so
$(|B_m^{\rm occ}|+|B_p^{\rm occ}|)/N\le C(m_0+E)^{-1/2}$.
Taking $\partial_{\theta_{\min}}-\partial_{\theta_{\max}}$,
both bump-derivative differences are $O(M_0-m_0)$.
The entire occupancy correction is therefore at most
$C\tau(M_0-m_0)(m_0+E)^{-1/2}\le C\tau L$, by the exact
expression for $L$ and the fixed upper bound on $M_0+E$.
Small fixed $\tau$ absorbs it into the positive projection.
The numerator of \eqref{eq:iso-extreme-ratio} is at most
$C[(m_0+E)(M_0+E)]^{-1/2}$, giving ratio
$C\sqrt{M_0+E}/(M_0-m_0)$.  If $P=0$, $E=\eps$; if $P>0$,
an actual positive angle gives $M_0\ge b$, and hence
$E\le C(M_0+\eps)$.  This proves \eqref{eq:iso-extreme-ratio}
with all actual occupancies, including on the enlarged rectangle.

After transport, bound every smooth coefficient by its proved
envelope and enlarge to this full rectangle.  All its occupancy
fractions lie in $[0,1]$, the balanced displacement preserves
$\Im W_i\ge\Im\mathsf S$, and Lemma~\ref{lem:iso-cell-gaussian}
has already placed its images in the regular domains.
Fix a constrained mean $\mu=b+(2\pi l+a_\sigma)/N$ and the two
extreme indices; at most $CN^3$ such choices occur.  Eliminate the
maximum coordinate with the constraint.  With the remaining $N-2$
coordinates fixed, conditions specifying the extrema, the shell,
and a nonnegative minimum are linear interval conditions in $m_0$,
with $M_0=A-m_0$.  Thus the fiber is one interval, and its projected
phase is injective by \eqref{eq:iso-extreme-ratio}.  Its phase range
is at most $CN$.  Equations \eqref{eq:iso-oblique-kernel} and
\eqref{eq:iso-extreme-ratio} bound this extreme-coordinate integral by
\[
 CN\frac{\sqrt T}{r_0}(1+|\log\eps|).
\]

If $m_0<0$, its branch attenuation from
\eqref{eq:iso-branch-normal} instead gives the direct bound
\[
 \frac{|(1+\varphi_w(w_{m_0}))(1+\varphi_w(w_{M_0}))|}
      {|1-\sigma Z|}
 \le\frac{C}{(|m_0|+\eps)\sqrt{|M_0|+\eps}}.
\]
With $a=|m_0|+\eps$ and $c=|M_0|+\eps$, $a+c\ge r_0$, and
\[
 \frac1{a\sqrt c}\le\frac1{\sqrt{r_0}}\left(\frac1a+\frac1c\right).
\]
Each reciprocal integrates to $C(1+|\log\eps|)$ on the bounded
slice $M_0=A-m_0$.  Since $T\ge r_0$, this gives the same extreme
pair envelope $C\sqrt T/r_0(1+|\log\eps|)$.
Only now replace the other $N-2$ Jacobian factors by the independent
majorants and integrate each over $[\mu-2r_0,\mu+2r_0]$.
Their total cost is $(Cr_0/\sqrt T)^{N-2}$.
Combining the extreme pair, the ordinary factors, and the polynomial
number of cells proves \eqref{eq:iso-shell-coarea}.
\end{proof}

\subsubsection{One sum for all higher particle orders}

Combining \eqref{eq:iso-transported-density} with
Lemma~\ref{lem:iso-coarea} bounds a right shell at a fixed mean by
\begin{equation}\label{eq:iso-one-shell}
 \frac{C_jN^{C_j}}{N!}e^{-\kappa N^2+C_jN}
 \Delta^{D_\ell+N-3-2j}(1+|\log\eps|).
\end{equation}
Its exponent is positive for every $N\ge n_*+2$ and $j\le k$:
it is at least $n_*+N-2$ in the odd case and $n_*+N-1$ in the
even case.  For fixed $\mu,\eps$,
\[
 \frac{\Delta(r_0/2)}{\Delta(r_0)}
 =\frac12\sqrt{\frac{|\mu|+r_0+\eps}{|\mu|+r_0/2+\eps}}
       \in[1/2,1/\sqrt2].
\]
The dyadic sum is therefore bounded uniformly in the mean by a geometric series.
Its largest possible scale is fixed; its power costs only $C^N$
times a fixed constant.  This is absorbed into the linear exponential
in \eqref{eq:iso-one-shell}.  Summing the at most $CN$ means adds
only a polynomial factor. The mixed, anchored, and right pieces
consequently obey the same bound for all higher orders,
\begin{equation}\label{eq:iso-every-order-tail}
 |\partial_{\mathsf S}^j\mathcal F_N^\sigma(\mathsf S_\eps)|
 \le C_j(1+|\log\eps|)
 \left(C_j^NN^{-N/2+C_j}
       +\frac{N^{C_j}}{N!}e^{-\kappa N^2+C_jN}\right).
\end{equation}
Both envelopes are summable.  The exact split
\eqref{eq:iso-exact-split} proves Theorem~\ref{thm:iso-tail}.
The parameter-change assertion follows from the finite chain rule.

\begin{proposition}[Finite-order transport interface]
\label{prop:iso-finite-transport}
For each fixed $N\ge2$ and every fixed derivative order, the mixed
piece is bounded and extends holomorphically across
$\mathsf S_*$ on its deformed contour.  A smooth one-sided term with
a compact left anchor has bounded derivatives of every fixed order.
A smooth right-supported term whose angular diameter is bounded
below by a fixed $d_0>0$ has derivatives
$O(1+|\log\eps|)$ of every fixed order.  The assertions allow
arbitrary fixed smooth angular cutoffs with those supports.
\end{proposition}
\begin{proof}
For the fixed mixed order, both true branches and the product
denominator have uniform gaps on the compact support, as proved
above; the same small parameter disk works for the whole integrand.
This includes $N=2,3$ directly.  The anchored field
\eqref{eq:iso-anchor-field} and its integrable independent Jacobian
majorants prove the second assertion at any fixed derivative budget.
On a right support with $d\ge d_0$, all jets of the covariance
field are bounded and there is no collision puncture.  Iterate
\eqref{eq:iso-transport-recurrence} with that budget, and then apply
Lemma~\ref{lem:iso-coarea}; the finitely many dyadic scales bounded
away from zero cost $O(1+|\log\eps|)$.  For $N=2$ the coarea keeps
both extreme factors and has zero remaining one-body integrations,
so the same argument applies.  All differentiated cutoff factors
are included in the exact recurrence.
\end{proof}


\subsection{All lower coefficients and the complete finite complement}

\begin{lemma}[Finite complement at every fixed derivative order]
\label{lem:iso-finite-complement}
At any of the signed positive-geometry points just specified, every
fixed $\mathsf S$ derivative of every smaller particle coefficient of
the relevant parity is bounded along \eqref{eq:iso-budget}.
At $n=n_*$ choose a fixed smooth local shape cutoff at
$\theta_1=\cdots=\theta_n=-\alpha_+$, and let $I_n^\sigma$ be
the resulting local integral.  Then every fixed derivative of
$\mathcal F_n^\sigma-2I_n^\sigma$ is bounded on the same path.
These conclusions also hold for the exterior-coordinate derivatives.
\end{lemma}
\begin{proof}
We first prove $O(1+|\log\eps|)$ at every fixed derivative order;
constants may depend on this finite particle order and derivative
order.  The case $n=1$ follows from \eqref{eq:iso-one-particle}.
For all $n\ge2$ use the whole-torus balanced deformation and exact
reflection decomposition from the preceding transport argument:
\[
 \mathcal F_n^\sigma=F_n^\sigma+2L_n^\sigma.
\]
The mixed-population piece has both occupancies at least $1/4$.
At this fixed $n$, both true branches are displaced by at least
$\tau/(4n)$, and the product kernel has a fixed damping gap.
The complete density is therefore holomorphic and bounded on a
fixed parameter neighborhood.  This includes $n=2,3$ by the same compactness argument.

Each left-root term in $L_n^\sigma$ has a named compact anchor with
$|\rho_a|\le q_0<1$ and $|U_{v,a}|\ge c>0$.
The anchor field $B_a=-\sum_iU_{\mathsf S,i}/U_{v,a}$, $B_i=0$
for $i\ne a$, makes the angular representative tangent to the
constraint.  The current identity \eqref{eq:iso-current} leaves all
fixed regular derivatives bounded and the product kernel separated.
Eliminate the compact anchor angle using the exact constraint.
Its Jacobian factor is bounded; every remaining factor is bounded by
$C(1+\dist(\theta,-b)^{-1/2})$, independently integrable.
This proves boundedness of every fixed derivative of the entire
left-root part, including its cutoff derivatives.

It remains to analyze the right part.  Choose its branch width below
one quarter of the finite mean gap for all relevant $n\le n_*$.
Every constrained equality configuration on this support is then a
compact regular point away from the branch.
At full equality the balanced shift is exactly zero: writing its
upper and lower bumps as $p,m$, one has
$P=np(\theta_0)$, $Q=nm(\theta_0)$ and
$P m(\theta_0)-Q p(\theta_0)=0$.
Its product denominator can vanish only when
\[
 \mathsf S_* = \cos\theta_0+\cos\varphi_0,
 \qquad n\theta_0,n\varphi_0\in2\pi\Z+a_\sigma.
\]
The prime arithmetic excludes every smaller order.  At $n_*$ it
leaves only $-\alpha_+$ on the one-sided support; $+\alpha_+$ is
in the excluded upper plateau.  Each noncritical equality admits a
fixed neighborhood on which all roots and pair factors are regular
and $1-\sigma Z$ is separated.  Its localized integral is
holomorphic through the parameter point.

Remove these finitely many neighborhoods and, at $n_*$, the critical
chart.  Take the latter wholly inside the lower plateau $m=1,p=0$,
where the deformation is identically zero, the one-sided weight is
one and the right cutoff is one.  The remaining right support and
all its cutoff jets have angular diameter $d\ge d_0>0$.
The collective field in Proposition~\ref{prop:iso-finite-transport}
has bounded fixed jets there, with no collision puncture or flux.
After any fixed number of applications of \eqref{eq:iso-current},
the density is bounded by
\[
 C_{n,j}\frac{\prod_i|1+\varphi_w(w_i)|}{|1-\sigma Z|}.
\]
Lemma~\ref{lem:iso-coarea} absorbs the two actual extreme-coordinate
Jacobian factors before any independent one-body majorization.
There are finitely many means and diameter shells with $d\ge d_0$,
so their integrals are $O(1+|\log\eps|)$.  When $n=2$ there are
no remaining one-body integrals, and the same two-extreme estimate
applies.  Smooth cutoff fibers are enlarged to the geometric
coarea intervals only after differentiation.  Thus the bound covers
all remaining phase zeros, not merely collision neighborhoods.

For the strengthening from logarithmic to bounded derivatives, let
$R$ be any complete lower coefficient or the complete remainder
$\mathcal F_{n_*}^\sigma-2I_{n_*}^\sigma$.
The contours and weights defining $R$ are fixed independently of
$\mathsf S$, so it is holomorphic on the one-sided sector.  For each
fixed $j$,
\[
 \partial_{\mathsf S}^{j}R(\mathsf S_\eps)
 =\partial_{\mathsf S}^{j}R(\mathsf S_{\eps_1})
  -\int_\eps^{\eps_1}
      \mathsf S'_u\,
      \partial_{\mathsf S}^{j+1}R(\mathsf S_u)\,du.
\]
The next-derivative logarithmic bound is integrable, and
$\mathsf S'_u$ is bounded.  This proves the assertion without
claiming a holomorphic extension of $R$ through the point.
The finite chain rule gives exterior-coordinate derivatives.
The factor two is exactly the reflection factor in
$\mathcal F_n^\sigma=F_n^\sigma+2L_n^\sigma$.
\end{proof}

\subsection{The first singular coefficient with its exact phase}

The following calculation uses only the complete product
\eqref{eq:iso-product}.  Let $N=n_*$, $\alpha=\alpha_+$ and
$N\alpha\in2\pi\Z+a_\sigma$.
In its lower chart set
\[
 \theta_i=-\alpha+u_i,\qquad
 u_N=-\sum_{i<N}u_i,\qquad Q(u)=\sum_i u_i^2,
 \qquad \mathcal V_N(u)=\prod_{i<j}(u_i-u_j).
\]
This is the literal angular constraint, with coordinate measure
$\prod_{i<N}du_i/(2\pi)$.  In this chart the physical root is
$\rho_i=e^{-\ii\varphi_i}$, where
$\varphi_i=\arccos(\mathsf S-\cos\theta_i)$.
At the central point its derivatives are
\[
 \varphi=\alpha,\qquad
 \varphi_{\mathsf S}=-\frac1{\sin\alpha},\qquad
 \varphi_\theta=1,\qquad
 \varphi_{\theta\theta}=-2\cot\alpha.
\]
Writing $h=\mathsf S-\mathsf S_*$ and $d_\alpha=\cot\alpha>0$,
the zero-sum constraint cancels both the linear shape term and its
mixed parameter derivative, so
\[
 \sum_i\varphi_i
 =N\alpha-\frac{N}{\sin\alpha}h-d_\alpha Q(u)
   +O\left(|h|^2+|h|Q(u)+\sum_i|u_i|^3\right).
\]
Consequently, for $q_\sigma=1-\sigma Z$,
\begin{align}
 q_\sigma(\mathsf S_\eps,u)
   &=2N\eps-\ii d_\alpha Q(u)
     +O\left(\eps^2+\eps Q(u)+Q(u)^2+
                                  \sum_i|u_i|^3\right),
                                  \label{eq:iso-local-q}\\
 \partial_{\mathsf S}q_\sigma(\mathsf S_*,0)
   &=-\frac{\ii N}{\sin\alpha}=:D_{\mathsf S},
                                  \label{eq:iso-local-derivative}\\
 \partial_t q_\sigma(\zeta,0)&=2N\zeta^{-1}=:D_t.\nonumber
\end{align}
All constants here may depend on the fixed $N$.  The two leading
terms in \eqref{eq:iso-local-q} have perpendicular phases.  By choosing
a small fixed real shape ball and then $\eps_0$, the remainder is
less than half their lower bound.  Thus
\begin{equation}\label{eq:iso-local-lower}
 |q_\sigma(\mathsf S_\eps,u)|\ge c(\eps+Q(u)).
\end{equation}
The cubic shape remainder is retained explicitly; no quadratic
normal-form assertion is required.

For the fixed shape cutoff $\chi$, the chart integral is exactly
\begin{equation}\label{eq:iso-shape-integral}
 I_N^\sigma(\mathsf S)=
   \int_{\R^{N-1}}
     \frac{\chi(u)\mathcal V_N(u)^2 A_N(\mathsf S,u)}
          {q_\sigma(\mathsf S,u)}\,du_1\cdots du_{N-1},
\end{equation}
with analytic amplitude and
\begin{equation}\label{eq:iso-exact-amplitude}
 A_N(\mathsf S_*,0)
 =\ii^{-N^2}
     \frac{2^{1-N(N-1)}}
          {N!(2\pi)^{N-1}\sin^{N^2}\alpha}
 =:\ii^{-N^2}a_N(\alpha),\qquad a_N(\alpha)>0.
\end{equation}
Indeed every sine numerator gives half an angle difference, every
pair denominator tends to $\ii\sin\alpha$, every one-body factor
to $(\ii\sin\alpha)^{-1}$, and $1+\sigma Z\to2$.
This includes all factorial and Haar normalizations.
At fixed shape coordinates every parameter derivative retains the
whole squared Vandermonde.  The reflected upper chart is exactly
equal under angle reflection.

\begin{lemma}[The nonzero scalar period]
\label{lem:iso-period}
Put
\[
 a=\frac{N^2-1}{2},\qquad k=\lfloor a\rfloor,\qquad
 \nu=k+1-a\in\{\tfrac12,1\},\qquad
 \mathcal S_N=\frac1{\sqrt N}
      \int_{\substack{\sum_i\omega_i=0\\\sum_i\omega_i^2=1}}
           \mathcal V_N(\omega)^2\,d\sigma(\omega)>0.
\]
Then
\begin{align}
 J_{N,\alpha}
  &:=\int_{\R^{N-1}}
     \frac{\mathcal V_N(v)^2}
          {(2N-\ii d_\alpha Q(v))^{k+1}}\,dv_1\cdots dv_{N-1}
          \nonumber\\
  &=\frac{\mathcal S_N}{2}(2N)^{-\nu}
      (-\ii d_\alpha)^{-a}\operatorname B(a,\nu)\ne0.
                                      \label{eq:iso-period}
\end{align}
The branch in the last line has $\arg(-\ii d_\alpha)=-\pi/2$.
For odd prime $N=p$, $k=(p^2-1)/2\equiv0\pmod4$ and this simplifies to
\begin{equation}\label{eq:iso-odd-period}
 J_{p,\alpha}=\frac{\mathcal S_p}{4pk\,d_\alpha^k}>0.
\end{equation}
\end{lemma}
\begin{proof}
The map $(v_1,\ldots,v_{N-1})\mapsto(v_1,\ldots,v_N)$,
$v_N=-\sum_{i<N}v_i$, has Gram determinant $N$.
Its coordinate measure is therefore Euclidean measure divided by
$\sqrt N$, exactly as in the constrained angular normalization.
Polar coordinates on that hyperplane give the radial power
$r^{N^2-2}=r^{2a-1}$.  With $x=r^2$,
\[
 J_{N,\alpha}=\frac{\mathcal S_N}{2}
   \int_0^\infty\frac{x^{a-1}}{(2N-\ii d_\alpha x)^{k+1}}\,dx.
\]
For a positive real coefficient in place of $-\ii d_\alpha$,
a real substitution gives the beta integral in
\eqref{eq:iso-period}.  Continue that coefficient through the fourth
quadrant.  Throughout the closed quadrant at fixed positive modulus,
$|2N+cx|\ge c_0(1+x)$, so dominated convergence extends the same
identity to the displayed imaginary coefficient.  The radial tail
is $r^{-2}$ for even $N$ and $r^{-3}$ for odd $N$.
The angular constant is positive because the zero-sum sphere has
an open set of vectors with distinct coordinates.

For odd $N=p$, an antiderivative gives the same period:
with $Q_0=2p$ and $c=-\ii d_\alpha$,
\[
 \frac{d}{dx}\left(\frac{cx}{Q_0+cx}\right)^k
   =kQ_0c^k\frac{x^{k-1}}{(Q_0+cx)^{k+1}}.
\]
The bracket has endpoint values zero and one; and
$c^k=d_\alpha^k$ because $k\equiv0\pmod4$.
This proves \eqref{eq:iso-odd-period} directly.
\end{proof}

\begin{theorem}[The complete first coefficient]
\label{thm:iso-first-coefficient}
For either parity and its signed positive geometry,
\begin{align}
 \partial_{\mathsf S}^{k}\mathcal F_N^\sigma(\mathsf S_\eps)
     &=\Lambda^{\mathsf S}_{N,\alpha}\eps^{-\nu}
                                      +o(\eps^{-\nu}),
                                        \label{eq:iso-first-S}\\
 (C_N^\sigma)^{(k)}(t_\eps)
     &=\Lambda^{t}_{N,\alpha}\eps^{-\nu}
                                      +o(\eps^{-\nu}),
                                        \label{eq:iso-first-t}\\
 \Lambda^{\mathsf S}_{N,\alpha}
     &=2\ii^{-N^2}a_N(\alpha)(-1)^k k!
                      D_{\mathsf S}^{k}J_{N,\alpha}\ne0,
     \nonumber\\
 \Lambda^{t}_{N,\alpha}
     &=2\ii^{-N^2}a_N(\alpha)(-1)^k k!
                      (2N\zeta^{-1})^kJ_{N,\alpha}\ne0.
                                        \label{eq:iso-first-constants}
\end{align}
All derivatives below $k$ are bounded for even $N$.
For odd $N$, they are bounded through $k-2$ and at most logarithmic
at $k-1$.
\end{theorem}
\begin{proof}
Lemma~\ref{lem:iso-finite-complement} disposes of the complete
complement.  Differentiate \eqref{eq:iso-shape-integral} $k$ times
at fixed real shape coordinates.  The unique highest pole is
\[
 \int\frac{\chi(u)\mathcal V_N(u)^2
           (-1)^k k!(\partial_{\mathsf S}q_\sigma)^k A_N}
               {q_\sigma^{k+1}}\,du.
\]
Under $u=\sqrt\eps\,v$ its total scale is
$\eps^{(N^2-1)/2-(k+1)}=\eps^{-\nu}$.
Equation~\eqref{eq:iso-local-lower} gives the integrable majorant
$C\mathcal V_N(v)^2/(1+Q(v))^{k+1}$.
Dominated convergence and \eqref{eq:iso-exact-amplitude} therefore
give the single-chart coefficient
$\ii^{-N^2}a_N(-1)^k k!D_{\mathsf S}^{k}J_{N,\alpha}$.
The reflected chart doubles it.  This reflection factor is distinct
from the factor $2$ already in $a_N$ from $1+\sigma Z$.

Every other term has denominator power at most $k$.
Its radial absolute bound is
$C\int_0^\delta r^{2a-1}(\eps+r^2)^{-m-1}\,dr$ with $m<k$.
It is bounded in the even case, and at worst $O(|\log\eps|)$ in
the odd case; both are $o(\eps^{-\nu})$.  The same radial powers
give the assertions about lower derivatives.
Finally $\mathsf S'(\zeta)=2\ii\sin\alpha\,\zeta^{-1}\ne0$.
The highest term in the finite chain rule is
$(\mathsf S')^k\partial_{\mathsf S}^{k}$; all remaining terms
are bounded or logarithmic.  The identity
$D_{\mathsf S}\mathsf S'(\zeta)=2N\zeta^{-1}$ proves the stated
exterior-coordinate constant.
\end{proof}

In particular the odd coefficient has
\[
 \Lambda^t_{p,\alpha}
   =-2\ii\,a_p(\alpha)k!(2p)^k\zeta^{-k}J_{p,\alpha}.
\]
For an original second-quadrant HT point, apply
\eqref{eq:iso-negative-jets}.  Since $k\equiv0\pmod4$,
$\sin\widehat\alpha=\sin\alpha$ and
$\widehat\zeta^{-k}=\overline\zeta^{-k}$, its constant has exactly
the same formula with
$J_{p,\alpha}=\mathcal S_p/(4pk|\cot\alpha|^k)>0$.
Thus both HT signs yield nonzero coefficients with the physical
ordinary constraint after the exact algebraic transfer.

\subsection{Noncancellation for the infinite physical sums}

\begin{theorem}[Isotropic natural boundary]
\label{thm:iso-internal-boundary}
The full LT normalized response has the unit circle as a natural
boundary in its exterior variable $s$.  The full HT normalized
response has the unit circle as a natural boundary in its disk
variable $s$.  These assertions concern the exact physical sums
\eqref{eq:iso-physical-L}--\eqref{eq:iso-physical-H}.
\end{theorem}
\begin{proof}
At a fixed prime point, normal convergence permits termwise
differentiation.  All smaller coefficients have bounded $k$th
derivatives by Lemma~\ref{lem:iso-finite-complement}.
The first one has the nonzero term in
Theorem~\ref{thm:iso-first-coefficient}, whereas
Theorem~\ref{thm:iso-tail} gives
\[
 \sum_{\substack{N>n_*\\N\ {\rm of\ the\ chosen\ parity}}}
       |\partial_{\mathsf S}^{j}\mathcal F_N^\sigma(\mathsf S_\eps)|
     =O(1+|\log\eps|),\qquad 0\le j\le k.
\]
The finite chain rule gives the same tail bound for the exterior
derivatives.  Since $\eps^\nu(1+|\log\eps|)\to0$, the entire
sum has precisely the nonzero leading singularity
$\Lambda^t\eps^{-\nu}$ in its $k$th derivative.
For negative HT dispersion use \eqref{eq:iso-negative-jets} for
every coefficient and every required derivative, including the
entire tail.  This proves the same obstruction for the physical
ordinary-sign series there.

A holomorphic continuation through such a point would bound every
fixed derivative on a short radial segment, a contradiction.
The physical multipliers $(1-s^{-4})^{1/4}$ and $Q_H(t)$ are
locally analytic and nonzero at all selected nonaxis points.
Division by those multipliers transfers the obstruction to the full
normalized observables.  Evenness and conjugation supply every
quadrant for LT.  The two separately proved upper HT quadrants
and conjugation supply every quadrant for HT.  Prime-grid density
then rules out continuation at every point of the circle, including
the four axis points: any proposed neighborhood of an axis point
would also continue through a nearby obstructing nonaxis point.
Finally $t=1/s$ is biholomorphic at every unit-circle point and
transfers the exterior HT assertion to its prescribed disk germ.
\end{proof}

\subsection{The energy variable, the physical field, and covering}

For $\lambda=1$ the explicit profiles are
\[
 H_1(y)=\operatorname{arsinh}(\cos y),\quad |y|<\pi/2,
 \qquad L_1(y)=\operatorname{arsinh}|\cos y|.
\]
The elementary identity
\begin{equation}\label{eq:iso-domain-identity}
 |\sinh(x+\ii y)|^2=\sinh^2x+\sin^2y
\end{equation}
shows that $s=\sinh z$ maps $\Omega_H$ onto $|s|<1$ and
$\Omega_L$ onto $|s|>1$ with the following exact mapping properties.
On the disk, choose the analytic square root of $1+s^2$ equal to
one at zero.  Then
$\operatorname{arsinh}s=\log(s+\sqrt{1+s^2})$, with value zero at
zero, is analytic; its image lies in the component of
$\{|\sinh z|<1\}$ containing zero, which is exactly $\Omega_H$.
The inverse identities, first near zero and then by analytic
continuation, prove that $\sinh:\Omega_H\to\{|s|<1\}$ is a
biholomorphism.

For LT set $w=e^z$ and $s=(w-w^{-1})/2$.
For every $|s|>1$, the two roots of
$w^2-2sw-1=0$ have product $-1$ and neither lies on $|w|=1$;
exactly one has modulus greater than one.  This unique outer root
depends holomorphically on $s$, since its only possible quadratic
branch points would be $s=\pm\ii$ on the excluded circle.
Thus $w\mapsto(w-w^{-1})/2$ is biholomorphic from
\[
 \mathcal W=\{|w|>1:\ |(w-w^{-1})/2|>1\}
\]
onto $|s|>1$.  Identity~\eqref{eq:iso-domain-identity} gives
$\exp^{-1}\mathcal W=\Omega_L$.
The epigraph $\Omega_L$ is simply connected, so the restriction
of the exponential is its universal covering of $\mathcal W$.
Consequently $\sinh:\Omega_L\to\{|s|>1\}$ is the universal
covering, with no omitted temperature sheets.

At every boundary point with $\cosh z\ne0$, this change of variable
is locally biholomorphic and takes the appropriate domain side to
the appropriate side of the unit circle.
Its only boundary ramification points have
$z=\ii(\pi/2+\pi l)$; they are isolated on the LT graph and are
the two HT tips.  The prime obstructions lift densely on each
remaining boundary arc.  A neighborhood continuation at an isolated
ramification point would therefore continue through a nearby lifted
prime obstruction as well.  The entire prescribed energy-variable
boundaries follow, including these tips and all LT sheets.

Finally retain the physical field convention from
Section~\ref{sec:physical}: the Hamiltonian field has energy units,
$z=2\beta J_x$, and
\[
 \chi_H(z)=\frac{z}{2J_x}\mathcal X_H(z),\qquad
 \chi_L(z)=\frac{z}{2J_x}\mathcal X_L(z).
\]
The preceding thermodynamic GHS argument identifies this with the
literal uniform-field derivative, with the plus phase fixed before
the zero-field limit in LT.  Its finite covariance sum and the
$M^2$ subtraction were already preserved in
\eqref{eq:iso-physical-L}.  The factor $z/(2J_x)$ is analytic and
nonzero at every boundary point.  It therefore preserves all the
local obstructions just proved.  In particular one must use this
single-valued physical energy variable on the LT covering, rather
than imposing an artificial globally single-valued inverse
$\operatorname{arsinh}s$ on the exterior $s$ plane.


\begin{proof}[Completion of the proof of Theorem~\ref{thm:aniso}]
The unequal-ratio part was proved in Section~\ref{sec:closure}.
At equal coupling Theorem~\ref{thm:iso-internal-boundary} proves
the full normalized obstructions, and the covering and
physical-field calculation above transfers them to the prescribed
energy-variable domains. No limit in $\lambda$ or phase duality is used. This proves the theorem for every finite positive ratio on
both physical temperature sides.
\end{proof}

\begin{corollary}[Inverse-temperature barriers]\label{cor:alllambda}
For each fixed $J_x>0$, the barriers in the $\beta$ plane are exactly
$\partial\OH/(2J_x)$ and $\partial\OL/(2J_x)$ for their respective
physical germs. At $\lambda=1$ they are the correctly interpreted
pullbacks of $|\sinh(2\beta J)|=1$, with the low-temperature covering
specified above.
\end{corollary}


\begin{thebibliography}{99}
\bibitem{BL} A. I. Bugrij and O. Lisovyy,
Correlation function of the two-dimensional Ising model on a finite lattice. II,
\emph{Theoretical and Mathematical Physics} \textbf{140} (2004), 987--1000.
\href{https://arxiv.org/abs/0708.3643v1}{arXiv:0708.3643v1}.

\bibitem{IL} N. Iorgov and O. Lisovyy,
Ising correlations and elliptic determinants,
\emph{Journal of Statistical Physics} \textbf{143} (2011), 33--59.
\href{https://doi.org/10.1007/s10955-011-0154-6}{doi:10.1007/s10955-011-0154-6}.
\href{https://arxiv.org/abs/1012.2856v2}{arXiv:1012.2856v2}.

\bibitem{ONGP} W. P. Orrick, B. G. Nickel, A. J. Guttmann and J. H. H. Perk,
The susceptibility of the square lattice Ising model: New developments,
\emph{Journal of Statistical Physics} \textbf{102} (2001), 795--841.
\href{https://doi.org/10.1023/A:1004850919647}{doi:10.1023/A:1004850919647}.
\href{https://arxiv.org/abs/cond-mat/0103074}{arXiv:cond-mat/0103074}.

\bibitem{BHM} S. Boukraa, S. Hassani and J.-M. Maillard,
Holonomic functions of several complex variables and singularities of
anisotropic Ising $n$-fold integrals,
\emph{Journal of Physics A: Mathematical and Theoretical}
\textbf{45} (2012), 494010.
\href{https://doi.org/10.1088/1751-8113/45/49/494010}{doi:10.1088/1751-8113/45/49/494010}.
\href{https://arxiv.org/abs/1207.1784v2}{arXiv:1207.1784v2}.

\bibitem{Raoufi} A. Raoufi,
Translation-invariant Gibbs states of Ising model: General setting,
preprint (2017).
\href{https://arxiv.org/abs/1710.07608v1}{arXiv:1710.07608v1}.

\bibitem{CGN} F. Camia, C. Garban and C. M. Newman,
The Ising magnetization exponent on $\mathbb Z^2$ is $1/15$,
preprint (2013).
\href{https://arxiv.org/abs/1205.6612v4}{arXiv:1205.6612v4}.

\end{thebibliography}
\end{document}
